\providecommand{\zerJ}{\mathcal{J}}
\providecommand{\zerK}{\mathcal{K}}
\providecommand{\zersinc}{\operatorname{sinc}}
\providecommand{\zergE}{\gamma_{\mathrm{E}}}
\providecommand{\zerLL}{\mathcal{L}}
\providecommand{\zerconv}{\operatorname{conv}}
\providecommand{\zerdiag}{\operatorname{diag}}
\providecommand{\zerSym}{\mathfrak{D}}
\providecommand{\zerLog}{\operatorname{Log}}
\providecommand{\anapos}[1]{\bigl(#1\bigr)_{+}}
\providecommand{\ananeg}[1]{\bigl(#1\bigr)_{-}}
\providecommand{\anaone}{\mathbf{1}}
\providecommand{\anaE}{\mathcal{E}}
\providecommand{\anaEx}{\mathbb{E}}
\providecommand{\anaPr}{\mathbb{P}}
\providecommand{\anaM}{\mathcal{M}}
\providecommand{\anaW}{\mathcal{W}}
\providecommand{\anaZ}{\mathcal{Z}}
\providecommand{\anaC}{\mathcal{C}}
\section{Prime errors at several scales: covariance, balancing, peaks and merging}\label{sec:flu}

This section studies the smoothed errors of the prime-power sums in residue classes, with the
weight $(1/n-1/x)$ of the source proof, as functions of the class and of the scale. Under
\cref{src:thm:dgrh} they are almost periodic in $\log x$, and the questions are about their joint
law: positive definiteness and oscillation of the covariance at several scales, with derivatives
or after a compact filter, without any independence hypothesis on the zeros
(\cref{zer:sec:cov}); simultaneous balancing of all races in one multiplicative interval
(\cref{zer:sec:balancing}); extreme values under positive spectral weights, including the
excursions of the source remainder (\cref{zer:sec:peaks}); and the variance of stencils when
scales merge, which reads off the multiplicities of zeros (\cref{zer:sec:coalesce}). The
algebraic input, rigidity of polynomially weighted zero divisors, comes first
(\cref{flu:sec:polyrigid}). The explicit formula used throughout is \cref{zer:lem:ramp}, and the
mass of the zero spectrum is \cref{zer:lem:mass}. Sign questions for a single race are in
\cref{sec:zer}, and positive filters that cancel prescribed zeros, among them a compact kernel
for all residue classes, in \cref{sec:fil}.

\subsection{Polynomially weighted multiplicity divisors}\label{flu:sec:polyrigid}
{\emergencystretch=6em\sources{grh-methods-polynomial-density; grh_cm_four_terms_and_prime_bias; grh_spectral_rigidity; grh_filtered_prime_races; grh_prime_error_balancing}}

The multiplicity rank statement \cref{zer:thm:rank} in \cref{sec:chz} localizes in height but
uses constant weights. The next statement allows
weights that are polynomials in the zero, characters of different conductors, and an
arbitrary finite exceptional set; it is the algebraic input of the covariance theorems of
\cref{zer:sec:cov}. Let $\chi_1,\dots,\chi_r$ be distinct primitive Dirichlet characters, of
conductors $q_j$, including the conductor-one character if desired. Write $a_j\in\set{0,1}$ for
the parity, $\chi_j(-1)=(-1)^{a_j}$, and $\delta_j=1$ for the conductor-one character and $0$
otherwise. The completed functions
\[
 X_j(s)=[s(s-1)]^{\delta_j}(q_j/\pi)^{(s+a_j)/2}\Gamma\bigl((s+a_j)/2\bigr)L(s,\chi_j)
\]
are entire of order one, do not vanish at $0$ and $1$, and their zeros are precisely the
nontrivial zeros. Put $m_j(\rho)=\ord_{s=\rho}X_j(s)$, with $m_j(\rho)=0$ when $X_j(\rho)\ne0$.
By Hadamard factorization,
\begin{equation}\label{zer:eq:hadamard}
 \frac{X_j'}{X_j}(s)=B_j+\sum_{\rho}m_j(\rho)\Bigl(\frac1{s-\rho}+\frac1\rho\Bigr),\qquad
 \sum_{\rho}\frac{m_j(\rho)}{\abs\rho^2}<\infty,
\end{equation}
with locally normal convergence away from the zeros and no assumption of simple zeros.

\begin{theorem}[Polynomial zero-divisor rigidity; \stU]\label{zer:thm:polyrigid}
If polynomials $P_1,\dots,P_r\in\C[s]$ satisfy
\[
 \sum_{j=1}^rP_j(\rho)m_j(\rho)=0
\]
at every point $\rho$ of the union of the zero sets of the $X_j$ except finitely many, then
$P_1=\dots=P_r=0$.
\end{theorem}

\begin{proof}
Let $d\ge0$ bound the degrees and put
\[
 F(s)=\sum_jP_j(s)\frac{X_j'}{X_j}(s),\qquad A(s)=\sum_jP_j(s)\frac{L'}{L}(s,\chi_j).
\]
First, $F$ is a rational function. Truncate the sums in \eqref{zer:eq:hadamard} to a common
disk $\abs\rho\le R$ containing the exceptional points. For each $\rho$,
$P_j(s)/(s-\rho)=P_j(\rho)/(s-\rho)+(P_j(s)-P_j(\rho))/(s-\rho)$ splits a simple pole and a
polynomial of degree at most $d-1$; the terms $P_j(s)/\rho$ and $B_jP_j(s)$ have degree at most
$d$. By hypothesis the pole part of the truncated $F$ does not change once $R$ is large, and the
remaining polynomial has degree at most $d$. Convergence at $d+1$ distinct regular points forces
convergence of its coefficients, by invertibility of the Vandermonde matrix. So $F$ is a fixed
finite sum of simple fractions plus a polynomial.

Second, write $g_a(s)=\tfrac12\psi((s+a)/2)$ with $\psi=\Gamma'/\Gamma$. Logarithmic
differentiation gives
\[
 F(s)=A(s)+\sum_jP_j(s)\Bigl[\tfrac12\log(q_j/\pi)+g_{a_j}(s)
 +\delta_j\Bigl(\frac1s+\frac1{s-1}\Bigr)\Bigr].
\]
Let $\Delta f(s)=f(s+2)-f(s)$. The recurrence of the Gamma function gives
$g_a(s+2k)=g_a(s)+\sum_{h=0}^{k-1}(s+a+2h)^{-1}$. In $\Delta^{d+1}(P_jg_{a_j})$ the
coefficient of $g_{a_j}(s)$ is $\Delta^{d+1}P_j=0$, and every other term is rational. Hence
$\Delta^{d+1}A$ is rational. For real $\sigma\to+\infty$ the Euler product gives
\begin{equation}\label{zer:eq:polydir}
 A(\sigma)=-\sum_{n\ge2}\Lambda(n)n^{-\sigma}\sum_jP_j(\sigma)\chi_j(n)
 =O\bigl((1+\sigma)^d2^{-\sigma}\bigr),
\end{equation}
for instance because $\sum_{n\ge2}(\log n)n^{-\sigma}\le2^{2-\sigma}\sum_{n\ge2}(\log n)n^{-2}$
for $\sigma\ge2$, and the same bound holds for the fixed difference $\Delta^{d+1}A$. A nonzero
rational function does not decay exponentially along the positive axis, so $\Delta^{d+1}A=0$ on
$\Re s>1$. For each real $\sigma>1$ the sequence $u_n=A(\sigma+2n)$ then has vanishing
difference of order $d+1$, so it is a polynomial in $n$; it tends to zero by
\eqref{zer:eq:polydir}, so it vanishes, and $A(\sigma)=0$.

Third, all coefficients $C_n(s)=\Lambda(n)\sum_jP_j(s)\chi_j(n)$ vanish identically: if $n_0$
were the least index with $C_{n_0}\ne0$, multiplying \eqref{zer:eq:polydir} by $n_0^\sigma$
would give $\abs{C_{n_0}(\sigma)}\ll(1+\sigma)^d(n_0/(n_0+1))^{\sigma-2}$, which is impossible
for a nonzero polynomial. Put $Q=\lcm(q_1,\dots,q_r)$. Dirichlet's theorem gives a prime
$p\equiv b\pmod Q$ for every unit $b$, and $\Lambda(p)>0$, so $\sum_jP_j(s)\chi_j(b)=0$ for
every $b\in(\Z/Q\Z)^\times$. The characters induced by the $\chi_j$ modulo $Q$ are distinct, by
uniqueness of the primitive inducing character, and orthogonality gives $P_j=0$ for all $j$.
\end{proof}

\begin{corollary}[Tail spanning; \stG]\label{zer:cor:tail}
Write $m_j(\gamma)=m_j(\tfrac12+i\gamma)$. For every $H\ge0$ and every integer $d\ge0$ the
vectors
\[
 v(\gamma)=\bigl(m_j(\gamma)(i\gamma)^\ell\bigr)_{1\le j\le r,\ 0\le\ell\le d},\qquad
 \abs\gamma>H,
\]
span $\C^{r(d+1)}$.
\end{corollary}
\begin{proof}
A linear functional with coefficients $c_{j\ell}$ annihilating these vectors gives polynomials
$P_j(s)=\sum_\ell c_{j\ell}(s-\tfrac12)^\ell$ that satisfy the hypothesis of
\cref{zer:thm:polyrigid}: by \cref{src:thm:dgrh} every zero is $\tfrac12+i\gamma$, and only
finitely many distinct ordinates have $\abs\gamma\le H$. Hence all $c_{j\ell}=0$.
\end{proof}

Common zeros, arbitrary multiplicities and both signs of the ordinates are allowed. For $d=0$
and characters to one modulus the statement also follows from \cref{zer:thm:rank} (applied
modulo the least common multiple of the conductors), which in addition localizes the spanning
zeros in every window of logarithmic length; the polynomial weights $(i\gamma)^\ell$ are what
the derivatives in \cref{zer:thm:jets} require.

\subsection{Covariance and oscillation of prime errors at several scales}\label{zer:sec:cov}
{\emergencystretch=6em\sources{multiscale_prime_fluctuations; grh_multiscale_covariance; grh_spectral_rigidity; grh_filtered_prime_races; grh_prime_error_balancing; prime_peak_rigidity; grh_extensions_2026-11-27}}

The kernel $(1/n-1/x)_+$ of the prime remainder in the source proof has Mellin denominator
$s(s-1)$, and $\rho(1-\rho)=\tfrac14+\gamma^2$ at $\rho=\tfrac12+i\gamma$. Under
\cref{src:thm:dgrh} the normalized smoothed errors are therefore absolutely convergent
almost periodic functions of $\log x$. This subsection proves that their joint covariances,
at several multiplicative scales, in all residue classes, with derivatives, or after a
compact filter, are positive definite even after any finite set of zero frequencies is
deleted, and derives uniform two-sided oscillation in every direction; the multiscale statement
also holds for the ordinary counts $\pi(x;q,a)$ (\cref{zer:thm:unsmoothed}). Simplicity of zeros,
disjointness of zero sets and rational independence of ordinates are never assumed:
coincident ordinates are grouped before squares are formed, and positivity comes from the
arithmetic of prime powers.

\paragraph{Setting.}
Fix $q\ge1$, put $G=(\Z/q\Z)^\times$ (the trivial group for $q=1$) and $n_q=\varphi(q)$, and let
$\chi$ run over the characters modulo $q$, extended by zero to nonunits, with primitive
inducing character $\chi^*$ and $\delta_\chi=1$ for the principal character, $0$ otherwise. For
$x\ge1$ and $a\in G$ put
\[
 S_\chi(x)=\sum_{n\le x}\Lambda(n)\chi(n)\Bigl(\frac1n-\frac1x\Bigr),\qquad
 S_a(x)=\sum_{\substack{n\le x\\n\equiv a\bmod q}}\Lambda(n)\Bigl(\frac1n-\frac1x\Bigr)
 =\frac1{n_q}\sum_\chi\overline{\chi(a)}S_\chi(x),
\]
and
\[
 c_\chi=\lim_{s\to1}\Bigl(-\frac{L'}{L}(s,\chi)-\frac{\delta_\chi}{s-1}\Bigr)-\delta_\chi,\qquad
 c_a=\frac1{n_q}\sum_\chi\overline{\chi(a)}c_\chi .
\]
For the principal character $c_\chi=-1-\zergE-\sum_{p\mid q}\log p/(p-1)$. The weight vanishes at
$n=x$, so $S_\chi$ and $S_a$ are continuous. Under \cref{src:thm:dgrh} let $m_\chi(\gamma)$ be the
order of $L(s,\chi^*)$ at $\tfrac12+i\gamma$, let $\Gamma_q$ be the set of distinct ordinates of
all the $L(s,\chi^*)$, and put
\begin{equation}\label{zer:eq:va}
 v_a(\gamma)=\frac1{n_q}\sum_\chi\overline{\chi(a)}\,m_\chi(\gamma),\qquad
 \beta_a=4v_a(0),\qquad
 Z_a(t)=e^{t/2}\Bigl(S_a(e^t)-\frac t{n_q}-c_a\Bigr)-\beta_a\quad(t\ge0).
\end{equation}
Since $m_{\overline\chi}(-\gamma)=m_\chi(\gamma)$, one has $v_a(-\gamma)=\overline{v_a(\gamma)}$, and
$c_a,\beta_a$ are real. The constant $\beta_a$ records possible central zeros; their absence is
not assumed. For a character put
\[
 K_\chi(t)=\sum_{\gamma}\frac{m_\chi(\gamma)e^{i\gamma t}}{1/4+\gamma^2},
\]
the sum over all ordinates of $L(s,\chi^*)$, the central one included. The local zero count
$O_q(\log(2+\abs y))$ in $[y-1,y+1]$ makes this series absolutely and uniformly convergent.

\begin{lemma}[Ramp explicit formula; \stU]\label{zer:lem:ramp}\label{rec:lem:ramp}
For every $\chi$ modulo $q$ and $x\ge2$,
\[
 S_\chi(x)=\delta_\chi\log x+c_\chi+\sum_\rho\frac{x^{\rho-1}}{\rho(1-\rho)}+O_q\bigl((1+\log x)/x\bigr),
\]
the sum running over the zeros $\rho$ of $L(s,\chi^*)$ in the open strip, with multiplicity; it
converges absolutely. Hence, as $t\to\infty$,
\[
 e^{t/2}\Bigl(S_a(e^t)-\frac t{n_q}-c_a\Bigr)=\frac1{n_q}\sum_{\chi}\overline{\chi(a)}\sum_\rho
 \frac{e^{(\rho-1/2)t}}{\rho(1-\rho)}+O_q\bigl((1+t)e^{-t/2}\bigr).
\]
For a single primitive character the difference between its normalized sum and its zero series is
real analytic in $t>0$. Under \cref{src:thm:dgrh} the zero sum is $x^{-1/2}K_\chi(\log x)$, so that
\[
 S_\chi(x)=\delta_\chi\log x+c_\chi+x^{-1/2}K_\chi(\log x)+O_q\bigl((1+\log x)/x\bigr),
\]
and consequently, uniformly in $a\in G$, for $t\ge\log2$,
\begin{equation}\label{zer:eq:Zexp}
 Z_a(t)=\sum_{\substack{\gamma\in\Gamma_q\\\gamma\ne0}}\frac{v_a(\gamma)}{1/4+\gamma^2}\,e^{i\gamma t}
 +O_q\bigl((1+t)e^{-t/2}\bigr),
\end{equation}
and the series converges absolutely and uniformly on $\R$.
\end{lemma}
\begin{proof}
First let $\chi$ be primitive and $A(s)=-L'(s,\chi)/L(s,\chi)$. For $c>1$, termwise Mellin
inversion of $y^s/(s(s-1))$, whose inverse transform is $(y-1)_+$, gives
\[
 S_\chi(x)=\frac1{2\pi i}\int_{c-i\infty}^{c+i\infty}A(s)\frac{x^{s-1}}{s(s-1)}\,\dd s .
\]
Move the line to $\Re s=-\tfrac12$. The functional equation and Stirling's formula give
$A(-\tfrac12+iy)=O_q(\log(2+\abs y))$. The local zero count lets the horizontal edges pass at
heights tending to $\pm\infty$ at distance $\gg_q1/\log(2+\abs y)$ from every ordinate; there
the local partial-fraction expansion bounds $A$ by $O_q(\log^2(2+\abs y))$, so these edges
vanish for fixed $x$, and the new vertical integral is $O_q(x^{-3/2})$. The residue at $s=1$
is $\delta_\chi\log x+c_\chi$. At a zero $\rho$ of multiplicity $m$ it is
$mx^{\rho-1}/(\rho(1-\rho))$, which is $mx^{-1/2+i\gamma}/(\tfrac14+\gamma^2)$ when
$\rho=\tfrac12+i\gamma$; since $\sum_\rho\abs{\rho(1-\rho)}^{-1}<\infty$, the zero sum converges
absolutely. At $s=0$ the function $A$
has at most a simple pole, so the residue is $x^{-1}$ times a polynomial of degree at most one
in $\log x$; no other pole is crossed. Shifting further to $\Re s=-J-\tfrac12$, a trivial zero
$-j$ adds $-x^{-j-1}/(j(j+1))$ times its order and the left integrals tend to zero as
$J\to\infty$; this series extends holomorphically to $\Re t>0$ in $x=e^t$, which gives the real
analyticity. For imprimitive $\chi$, each prime $p\mid q$ with
$p\nmid f_\chi$ removes
\[
 \sum_{p^k\le x}\chi^*(p)^k\log p\Bigl(\frac1{p^k}-\frac1x\Bigr)
 =\frac{\chi^*(p)\log p}{p-\chi^*(p)}+O_p\Bigl(\frac{1+\log x}x\Bigr),
\]
by the geometric tail and the at most $\log x/\log p$ terms of the $1/x$ sum. The constant is
the change of $c_\chi$ caused by the removed Euler factor, and the nontrivial zeros are those
of $\chi^*$. Character orthogonality gives \eqref{zer:eq:Zexp}; the frequency $\gamma=0$
contributes $4v_a(0)=\beta_a$.
\end{proof}

All the statements below rest on one elementary principle.

\begin{lemma}[Almost periodic vectors; \stU]\label{zer:lem:ap}
Let $F\subset\R\setminus\set0$ be countable with $F=-F$, let $z_\gamma\in\C^N$ satisfy
$z_{-\gamma}=\overline{z_\gamma}$ and $B=\sum_{\gamma\in F}\norm{z_\gamma}<\infty$, and put
$V(t)=\sum_\gamma z_\gamma e^{i\gamma t}$. Let $X:[t_0,\infty)\to\R^N$ be locally bounded and
measurable.
\begin{enumerate}[label=\textup{(\roman*)}]
\item If $\norm{X(t)-V(t)}\to0$, or if $X-V$ is bounded and $\int_{t_0}^\infty\norm{X-V}\,\dd t<\infty$,
then, uniformly in $u\ge t_0$ as $L\to\infty$,
\[
 \frac1L\int_u^{u+L}X\,\dd t\to0,\qquad
 \frac1L\int_u^{u+L}XX^{\mathsf T}\dd t\to\mathcal G:=\sum_\gamma z_\gamma z_\gamma^*,
\]
and $\mathcal G$ is real symmetric positive semidefinite. For every $w\in\R^N$ and $c\in\R$ the
mean of $(w\cdot X-c)^2$ over $[u,u+L]$ tends to $w^{\mathsf T}\mathcal Gw+c^2$.
\item Suppose $\lambda=\lambda_{\min}(\mathcal G)>0$ and $\norm{X-V}\to0$. Put
$\eps=\lambda/(32B)$ and $\eta=\lambda/(16B^2)$. There are $L_0,T_0$ such that for every
unit $w\in\R^N$, every $T\ge T_0$ and every $L\ge L_0$,
\[
 \abs{\set{t\in[T,T+L]:w\cdot X(t)\ge\eps}}\ge\eta L,\qquad
 \abs{\set{t\in[T,T+L]:w\cdot X(t)\le-\eps}}\ge\eta L .
\]
If $X$ is continuous, the convex hull of $\set{X(t):T\le t\le T+L}$ contains the closed ball
of radius $\eps$ about $0$. In any case the closed convex hull of the cluster set of $X(t)$ as
$t\to\infty$ contains that ball.
\item Suppose $\lambda>0$, $X$ is continuous, $X-V$ is bounded and integrable on
$[t_0,\infty)$, and put $B'=\max(1,\sup\norm X)$. There is $L_1$ such that every interval
$[u,u+L]$ with $u\ge t_0$ and $L\ge L_1$ contains points $t_\pm$ with
$\pm w\cdot X(t_\pm)\ge\lambda/(8B')$, for every unit $w$.
\end{enumerate}
\end{lemma}
\begin{proof}
(i) The mean of $e^{i\omega t}$ over $[u,u+L]$ has modulus at most $2/(L\abs\omega)$ for
$\omega\ne0$. This proves the limits for finite truncations of $V$, the second by retaining
only the pairs $(\gamma,-\gamma)$. Absolute summability bounds the truncation error uniformly,
and the hypothesis on $X-V$ makes its contribution to both means $o(1)$ uniformly in $u$. The
matrix is real because $z_{-\gamma}=\overline{z_\gamma}$. The last assertion follows by
expanding the square.

(ii) Let $g=w\cdot V$, so $\abs g\le B$. By (i), uniformly in $u$ and in unit $w$ (the
convergence is that of finitely many vector and matrix entries), there is $L_0$ with
$\abs{\mathbb E g}\le\lambda/(4B)$ and $\mathbb Eg^2\ge\lambda/2$ for windows of length
$L\ge L_0$, where $\mathbb E$ is the window average. Since $g^2\le B\abs g$, we get
$\mathbb E\abs g\ge\lambda/(2B)$ and $\mathbb Eg_\pm=(\mathbb E\abs g\pm\mathbb Eg)/2\ge\lambda/(8B)$.
With $\theta=2\eps=\lambda/(16B)$, the inequality $g_+\le\theta+B\mathbf 1_{g\ge\theta}$ gives
$\mathbb P(g\ge\theta)\ge\eta$, and likewise for $-g$. Choose $T_0$ with $\norm{X-V}\le\eps$
beyond $T_0$. For continuous $X$ the hull of a compact image is compact; if a point of the
ball were outside it, strict separation would give a unit $w$ whose maximum on the hull is
less than $\eps$, contradicting what was just proved. For the cluster set, each unit $w$ has
arbitrarily late $t$ with $w\cdot X(t)\ge\eps$; a convergent subsequence gives a cluster point
with the same inequality, and separation again gives the inclusion.

(iii) Let $f=w\cdot X$, so $\abs f\le B'$. By (i) there is $L_1$ with $\mathbb Ef^2\ge\lambda/2$
and $\abs{\mathbb Ef}\le\lambda/(4B')$ on every window $[u,u+L]$, $u\ge t_0$, $L\ge L_1$. Then
$\mathbb E\abs f\ge\lambda/(2B')$ and $\mathbb Ef_\pm\ge\lambda/(8B')$, and a continuous function
attains a maximum at least its mean.
\end{proof}

\paragraph{Several scales.}
Fix $r\ge1$ distinct real numbers $h_1,\dots,h_r$; the observed multiplicative scales are
$e^{h_j}$. Put $D=rn_q$. For $w=(w_{j,a})\in\C^D$, real $\gamma$, and $1<T<T+H$, define
\[
 A_w(\gamma)=\sum_{j=1}^r\sum_{a\in G}w_{j,a}e^{i\gamma h_j}v_a(\gamma),\qquad
 V_{T,H}(w)=2\sum_{\substack{\gamma\in\Gamma_q\\T<\gamma\le T+H}}
 \frac{\abs{A_w(\gamma)}^2}{(1/4+\gamma^2)^2}.
\]
For real $w$, $V_{T,H}(w)$ is the mean square of the real signal
$2\Re\sum_{T<\gamma\le T+H}A_w(\gamma)(\tfrac14+\gamma^2)^{-1}e^{i\gamma t}$ carried by the band.

\begin{theorem}[Multiscale spectral coercivity; \stG]\label{zer:thm:coercive}
For fixed $q$ and distinct $h_1,\dots,h_r$ there are positive constants $c,C,K,T_0$,
depending only on $q$ and the shifts, such that for every $T\ge T_0$, every $H$ with
$K\log(2T)\le H\le T$ and every $w\in\C^D$,
\begin{equation}\label{zer:eq:coercive}
 \frac{cH}{T^4\log(2T)}\norm w_2^2\ \le\ V_{T,H}(w)\ \le\ \frac{CH\log^2(2T)}{T^4}\norm w_2^2 .
\end{equation}
In particular each such band contains at least $D$ distinct ordinates, and the vectors
$(e^{i\gamma h_j}v_a(\gamma))_{j,a}$ for the ordinates in the band span $\C^D$.
\end{theorem}

No uniformity in $q$, in the number of scales, or as two shifts approach each other is
asserted; the dependence on merging scales is the subject of \cref{zer:sec:coalesce}. The proof
uses two ingredients. The first is Landau's formula at a fixed argument: for real $u\ge2$,
uniformly for $1<T<T+H\le2T$,
\begin{equation}\label{zer:eq:landaufixed}
 \sum_{T<\gamma\le T+H}m_\chi(\gamma)e^{i\gamma\log u}
 =-\frac{H}{2\pi\sqrt u}\Lambda(u)\chi^*(u)+O_{q,u}(\log(2T)),
\end{equation}
where the main term is zero unless $u$ is an integral prime power. This is Banks's extension of
the Landau--Gonek formula \cite[Theorem~1.2]{Banks2023Comparing}, divided by $\sqrt u$ using
\cref{src:thm:dgrh}; for fixed $u$ the distance to the nearest other prime power is positive,
and half-open instead of open endpoints change the sum by $O_{q,u}(\log(2T))$. Banks's paper
also supplies the contour estimates used in the proof of \cref{zer:lem:rational} in \cref{sec:chz}, which is uniform
in rational arguments of bounded height but has the weaker remainder $O(Q^4\log^2(2T))$; the fixed
form \eqref{zer:eq:landaufixed} is needed here because the probe ratios $u$ may be irrational.

\begin{lemma}[Prime probes; \stU]\label{zer:lem:probes}
For a real $R>0$ with $R\ne1$, only finitely many primes $p$ make $Rp$ an integral prime power.
Consequently, for each $a\in G$ there is a prime $p_a\equiv a\pmod q$ such that
$u_{a,j,k}=p_ae^{h_k-h_j}\ge2$ for all $j,k$, and $u_{a,j,k}$ is an integral prime power only for
$j=k$.
\end{lemma}
\begin{proof}
If $Rp$ is an integer, $R=u/v$ is rational in lowest terms. If $v>1$ then $v\mid p$, which
allows one prime at most; if $v=1$ then $u\ge2$ and $up$ is a prime power only if $p\mid u$.
For the probes, exclude the finitely many primes dividing a numerator or denominator of a
rational ratio $e^{h_k-h_j}\ne1$, require $p_a\min_{j,k}e^{h_k-h_j}\ge2$, and use Dirichlet's
theorem.
\end{proof}

\begin{proof}[Proof of \cref{zer:thm:coercive}]
Let $B_{T,H}$ be the set of ordinates in $(T,T+H]$ and $L=\log(2T)$. For fixed $j,a$,
\[
 \sum_{\gamma\in B_{T,H}}A_w(\gamma)e^{i\gamma(\log p_a-h_j)}
 =\frac1{n_q}\sum_{k,b,\chi}w_{k,b}\overline{\chi(b)}\sum_{\gamma\in B_{T,H}}
 m_\chi(\gamma)e^{i\gamma\log u_{a,j,k}} .
\]
By \eqref{zer:eq:landaufixed} and \cref{zer:lem:probes} only $k=j$ gives a main term. There
$u=p_a$ is prime to $q$, so $\chi^*(p_a)=\chi(a)$, and orthogonality isolates $w_{j,a}$:
\begin{equation}\label{zer:eq:recover}
 \sum_{\gamma\in B_{T,H}}A_w(\gamma)e^{i\gamma(\log p_a-h_j)}
 =-\frac{H\log p_a}{2\pi\sqrt{p_a}}\,w_{j,a}+O_{q,h}\bigl(\norm w_2L\bigr),
\end{equation}
uniformly in $j,a,w$, since the errors total at most a constant times
$\sum_{k,b}\abs{w_{k,b}}L\le\sqrt D\norm w_2L$. Put $\alpha=\min_a\log p_a/(2\pi\sqrt{p_a})$,
choose $(j,a)$ with $\abs{w_{j,a}}\ge\norm w_2/\sqrt D$, and take $K$ so large that for
$H\ge KL$ the left side of \eqref{zer:eq:recover} has modulus at least
$H\alpha\norm w_2/(2\sqrt D)$. With $M_\gamma=\sum_\chi m_\chi(\gamma)$ the zero count gives
$\abs{B_{T,H}}\le\sum_{\gamma\in B_{T,H}}M_\gamma\le C_0HL$ and $\max M_\gamma\le C_1L$.
Cauchy--Schwarz over the distinct ordinates gives
\[
 \sum_{\gamma\in B_{T,H}}\abs{A_w(\gamma)}^2\ge\frac{\alpha^2}{4DC_0}\frac HL\norm w_2^2 ,
\]
and $\tfrac14+\gamma^2\le5T^2$ in the band gives the lower bound with $c=\alpha^2/(50DC_0)$. For
the upper bound, $\abs{v_a(\gamma)}\le M_\gamma/n_q$ gives
$\abs{A_w(\gamma)}\le(\sqrt D/n_q)M_\gamma\norm w_2$, so the sum of squares is at most
$(D/n_q^2)C_0C_1HL^2\norm w_2^2$, and $(\tfrac14+\gamma^2)^2\ge T^4$. The lower bound makes
$w\mapsto(A_w(\gamma))_{\gamma\in B_{T,H}}$ injective, which gives the last assertion.
\end{proof}

Thus any coefficient array is recovered from aggregated zero data of one band by
\eqref{zer:eq:recover}, with error $O(\norm w_2\log(2T)/H)$ per entry. For real arrays the
matrix $Q_{T,H}$ of the form $V_{T,H}$ has all its eigenvalues between the two bounds of
\eqref{zer:eq:coercive}, so $\det Q_{T,H}\ge(cH/(T^4\log(2T)))^D$ and its condition number is
at most $(C/c)\log^3(2T)$.

\begin{corollary}[All-lag strict positivity; \stG]\label{zer:cor:multiscale}
Let $J$ be a finite subset of $\Gamma_q$ with $J=-J$, and define the real vector
$U_J(t)\in\R^D$ by
\[
 U_J(t)_{j,a}=Z_a(t+h_j)-\sum_{\gamma\in J\setminus\set0}\frac{v_a(\gamma)}{1/4+\gamma^2}e^{i\gamma(t+h_j)} .
\]
Its long-time covariance in the sense of \cref{zer:lem:ap}\textup{(i)} is
\[
 \mathcal G_J=\sum_{\gamma\in\Gamma_q\setminus(J\cup\set0)}z(\gamma)z(\gamma)^*,\qquad
 z(\gamma)=\Bigl(\frac{e^{i\gamma h_j}v_a(\gamma)}{1/4+\gamma^2}\Bigr)_{j,a},
\]
and it is positive definite: $\lambda_{\min}(\mathcal G_J)\ge cH/(T^4\log(2T))$ for every band
$(T,T+H]$ of \cref{zer:thm:coercive} lying above $\max J$. Consequently
\cref{zer:lem:ap}\textup{(ii)} applies to $U_J$ with constants depending only on $q$, the
shifts and $J$: for every real unit array $w$, each of $w\cdot U_J(\log x)\ge\eps$ and
$w\cdot U_J(\log x)\le-\eps$ holds on a set of logarithmic measure at least $\eta\log K$ in
every interval $[X,KX]$ with $X$ and $K$ sufficiently large; the convex hull of the values on
each such interval contains the ball of radius $\eps$. Moreover, for one scale $j$, a nonzero
$w\in\R^{n_q}$, arbitrary real predictor coefficients $b_{k,a}$ $(k\ne j)$ and an arbitrary
intercept $c$,
\[
 \lim_{L\to\infty}\frac1L\int_u^{u+L}\Bigl(\sum_aw_aZ_a(t+h_j)-\sum_{k\ne j,\,a}b_{k,a}Z_a(t+h_k)-c\Bigr)^2\dd t
 \ \ge\ \lambda_{\min}(\mathcal G_\varnothing)\norm w_2^2 .
\]
\end{corollary}
\begin{proof}
By \eqref{zer:eq:Zexp}, $U_J(t)=\sum_{\gamma\notin J\cup\set0}z(\gamma)e^{i\gamma t}+O((1+t)e^{-t/2})$,
with $z(-\gamma)=\overline{z(\gamma)}$ and $\sum\norm{z(\gamma)}<\infty$, so \cref{zer:lem:ap}(i)
gives $\mathcal G_J$. For real $w$, $w\cdot z(\gamma)=A_w(\gamma)/(\tfrac14+\gamma^2)$; pairing $\gamma$
with $-\gamma$ gives $w^{\mathsf T}\mathcal G_Jw\ge V_{T,H}(w)$ for every band above $\max J$, and
\eqref{zer:eq:coercive} gives the bound. The ramp sums are continuous, so the hull statement
applies. For prediction, collect the coefficients in an array whose block at scale $j$ is
$w$; its squared norm is at least $\norm w_2^2$, and \cref{zer:lem:ap}(i) gives the limit
$v^{\mathsf T}\mathcal G_\varnothing v+c^2$.
\end{proof}

The theorem says more than distinctness of zero sets: common ordinates contribute to the
same vector $z(\gamma)$, and positivity survives the deletion of any bounded part of the
spectrum. Arbitrary real shifts are allowed, including logarithms of rational numbers
arbitrarily close to $1$, because the probes separate scaled prime powers rather than zeros.
The characters attached to a modulus give an invertible change of coordinates of the vector
$(Z_a)_a$, so the same conclusions hold for the errors of any family of distinct primitive
characters whose conductors divide $q$ (their imprimitive Euler factors only add terms
$O_q((1+t)e^{-t/2})$). If $\sum_jw_{j,a}=0$ for every $a$, the central constants cancel, so
with $R_a(x)=S_a(x)-\log x/n_q-c_a$ the combination
$\sum_{j,a}w_{j,a}\sqrt{e^{h_j}x}\,R_a(e^{h_j}x)$ oscillates in both directions without any
information on central zeros; for instance $\sqrt{cx}\,R_a(cx)-\sqrt x\,R_a(x)$, for fixed
$c>1$, exceeds a positive constant and falls below its negative on sets of positive lower
logarithmic density.

\begin{example}\label{zer:ex:probes}
For $q=12$ and the four shifts $0,\log2,\log(3/2),\log5$ (scales $1,2,3/2,5$), the primes
$13,17,19,11$ in the classes $1,5,7,11$ are probes in the sense of \cref{zer:lem:probes}: each
product with an off-diagonal scale ratio is either not an integer or has two distinct prime
factors. \Cref{zer:thm:coercive} therefore gives rank $16$ in every band of the stated sizes.
For $q=3$ and scales $1,2,3/2$ the primes $7$ and $5$ serve the classes $1$ and $2$. These
exact rational checks are in the files
\begin{center}
\artifact{multiscale_prime_fluctuations/verify_examples.py}\\
\artifact{grh_multiscale_covariance/numerical_check.py}
\end{center}
\end{example}

\paragraph{Unsmoothed prime counts.}
For the ordinary counting function $\pi(x;q,a)$ the zero series converges only in a local mean, but
the positivity of \cref{zer:cor:multiscale} persists once the prime squares and the central zeros are
removed. The two statements do not imply each other: the smoothed and the unsmoothed vectors have
different coefficients, $1/(\tfrac14+\gamma^2)$ against $-1/(\tfrac12+i\gamma)$, and different
centering constants, because the Mellin denominator is now $s$ instead of $s(s-1)$. Stepanov almost
periodicity of normalized prime-counting remainders in $L^2$ is classical
\cite[Theorem~1(e)]{KaczorowskiRamare2003Almost}; the additions are an explicit local $L^4$ bound, the
transfer to $\pi(x;q,a)$, and the arithmetic separation of the scales.

With $v_a$ as in \eqref{zer:eq:va}, put $r_q(a)=\#\set{b\in G:b^2=a}$ and
\[
 \mathcal X_a(t)=te^{-t/2}\Bigl(\pi(e^t;q,a)-\frac{\li(e^t)}{n_q}\Bigr)+\frac{r_q(a)}{n_q}+2v_a(0),
 \qquad t\ge\log2,
\]
where $\li(x)=\int_2^x\dd u/\log u$; the term $r_q(a)/n_q$ removes the prime squares and $2v_a(0)$ the
possible central zeros, whose absence is not assumed. For distinct real $h_1,\dots,h_r$, $D=rn_q$ and
$\gamma\ne0$ put $\widetilde z(\gamma)_{j,a}=-e^{i\gamma h_j}v_a(\gamma)/(\tfrac12+i\gamma)$, and for a
finite set $J\subset\Gamma_q\setminus\set0$ with $J=-J$ let
\[
 \mathcal X_J(t)_{j,a}=\mathcal X_a(t+h_j)-\sum_{\gamma\in J}\widetilde z(\gamma)_{j,a}e^{i\gamma t},
\]
a real vector defined for $t\ge u_0$, where $u_0+h_j\ge\log2$ for every $j$.

\begin{theorem}[Unsmoothed multiscale covariance and occupation; \stG]\label{zer:thm:unsmoothed}
The series $\widetilde{\mathcal G}_J=\sum_{\gamma\in\Gamma_q\setminus(J\cup\set0)}\widetilde z(\gamma)\widetilde z(\gamma)^*$
converges to a real symmetric positive definite matrix, and uniformly in $u\ge u_0$, as $L\to\infty$,
\[
 \frac1L\int_u^{u+L}\mathcal X_J(t)\,\dd t\to0,\qquad
 \frac1L\int_u^{u+L}\mathcal X_J(t)\mathcal X_J(t)^{\mathsf T}\dd t\to\widetilde{\mathcal G}_J .
\]
There are $\eps,\eta,L_0>0$ such that for every real unit $w\in\R^D$, every $u\ge u_0$ and every
$L\ge L_0$,
\[
 \abs{\set{t\in[u,u+L]:w\cdot\mathcal X_J(t)\ge\eps}}\ge\eta L,\qquad
 \abs{\set{t\in[u,u+L]:w\cdot\mathcal X_J(t)\le-\eps}}\ge\eta L,
\]
so the closed convex hull of the essential values of $\mathcal X_J$ on each such interval contains the
closed ball of radius $\eps$ about $0$. Moreover, with $c,K,T_0>0$ depending only on $q$ and the
shifts,
\[
 \lambda_{\min}(\widetilde{\mathcal G}_J)\ge\frac{cH}{T^2\log(2T)}
\]
whenever $T\ge T_0$, $K\log(2T)\le H\le T$ and $T>\max_{\gamma\in J}\abs\gamma$.
\end{theorem}

In terms of $x=e^t$ the intervals are $[e^u,e^{u+L}]$ with the measure $\dd x/x$: every centered
linear contrast of the unsmoothed counts at the scales $e^{h_j}$ lies above $\eps$ and below $-\eps$ on
a positive proportion of every sufficiently long logarithmic interval. The constants may deteriorate
as shifts coalesce, and no uniformity in the number of coordinates is asserted. The proof needs a
fourth-moment form of \cref{zer:lem:bins}.

\begin{lemma}[Local fourth moments; \stU]\label{zer:lem:bins4}
Let $P(t)=\sum_\lambda c_\lambda e^{i\lambda t}$ be a finite sum with real frequencies and
$b_j=\sum_{\lambda\in[j,j+1)}\abs{c_\lambda}$. With an absolute constant $A$,
$\norm P_{S^4}\le A\norm b_{\ell^{4/3}}$, where
$\norm f_{S^4}=\sup_u(\int_u^{u+1}\norm{f(t)}^4\dd t)^{1/4}$. Consequently, if a locally finite family
of coefficients has $b\in\ell^{4/3}(\Z)$, its finite symmetric truncations converge in $S^4$.
\end{lemma}
\begin{proof}
\Cref{zer:lem:bins} with $h=1$ gives $\norm P_{S^2}\ll\norm b_{\ell^2}$. Apply it to $P^2$: adding
frequencies in the bins $j$ and $\ell$ gives a frequency in the bin $j+\ell$ or $j+\ell+1$, and
grouping coincident frequencies can only decrease absolute mass, so the bin masses of $P^2$ satisfy
$d_k\le(b*b)_k+(b*b)_{k-1}$. Young's inequality gives
$\norm P_{S^4}^4=\norm{P^2}_{S^2}^2\ll\norm{b*b}_{\ell^2}^2\le\norm b_{\ell^{4/3}}^4$. Apply this to
differences of truncations.
\end{proof}

By the local zero count, the bin masses of
\[
 F_a(t)=-\sum_{\substack{\gamma\in\Gamma_q\\\gamma\ne0}}\frac{v_a(\gamma)}{\tfrac12+i\gamma}e^{i\gamma t}
\]
are $O_q(\log(2+\abs j)/(1+\abs j))$, so the series converges in $S^4$; comparing the tails of
$\sum b_j^{4/3}$ and $\sum b_j^2$ with integrals, its truncation to $0<\abs\gamma\le Y$ has error
$O_q(Y^{-1/4}\log(2Y))$ in $S^4$ and $O_q(Y^{-1/2}\log(2Y))$ in $S^2$. No lower bound on gaps between
ordinates enters.

\begin{lemma}[The unsmoothed expansion; \stG]\label{zer:lem:unsmoothed}
For each $a\in G$, almost everywhere on $t\ge\log2$, $\mathcal X_a(t)=F_a(t)+R_a(t)$, where $R_a$ is
locally bounded, square integrable on $[\log2,\infty)$, has uniformly bounded local fourth moments, and
satisfies $\norm{R_a}_{L^4([U,U+1])}\ll_qU^{-1}$ for $U\ge2$.
\end{lemma}
\begin{proof}
For a primitive $\chi$, with $\delta_\chi=1$ in conductor one and $0$ otherwise, the explicit formula,
read as an identity of distributions on $t>0$ (integrate $-L'/L(s,\chi)$ against a smooth compactly
supported Mellin test and move the contour; \cref{src:thm:dgrh} makes each nontrivial zero term a
pure phase), is
\[
 e^{-t/2}\bigl(\psi(e^t,\chi)-\delta_\chi e^t\bigr)=-\sum_\gamma\frac{m_\chi(\gamma)e^{i\gamma t}}{\tfrac12+i\gamma}
 +O_\chi\bigl((1+t)e^{-t/2}\bigr)\qquad(t\ge\log2),
\]
the remainder collecting the constant, the pole at $s=0$ and the trivial zeros. The zero series
converges in $S^4$ by \cref{zer:lem:bins4}, hence in $L^1_{\mathrm{loc}}$, so the identity holds almost
everywhere for the arithmetic step function. Imprimitive characters lose prime powers of total size
$O_q(1+t)$ before normalization. Orthogonality, with the central term $-(\tfrac12)^{-1}v_a(0)$, gives
\[
 \begin{aligned}E_{\psi,a}(t)&:=e^{-t/2}\Bigl(\psi(e^t;q,a)-\frac{e^t}{n_q}\Bigr)\\
 &=F_a(t)-2v_a(0)+O_q\bigl((1+t)e^{-t/2}\bigr).\end{aligned}
\]
Prime squares contribute $\sum_{b^2=a}\vartheta(\sqrt x;q,b)$ to the difference
$\psi(x;q,a)-\vartheta(x;q,a)$ and higher powers $O_q(x^{1/3}\log^2(2x))$; the pointwise bound
$\vartheta(x;q,b)=x/n_q+O_q(\sqrt x\log^2(2x))$, used only for such deterministic remainders, gives
\[
 \begin{aligned}E_{\vartheta,a}(t)&:=e^{-t/2}\Bigl(\vartheta(e^t;q,a)-\frac{e^t}{n_q}\Bigr)\\
 &=E_{\psi,a}(t)-\frac{r_q(a)}{n_q}+O_q\bigl((1+t)^2e^{-t/6}\bigr).\end{aligned}
\]
Each coordinate is centered against $\li(x)/n_q$, the principal direction included, which is why the
correction is $r_q(a)/n_q$. Partial summation gives exactly
\[
 te^{-t/2}\Bigl(\pi(e^t;q,a)-\frac{\li(e^t)}{n_q}\Bigr)=E_{\vartheta,a}(t)+V_a(t)+\frac{2te^{-t/2}}{n_q\log2},\qquad
 V_a(t)=te^{-t/2}\int_{\log2}^t\frac{e^{v/2}E_{\vartheta,a}(v)}{v^2}\,\dd v .
\]
Split the integral at $v=t/2$. On the earlier part $E_{\vartheta,a}(v)\ll_q(1+v)^2$, which contributes
$O_q(t^2e^{-t/4})$. On the later part put $s=t-v$ and use $t/(t-s)^2\le4/t$: for $t\in[U,U+1]$,
$\abs{V_{a,\mathrm{late}}(t)}\le(4/U)\int_0^{(U+1)/2}e^{-s/2}\abs{E_{\vartheta,a}(t-s)}\,\dd s$, and
Minkowski's inequality with the bounded local fourth moments of $E_{\vartheta,a}$ gives
$\norm{V_a}_{L^4([U,U+1])}\ll_qU^{-1}$. The remainders are locally bounded, and
$\int_U^{U+1}R_a^2\le(\int_U^{U+1}R_a^4)^{1/2}\ll_qU^{-2}$ summed over integers $U$ gives square
integrability.
\end{proof}

\begin{lemma}[Occupation from a fourth moment; \stU]\label{zer:lem:occupation4}
Let $Y$ be a real vector-valued function with uniform long-window mean $0$ and uniform long-window
second moment $\mathcal G$, $\lambda=\lambda_{\min}(\mathcal G)>0$, and suppose that
$\abs I^{-1}\int_I\abs{w\cdot Y}^4\le M^4$ for every real unit $w$ and every sufficiently long allowed
interval $I$, with $M\ge1$. Then the two occupation inequalities of \cref{zer:thm:unsmoothed} hold
uniformly with
\[
 a=\frac{(\lambda/2)^{3/2}}{M^2},\qquad\eps=\frac a8,\qquad\eta=\Bigl(\frac a{8M}\Bigr)^2 .
\]
\end{lemma}
\begin{proof}
Let $\mathbb E$ be the window average and $f=w\cdot Y$. For long windows, uniformly in unit $w$,
$\mathbb Ef^2\ge\lambda/2$ and $\abs{\mathbb Ef}\le a/2$. H\"older's inequality between $L^1$ and $L^4$
gives $\mathbb E\abs f\ge(\mathbb Ef^2)^{3/2}/(\mathbb Ef^4)^{1/2}\ge a$, so $\mathbb Ef_\pm\ge a/4$. With
$\eps=a/8$, $a/8\le\mathbb E[f_+\mathbf 1_{f\ge\eps}]\le(\mathbb Ef^2)^{1/2}\mathbb P(f\ge\eps)^{1/2}
\le M\,\mathbb P(f\ge\eps)^{1/2}$, and likewise for $-f$.
\end{proof}

\begin{proof}[Proof of \cref{zer:thm:unsmoothed}]
By \cref{zer:lem:unsmoothed}, $\mathcal X_J=F_J+R$ with $F_J$ the $S^4$ limit of exponential sums with
coefficients $\widetilde z(\gamma)$, $\gamma\notin J\cup\set0$, square summable by the bin bound, and
$R\in L^2$ on a half-line. For a finite sum the long-window limits follow by expanding products,
since the mean of $e^{i\omega t}$ over $[u,u+L]$ has modulus at most $2/(L\abs\omega)$. If $P$
approximates $F_J$ in $S^2$, covering $[u,u+L]$ by at most $L+1$ unit intervals gives
$L^{-1}\int_u^{u+L}\norm{F_J-P}^2\le2\norm{F_J-P}_{S^2}^2$, and $L^{-1}\int_u^{u+L}\norm R^2\le
L^{-1}\norm R_{L^2}^2$; Cauchy--Schwarz then gives the two limits uniformly in $u$, with
$\widetilde{\mathcal G}_J$ real because $\widetilde z(-\gamma)=\overline{\widetilde z(\gamma)}$.

For positivity, the proof of \cref{zer:thm:coercive} shows, for $T\ge T_0$ and $K\log(2T)\le H\le T$,
that $\sum_{T<\gamma\le T+H}\abs{A_w(\gamma)}^2\ge\alpha^2H\norm w_2^2/(4DC_0\log(2T))$; it uses only the
fixed-argument Landau formula \eqref{zer:eq:landaufixed}, the probes of \cref{zer:lem:probes} and
Cauchy--Schwarz over the distinct ordinates of the band. For real $w$,
$w\cdot\widetilde z(\gamma)=-A_w(\gamma)/(\tfrac12+i\gamma)$, and pairing $\gamma$ with $-\gamma$ gives
$w^{\mathsf T}\widetilde{\mathcal G}_Jw\ge2\sum_{T<\gamma\le T+H}\abs{A_w(\gamma)}^2/(\tfrac14+\gamma^2)$ for
every band above $\max J$; with $\tfrac14+\gamma^2\le5T^2$ this is the eigenvalue bound. The probes
separate scaled prime powers, not individual zeros.

The Fourier part has bounded $S^4$ norm and the remainder bounded local fourth moments, so covering a
long interval by unit intervals gives a common $M$ in \cref{zer:lem:occupation4}, which gives the
occupation inequalities. If the closed convex hull of the essential values missed a point of the
$\eps$-ball, separation would give a unit direction with essential supremum less than $\eps$.
\end{proof}

\begin{corollary}[No exact finite linear explanation; \stG]\label{zer:cor:unsmoothedpred}
For a nonzero real array $w=(w_{j,a})$, a real finite trigonometric polynomial $P$ and a real
constant $b$, the long-window mean square of $\sum_{j,a}w_{j,a}\mathcal X_a(t+h_j)-P(t)-b$ has a
positive limit, uniformly in the origin. In particular no nonzero contrast at one scale is predicted
exactly in mean square by real linear combinations of the other scales and a finite trigonometric
correction.
\end{corollary}
\begin{proof}
Let $J$ consist of the nonzero frequencies of $P$ in $\Gamma_q$ and their negatives. The frequencies
outside $J$ keep the coefficients of $w\cdot\mathcal X_J$, and the mean-square computation in the proof
above, which allows extra frequencies of $P$, bounds the limit below by
$w^{\mathsf T}\widetilde{\mathcal G}_Jw>0$. For a predictor, collect all coefficients in one array, whose
block at the predicted scale is the given nonzero contrast.
\end{proof}

The lower bound is not uniform over all corrections, since longer Fourier truncations approximate the
error in mean square. Positive covariance and two-sided directional occupation do not force the
limiting values to fill an open set, so they do not give every ordering of the counts or inclusivity
of ordinary prime races.

\paragraph{A certificate from one isolated prime atom.}
The probe argument also has a distributional form. For a primitive character $\eta$ of conductor
$Q$, parity $\kappa$ and $\delta=1$ for conductor one, the distribution
$U_\eta(t)=\sum_\gamma m_\eta(\gamma)e^{i\gamma t}$ (pairing with test functions converges by the
zero count) satisfies, on $(0,\infty)$,
\begin{equation}\label{zer:eq:atoms}
 U_\eta=-\sum_{n\ge2}\frac{\Lambda(n)\eta(n)}{\sqrt n}\,\delta_{\log n}+g_\eta,\qquad
 g_\eta(t)=-\frac{e^{-(\kappa+1/2)t}}{1-e^{-2t}}+\delta\bigl(e^{-t/2}+e^{t/2}\bigr).
\end{equation}
Indeed the second logarithmic derivative of the completed function $\Xi_\eta$ is
$\sum_\gamma m_\eta(\gamma)(s-\tfrac12-i\gamma)^{-2}$, and for $\Re s>1$ it equals
$-\sum_n\Lambda(n)\eta(n)(\log n)n^{-s}-\sum_{k\ge0}(s+\kappa+2k)^{-2}+\delta(s^{-2}+(s-1)^{-2})$.
At $s=\tfrac12+c+iy$, $c>\tfrac12$, the function $(c+i(y-\gamma))^{-2}$ is the Fourier transform of
$\mathbf 1_{t>0}te^{-ct}e^{i\gamma t}$; testing against $e^{ct}\phi(t)/t$ with
$\phi\in C_c^\infty(0,\infty)$, and using
$\sum_\gamma m_\eta(\gamma)/(c^2+(y-\gamma)^2)=O(\log(2+\abs y))$, the exchanges are justified, and
dividing by $te^{-ct}$ gives \eqref{zer:eq:atoms}.

\begin{proposition}[Atom certificate; \stG]\label{zer:prop:atoms}
Let $J\subset\Gamma_q$ be finite. For each coordinate $l=(j,a)$ take an isolated probe
$p_l=p_a$ as in \cref{zer:lem:probes}, a bounded interval $I_l$ about
$t_l=\log p_l-h_j$ containing no other atom location of the translates
$U_{\chi^*}(\cdot+h_k)$ and with $t+h_k>0$ on it, and $\phi_l\in C_c^\infty(I_l)$ with
$\phi_l(t_l)=1$ that annihilates the functions $g_{\chi^*}(t+h_k)$ and $e^{i\gamma t}$,
$\gamma\in J\cup\set0$. Such tests exist, the numbers
\[
 K_l=\frac{p_l}{(\log p_l)^2}\sum_{\gamma\in\Gamma_q\setminus J}(1/4+\gamma^2)^2
 \Bigl|\int\phi_l(t)e^{i\gamma t}\dd t\Bigr|^2
\]
are finite, and
\[
 \lambda_{\min}\Bigl(\sum_{\gamma\in\Gamma_q\setminus(J\cup\set0)}z(\gamma)z(\gamma)^*\Bigr)\ge\Bigl(\sum_lK_l\Bigr)^{-1}.
\]
\end{proposition}
\begin{proof}
Existence: if evaluation at $t_l$ vanished on the common kernel of the finitely many
functionals $\phi\mapsto\int\phi f$, it would be a combination of them, and a point mass would be
a smooth function on $I_l$. For $v\in\C^D$ put $d_\gamma=\sum_lv_lz(\gamma)_l$ and
$c_{j,\chi}=n_q^{-1}\sum_av_{j,a}\overline{\chi(a)}$. The distribution
$\sum_{j,\chi}c_{j,\chi}U_{\chi^*}(t+h_j)$ has coefficient $(\tfrac14+\gamma^2)d_\gamma$ at each
$\gamma\ne0$. Pairing it with $\phi_l$ removes the frequencies in $J\cup\set0$ and the smooth part
of \eqref{zer:eq:atoms}, and leaves only the isolated atom, whose coefficient is
$-(\log p_l)p_l^{-1/2}\sum_\chi c_{j,\chi}\chi^*(p_l)=-(\log p_l)p_l^{-1/2}v_l$. Cauchy--Schwarz gives
$\abs{v_l}^2\le K_l\sum_{\gamma\notin J}\abs{d_\gamma}^2$; summing over $l$ bounds the Rayleigh
quotient. Finiteness of $K_l$ follows from the zero count and the rapid decay of
$\int\phi_le^{i\gamma t}\dd t$.
\end{proof}

\paragraph{Derivatives.}
Let $\chi_1,\dots,\chi_r$ be distinct primitive characters as in \cref{zer:thm:polyrigid}, with
$\delta_j=1$ for conductor one. For $t\ge0$ and integers $k\ge1$ define the logarithmic Riesz
means and their errors
\[
 S_{j,k}(t)=\frac1{k!}\sum_{n\le e^t}\Lambda(n)\chi_j(n)(t-\log n)^k,\qquad
 E_{j,k}(t)=e^{-t/2}\bigl(S_{j,k}(t)-\delta_je^t\bigr).
\]
Let $\mathcal G$ be the set of distinct ordinates of the family and fix $d\ge0$, $k\ge d+1$.

\begin{lemma}[Riesz-mean explicit formula; \stG]\label{zer:lem:riesz}
For $k\ge1$ and $t\ge0$,
$S_{j,k}(t)=\delta_je^t-\sum_\rho m_j(\rho)e^{\rho t}\rho^{-k-1}+O_{j,k}((1+t)^{k+1})$. Consequently,
for $0\le\ell\le d$,
\begin{equation}\label{zer:eq:jetexp}
 E_{j,k}^{(\ell)}(t)=-\sum_{\gamma\in\mathcal G}\frac{m_j(\gamma)(i\gamma)^\ell}{(1/2+i\gamma)^{k+1}}e^{i\gamma t}
 +O_{j,k,d}\bigl((1+t)^{k+1}e^{-t/2}\bigr),
\end{equation}
with absolutely and uniformly convergent series ($0^0=1$).
\end{lemma}
\begin{proof}
The Laplace transform of $S_{j,k}$ is $-L'(s,\chi_j)/(s^{k+1}L(s,\chi_j))$ for $\Re s>1$. Write
$-L'/L$ through \eqref{zer:eq:hadamard} and the digamma series: it is a constant, plus
$\delta/(s-1)$, plus a multiple of $1/s$, minus $\sum_\rho m(\rho)(1/(s-\rho)+1/\rho)$, minus
$\sum_{b_n>0}(1/(s+b_n)-1/b_n)$ with $b_n=a+2n$. For $z\ne0$ the inverse Laplace transform of
$(1/(s-z)+1/z)s^{-k-1}$ is $\bigl(e^{zt}-\sum_{h<k}(zt)^h/h!\bigr)z^{-k-1}$. Since
$\sum\abs z^{-2}$ converges over both families of $z$, and $\Re z\le\tfrac12$, termwise inversion
is justified by uniqueness of Laplace transforms of continuous functions of exponential
order. The pole terms give $\delta(e^t-\sum_{h\le k}t^h/h!)$ and polynomials of degree at most
$k+1$, the polynomial parts are $O((1+t)^{k-1})$ and the trivial zeros are bounded. Since
$S_{j,k}'=S_{j,k-1}$ for $k\ge2$, one has
$E_{j,k}^{(\ell)}=\sum_{h\le\ell}\binom\ell h(-\tfrac12)^{\ell-h}E_{j,k-h}$; inserting the formula
at each order $k-h\ge1$ turns the zero numerator into $(\rho-\tfrac12)^\ell=(i\gamma)^\ell$ by
\cref{src:thm:dgrh}, and $k+1-\ell\ge2$ gives absolute convergence.
\end{proof}

For a cut $H\ge0$ remove the low frequencies completely:
\[
 Y_{j\ell,H}(t)=E_{j,k}^{(\ell)}(t)+\sum_{\abs\gamma\le H}
 \frac{m_j(\gamma)(i\gamma)^\ell}{(1/2+i\gamma)^{k+1}}e^{i\gamma t},
\]
and stack these coordinates in $Y_H(t)\in\C^N$, $N=r(d+1)$.

\begin{theorem}[High-pass jet covariance; \stG]\label{zer:thm:jets}
For every finite family of distinct primitive characters, every $d\ge0$, $k\ge d+1$ and
$H\ge0$, uniformly in $u\ge0$,
\[
 \lim_{L\to\infty}\frac1L\int_u^{u+L}Y_H(t)Y_H(t)^*\dd t
 =\mathcal G_H:=\sum_{\abs\gamma>H}\frac{v(\gamma)v(\gamma)^*}{\abs{1/2+i\gamma}^{2k+2}},
\]
with $v(\gamma)$ as in \cref{zer:cor:tail}, and $\mathcal G_H$ is positive definite. The mean of
$Y_H$ tends to zero in the same sense. The same holds if any finite set of ordinates, zero
included, is removed instead of a full band.
\end{theorem}
\begin{proof}
By \eqref{zer:eq:jetexp}, $Y_H$ is an absolutely convergent series with frequencies $\abs\gamma>H$
plus a bounded integrable remainder, so the complex analogue of \cref{zer:lem:ap}(i) gives the
limits. For $c\ne0$, $c^*\mathcal G_Hc=\sum_{\abs\gamma>H}\abs{c^*v(\gamma)}^2\abs{1/2+i\gamma}^{-2k-2}>0$ by
\cref{zer:cor:tail}. A finite deletion leaves all sufficiently high frequencies.
\end{proof}

In particular, for coefficients $c_{j\ell}$ not all zero and any finite trigonometric polynomial
$P(t)$ with real frequencies, the mean square of $\sum c_{j\ell}E_{j,k}^{(\ell)}-P$ has a positive
limit: choose $H$ above the frequencies of $P$. No finite set of frequencies explains a
nonzero combination of the smoothed errors and their derivatives. If $\gamma_1,\dots,\gamma_N$
are tail ordinates whose vectors form a basis and $W$ is the matrix with columns
$v(\gamma_h)(1/2+i\gamma_h)^{-k-1}$, then $\mathcal G_H\succeq WW^*$ and
\begin{equation}\label{zer:eq:jetcert}
 \lambda_{\min}(\mathcal G_H)\ge\frac{\abs{\det W}^2}{(\tr WW^*)^{N-1}}>0,
\end{equation}
since the product of the $N-1$ largest eigenvalues of $WW^*$ is at most $(\tr WW^*)^{N-1}$;
computing this bound needs certified multiplicities.

For residue classes take the characters modulo $q$ with their distinct primitive inducing
characters, and for $b\in G$ put
\[
 E^{[q]}_{b,k}(t)=e^{-t/2}\Bigl(\frac1{k!}\sum_{n\le e^t,\ n\equiv b}\Lambda(n)(t-\log n)^k-\frac{e^t}{n_q}\Bigr).
\]
Passing to primitive characters changes $S_k$ by $O_q((1+t)^{k+1})$, since each prime $p\mid q$
has at most $t/\log p$ powers of weight at most $t^k\log p/k!$. Let $Z_H(t)\in\R^{n_q(d+1)}$ be the
vector of $(E^{[q]}_{b,k})^{(\ell)}$ with the low-frequency terms
$n_q^{-1}\sum_\chi\overline{\chi(b)}\sum_{\abs\gamma\le H}m_{\chi^*}(\gamma)(i\gamma)^\ell(1/2+i\gamma)^{-k-1}e^{i\gamma t}$
added; it is real, because conjugate characters have reflected conjugate terms and the cut is
symmetric. Up to a decaying error it is $MY_H$ with $M=(\overline{\chi(b)}/n_q)\otimes I_{d+1}$ and
$MM^*=I/n_q$, so its covariance $M\mathcal G_HM^*$ is real symmetric positive definite.

\begin{corollary}[Uniform directional oscillation; \stG]\label{zer:cor:jetosc}
For fixed $q,k,d,H$ as above there are $\eta>0$ and $T_0$ such that for every $u\ge0$, every
$T\ge T_0$ and every real unit vector $c$, there are $t_\pm\in[u,u+T]$ with
$\pm c^{\mathsf T}Z_H(t_\pm)\ge\eta$. Equivalently, the convex hull of
$\set{Z_H(t):u\le t\le u+T}$ contains the closed ball of radius $\eta$ about $0$.
\end{corollary}
\begin{proof}
Apply \cref{zer:lem:ap}(iii); the remainder in \eqref{zer:eq:jetexp} is bounded and
integrable on $[0,\infty)$. The converse direction of the equivalence follows by maximizing a
linear functional over the hull.
\end{proof}

For $d=0$ and $H=0$ every nonzero real contrast of centred smoothed residue-class counts
changes sign with bounded waiting time on the logarithmic scale. A covariance statement does
not say that every ordering of the unsmoothed counts occurs, nor that the limiting law has full
support: a distribution on a curve can have positive definite covariance and miss open sets.

\paragraph{A filter computed from primes alone.}
The constants $c_a$ and $\beta_a$ require $L'/L(1,\chi)$ and the central orders. A compact filter
removes them. Fix $h>0$ and put
\[
 (A_hf)(t)=\frac1h\int_0^hf(t-u)\,\dd u,\qquad d_h=\frac{2(1-e^{-h/2})}h\in(0,1),
\]
\[
 \mathcal D_a(t)=(I-A_h)(d_hI-A_h)^2\bigl[e^{t/2}S_a(e^t)\bigr].
\]
Only prime powers up to $e^t$ enter $\mathcal D_a(t)$, and only the values on $[t-3h,t]$. The
operator sends $e^{zt}$ to $M_h(z)e^{zt}$ with
\[
 a_h(z)=\frac{1-e^{-hz}}{hz},\quad a_h(0)=1,\qquad M_h(z)=(1-a_h(z))(d_h-a_h(z))^2 .
\]
Since $M_h(0)=M_h(\tfrac12)=M_h'(\tfrac12)=0$ (note $a_h(\tfrac12)=d_h$), it annihilates $1$,
$e^{t/2}$ and $te^{t/2}$.

\begin{lemma}\label{zer:lem:multiplier}
For $h>0$ and real $\gamma\ne0$, $M_h(i\gamma)\ne0$.
\end{lemma}
\begin{proof}
$a_h(i\gamma)$ is the mean of $e^{-i\gamma u}$ over $[0,h]$, of modulus less than $1$, so
$a_h(i\gamma)\ne1$. If it is real, its imaginary part $-(1-\cos\gamma h)/(\gamma h)$ vanishes, so
$e^{i\gamma h}=1$ and $a_h(i\gamma)=0\ne d_h$.
\end{proof}

\begin{theorem}[Filtered prime races; \stG]\label{zer:thm:filtered}
For every $q\ge1$ and $h>0$ there are $\eps,\eta,L_0>0$ and $T_0$, depending only on $q,h$, such
that for every real unit $w\in\R^G$, every $T\ge T_0$ and every $L\ge L_0$,
\[
 \abs{\set{t\in[T,T+L]:w\cdot\mathcal D(t)\ge\eps}}\ge\eta L,\qquad
 \abs{\set{t\in[T,T+L]:w\cdot\mathcal D(t)\le-\eps}}\ge\eta L,
\]
and the convex hull of $\set{\mathcal D(t):T\le t\le T+L}$ contains the ball of radius $\eps$. The
same holds, with constants depending also on the deleted set, after subtracting from $\mathcal D$
the terms $M_h(i\gamma)v_a(\gamma)(\tfrac14+\gamma^2)^{-1}e^{i\gamma t}$ of any finite symmetric set of
nonzero ordinates.
\end{theorem}
\begin{proof}
By \cref{zer:lem:ramp},
$e^{t/2}S_a(e^t)=te^{t/2}/n_q+c_ae^{t/2}+\beta_a+\sum_{\gamma\ne0}v_a(\gamma)(\tfrac14+\gamma^2)^{-1}e^{i\gamma t}+O((1+t)e^{-t/2})$.
The filter has operator norm at most $2(1+d_h)^2$ on bounded functions and memory $3h$; it
removes the first three terms and multiplies each frequency by $M_h(i\gamma)$. So
$\mathcal D(t)=\sum_{\gamma\ne0}M_h(i\gamma)u_\gamma e^{i\gamma t}+O((1+t)e^{-t/2})$ with
$u_\gamma=(v_a(\gamma)/(\tfrac14+\gamma^2))_a$. If
$\sum_{\abs\gamma>H}\abs{M_h(i\gamma)}^2\abs{w\cdot u_\gamma}^2=0$ for a real $w$, then by
\cref{zer:lem:multiplier}
$\sum_\chi\bigl(\sum_aw_a\overline{\chi(a)}\bigr)m_\chi(\gamma)=0$ for $\abs\gamma>H$, and
\cref{zer:thm:polyrigid} with constant polynomials gives $w=0$. \Cref{zer:lem:ap}(ii) completes
the proof.
\end{proof}

For evaluation, $e^{t/2}S_a(e^t)=\sum_{n\equiv a}\Lambda(n)n^{-1/2}b(t-\log n)$ with
$b(v)=\mathbf 1_{v\ge0}(e^{v/2}-e^{-v/2})$, and the filter expands as
$d_h^2I-(d_h^2+2d_h)A_h+(1+2d_h)A_h^2-A_h^3$ with
$A_h^jb(v)=h^{-j}\sum_{k\le j}(-1)^k\binom jkb_j(v-kh)$, where $b_j$ is the $j$-fold primitive
of $b$ vanishing on $v<0$,
\[
 b_j(v)=2^j\Bigl(e^{v/2}-\sum_{k<j}\frac{(v/2)^k}{k!}\Bigr)
 -(-2)^j\Bigl(e^{-v/2}-\sum_{k<j}\frac{(-v/2)^k}{k!}\Bigr)\quad(v\ge0).
\]
For $v\ge3h$ the filtered kernel is exactly $-M_h(-\tfrac12)e^{-v/2}$, so each prime power far
below the cutoff contributes a decaying term and no growing terms cancel. Floating-point
samples for $q=1,4,5,8$ are in \artifact{grh_filtered_prime_races/experiment_output.txt}.





\subsection{Balancing all prime-power races in one multiplicative interval}\label{zer:sec:balancing}
\sources{grh_cm_four_terms_and_prime_bias; grh_prime_error_balancing; multiscale_prime_fluctuations}

Fix a positive integer $q$, put $G=(\Z/q\Z)^\times$ and $n_q=\abs G$, and let $\chi$ run
over the characters of $G$. Write $\chi^*$ for the primitive character inducing $\chi$ and
$m_\chi$ for the order of vanishing of $L(s,\chi^*)$ at $1/2$. For $a\in G$ define
\begin{equation}\label{zer:eq:racevec}
\begin{aligned}
 \psi(x;q,a)&=\sum_{n\le x,\ n\equiv a\bmod q}\Lambda(n),&
 b_a&=\frac2{n_q}\sum_\chi\overline{\chi(a)}\,m_\chi,\\
 D_a(x)&=\frac{\psi(x;q,a)-x/n_q}{\sqrt x}+b_a,&
 D(x)&=(D_a(x))_{a\in G}.
\end{aligned}
\end{equation}
The constants $b_a$ are real by conjugation. They incorporate possible central zeros, so
even characters with unknown central nonvanishing are included. The von Mangoldt weight
counts prime powers as well as primes.

\begin{theorem}[Fixed-interval balancing; \stG]\label{zer:thm:balancing}
There exist $C_q>1$ and $r_q>0$ such that, for every sufficiently large real $X$,
\begin{equation}\label{zer:eq:ballrace}
 \set{y\in\R^{n_q}:\norm y_2\le r_q}\ \subseteq\
 \zerconv\set{D(n):\ n\in\Z,\ X\le n\le C_qX}.
\end{equation}
Consequently every such $y$, including zero, is exactly a convex combination of at most
$n_q+1$ integer samples from that interval, with the same sample points and weights for
every residue class, and every unit linear combination of the $n_q$ coordinates attains
values at least $r_q$ and at most $-r_q$ in the interval.
\end{theorem}

The theorem gives existence with constants depending on the fixed modulus. It requires
neither simplicity of zeros nor rational independence of their ordinates. A ball in the
convex hull allows thresholds to be mixed; a prescribed ranking of raw prime counts asks
for all of its inequalities to hold at one threshold, which is a stronger demand (even a
circle has a convex hull containing points that the circle never visits). For the
single-threshold problem under hypotheses on the zero ordinates see
\cite{MartinNg2020Inclusive}.

The first step is a vector version of the local averaging in the source proof, where high
Laplace moments find a one-sided small value on a window (\code{IsPolar.eventually_moment_le}
and \code{IsPolar.exists_window_sub_mul_le} in
\code{New/Endgame/OneSidedCriterion.lean}). A common choice of positive weights replaces a
separate choice of a point for each direction.

\begin{lemma}[Laplace averaging; \stU]\label{zer:lem:laplace}
Let $E:[0,\infty)\to\R^n$ be measurable, with $\norm{E(v)}\le Me^{Av}$, $A\ge0$. Suppose
its coordinate Laplace transforms $\Phi(z)=\int_0^\infty E(v)e^{-zv}\,\dd v$, initially
defined for $\Re z>A$, extend holomorphically to $\Re z>0$ and to a neighborhood of $0$.
For every $\tau>0$ there is $\delta>0$ such that, for all large $T$, some probability
density supported on $[T,(1+\tau)T]$ has $E$-average of norm $O(e^{-\delta T})$. The same
estimate holds for some convex combination of at most $n+1$ values of $E$ in the interval.
\end{lemma}
\begin{proof}
Take $c=1+\tau/2$. Choose $\lambda>A$ sufficiently large as specified below, set
$k=\floor{\lambda cT}$, and let $V$ have density $\lambda^{k+1}v^ke^{-\lambda v}/k!$ on
$v>0$. The disk $\abs{z-\lambda}\le\lambda$ meets the imaginary axis only at $0$.
Compactness gives $R>\lambda$ with $\Phi$ holomorphic on a neighborhood of the closed
concentric disk of radius $R$. Cauchy's estimate yields
\[
 \norm{\mathbb E\,E(V)}=\lambda^{k+1}\norm{\Phi^{(k)}(\lambda)}/k!\ll\lambda(\lambda/R)^k .
\]
For the tails, exponential tilting gives
\[
 \mathbb E\bigl[e^{AV};\ V\notin I\bigr]=\Bigl(\frac\lambda{\lambda-A}\Bigr)^{k+1}
 \Pr(W\notin I),\qquad I=[T,(1+\tau)T],
\]
where $W$ has the same shape $k+1$ and rate $\lambda-A$. Its mean divided by $T$ tends to
$c\lambda/(\lambda-A)$. For all sufficiently large $\lambda$ this limit stays a fixed
positive distance from both endpoints $1,1+\tau$. The Chernoff bound for a gamma variable
at $y$ times its mean is $\exp[-(k+1)(y-1-\log y)]$ for the corresponding tail, by
optimizing its moment generating function. The two endpoint ratios stay bounded away from
$1$, so the probability is at most $2e^{-\kappa_\tau(k+1)}$ with $\kappa_\tau>0$
independent of all sufficiently large $\lambda,T$. Choose $\lambda$ large enough that
$\log(\lambda/(\lambda-A))<\kappa_\tau/2$. Truncating the density to $I$ thus changes its
vector integral by $O(e^{-\delta T})$, and its mass is $1-O(e^{-\delta T})$, after
decreasing $\delta$ if necessary. The normalized integral belongs to the closed convex
hull. Approximate it to within $e^{-\delta T}$ by a finite convex combination and apply
Carath\'eodory's theorem.
\end{proof}

\begin{corollary}[Exponentially accurate balance in log-windows; \stG]\label{zer:cor:expbalance}
For every $\tau>0$ there is $\delta>0$ such that for all large $T$ some convex combination
of at most $n_q+1$ values $D(e^v)$, $v\in[T,(1+\tau)T]$, has norm $O_q(e^{-\delta T})$.
\end{corollary}
\begin{proof}
For the character coordinate of \eqref{zer:eq:racevec}, write
$\psi(x,\chi)=\sum_{n\le x}\chi(n)\Lambda(n)$, $e_\chi=1$ for the principal character and
$0$ otherwise, and $s=z+1/2$. The Laplace transform of
$e^{-v/2}(\psi(e^v,\chi)-e_\chi e^v)+2m_\chi$ is
\[
 -\frac1s\frac{L'}L(s,\chi)-\frac{e_\chi}{s-1}+\frac{2m_\chi}z .
\]
The pole at $s=1$ cancels for the principal character. The residue at $z=0$ before adding
the final term is $-2m_\chi$. Finite Euler factors removed from primitive $L$-functions
have zeros only on $\Re s=0$. \Cref{src:thm:dgrh} therefore gives the continuation required
by \cref{zer:lem:laplace}, and the finite Fourier transform on $G$ gives all the coordinates
in \eqref{zer:eq:racevec}.
\end{proof}

For exact balance in a log-window of fixed length the individual nonzero frequencies are
used, through \cref{zer:thm:rank}.

\begin{proof}[Proof of \cref{zer:thm:balancing}]
For a nonzero ordinate $\gamma$ in the union of the spectra put
$\rho_\gamma=\tfrac12+i\gamma$ and
\[
 c_{\gamma,a}=-\frac1{n_q\rho_\gamma}\sum_\chi\overline{\chi(a)}\,M(\rho_\gamma)_\chi,
 \qquad c_\gamma=(c_{\gamma,a})_{a\in G},
\]
with $M(\rho)$ as in \cref{zer:thm:rank}. The finite Fourier transform and
\cref{zer:thm:rank} show that the vectors $c_\gamma$ for positive $\gamma$ span
$\C^{n_q}$. Thus the real and imaginary parts of finitely many of them span $\R^{n_q}$.

On $v>1$, the explicit formula, as an identity of distributions, reads
\begin{equation}\label{zer:eq:explicitvec}
 D(e^v)=\sum_{\gamma\ne0}c_\gamma e^{i\gamma v}+O_q\bigl((1+v)e^{-v/2}\bigr).
\end{equation}
To see the error term, first use a primitive character. The main pole gives $e_\chi e^v$,
each nontrivial zero contributes $-e^{\rho v}/\rho$ with multiplicity, and the central
contribution is canceled by the constants in \eqref{zer:eq:racevec}. The pole of the
Perron kernel at zero and a possible trivial zero there contribute $O_q(1+v)$ before
normalization. The other trivial-zero terms form a geometrically decreasing series for
$v>1$. Replacing a primitive character by the one modulo $q$ changes $\psi$ by
$O_q(v+1)$, since only powers of primes dividing $q$ are affected. The choice at a jump is
immaterial to integration. For the boundary-value formula and its convergence for smooth
averages see \cite[Theorem~1(d),(e)]{KaczorowskiRamare2003Almost}.

Take a smooth probability density $w$, positive on $(0,1)$ and supported on $[0,1]$, and
let $w_n$ be its $n$-fold convolution, supported on $[0,n]$. Put
$\widehat w(t)=\int w(y)e^{ity}\,\dd y$. For each nonzero real $t$,
$\abs{\widehat w(t)}<1$, and $\widehat w$ decreases rapidly. For a fixed $\omega$,
integration of \eqref{zer:eq:explicitvec} shows
\begin{multline*}
 \int_0^nw_n(y)\,D(e^{V+y})\,e^{-i\omega(V+y)}\,\dd y\\
 =\sum_{\gamma\ne0}c_\gamma e^{i(\gamma-\omega)V}\,\widehat w(\gamma-\omega)^n
 +O_q\bigl((1+V)e^{-V/2}\bigr)\qquad(V\ge2).
\end{multline*}
The error is uniform in $n$ since $(1+v)e^{-v/2}$ decreases for $v\ge2$. The sum converges
absolutely for $n\ge1$ by rapid decay and the classical $O_q(Y\log(2Y))$ zero count up to
height $Y$. The union of the spectra is locally finite. Hence for each fixed $\omega$ the
supremum of $\abs{\widehat w(\gamma-\omega)}$ over $\gamma\ne\omega$ is strictly below $1$.
All off-frequency terms tend uniformly in $V$ to zero as $n\to\infty$. At $\omega=0$ the
sum therefore tends to zero. At a selected $\omega=\gamma$ it tends to $c_\gamma$. For the
arithmetic remainder take first a sufficiently large fixed $n$, then $V$ sufficiently
large.

Average $D(e^v)$ against the nonnegative densities proportional to
$w_n(v-V)(1\pm\cos(\gamma v))$ and $w_n(v-V)(1\pm\sin(\gamma v))$. Their masses tend
uniformly to $1$ because $\abs{\widehat w(\gamma)}<1$. The averages approach, to any
prescribed accuracy, the vertices $\pm\Re c_\gamma$ and $\pm\Im c_\gamma$ for the selected
finite set of positive ordinates (the opposite signs for sine merely interchange the two
vertices). The convex hull of these ideal vertices contains a ball about the origin because
their real span has full dimension. Concretely, if its inradius is $r>0$, approximating
every vertex within $r/4$ preserves a ball of radius $3r/4$, by the support-function
characterization of convex inclusion.

Finally replace $D(e^v)$ in the averages by $D(\floor{e^v})$. Between consecutive integers
the arithmetic sum in the numerator is constant, and $\psi(x;q,a)\le x\log(2x)$. The mean
value theorem bounds the change by $O_q((1+v)e^{-v/2})$, uniformly for large $v$. The
averages now belong to the convex hull of a finite set of integer samples. Set
$V=\log(2X)$, so those integers lie in $[X,2e^nX]$. Taking $C_q=2e^n$ and $r_q=r/2$ proves
\eqref{zer:eq:ballrace}, and Carath\'eodory's theorem gives the sample bound.
\end{proof}

\begin{remark}\label{zer:rem:stepanov}
The conditional almost-periodicity theory of \cite[Theorem~1(e)]{KaczorowskiRamare2003Almost}
gives a second perspective. The corrected vector has Stepanov almost-periodic coordinates,
zero mean after removing the central terms, and the nonzero Fourier coefficients displayed
in \eqref{zer:eq:explicitvec}, up to an overall sign. \Cref{zer:thm:rank} makes every
nonzero real projection nonconstant. Positive and negative values of a nonzero mean-zero
projection then recur on sets of positive measure in intervals of bounded log-length: an
interval on which it exceeds a positive level on a positive-measure subset can be repeated
using relatively dense Stepanov almost-periods, and the same applies to its negative. The
positive compact averages in the proof give a direct argument uniform across directions.
\end{remark}

\begin{example}[\stU]\label{zer:ex:triangle}
Take $q=4$. Its two central orders vanish: the nonprincipal central value is a positive
alternating series, and for the principal character $\zeta(1/2)=\eta(1/2)/(1-\sqrt2)$,
where the alternating eta series is positive. The three integers $293,283,228$ lie in
$[200,400]$. For $v(n)=(\psi(n;4,1)-n/2,\ \psi(n;4,3)-n/2)$ the certified lower bounds for
the cyclic determinants are
\[
 \det(v(293),v(283))\ge5.625,\quad\det(v(283),v(228))\ge41.693,\quad
 \det(v(228),v(293))\ge14.050 .
\]
All three are positive, and the identity
$\det(v_2,v_3)v_1+\det(v_3,v_1)v_2+\det(v_1,v_2)v_3=0$ therefore gives strictly positive
balancing weights. On replacing $v(n)$ by $D(n)=v(n)/\sqrt n$, multiply the associated
weight by $\sqrt n$ and renormalize. This example is independent of GRH. The three bounds
are verified with integer sieving and rational estimates for logarithms in
\begin{center}
\artifact{grh_cm_four_terms_and_prime_bias/checks/prime_race_triangle.py}
\end{center}
\end{example}

\paragraph{Smoothed sums in bounded logarithmic windows.}
For the ramp sums of \cref{zer:sec:cov} the balance holds at real scales in windows of bounded
logarithmic length centred anywhere. With the notation \eqref{zer:eq:va} put
$R_a(x)=S_a(x)-\log x/n_q-c_a$, $F(v)=(e^{v/2}R_a(e^v))_{a\in G}$ and $\beta=(\beta_a)_{a\in G}$.

\begin{corollary}[Uniform balancing of smoothed errors; \stG]\label{zer:cor:smoothbalance}
For each $q$ there are $H>0$, $r>0$ and $T_0$ such that for every real $T\ge T_0$
\[
 \beta+\set{y\in\R^{n_q}:\norm y_2\le r}\ \subseteq\ \zerconv\set{F(v):T-H\le v\le T+H}.
\]
Hence every target $y$ with $\norm{y-\beta}\le r$, in particular $y=\beta$, equals
$\sum_{j\le k}w_j\sqrt{x_j}R_a(x_j)$ simultaneously for all $a\in G$, for some $k\le n_q+1$,
scales $x_j\in[e^{T-H},e^{T+H}]$ and convex weights $w_j$.
\end{corollary}
\begin{proof}
$F(v)-\beta=Z(v)$. Apply \cref{zer:cor:multiscale} with $r=1$, $h_1=0$ and $J=\varnothing$,
and \cref{zer:lem:ap}(ii) with windows of length $2H\ge L_0$; finish with Carath\'eodory's
theorem.
\end{proof}

The scales and weights depend on the centre and the target, and no single scale balances all
classes; for $q=1$ the intermediate value theorem gives one scale for every target in the
ball. This ball concerns the smoothed sums at real scales and \cref{zer:thm:balancing} the
unsmoothed $\psi(x;q,a)$ at integers; neither statement implies the other. The radius can be
certified from finite data: for a finite symmetric set $E$ of nonzero ordinates with tail
$\tau=\sum_{\gamma\notin E}\norm{c(\gamma)}$, where $c(\gamma)=(v_a(\gamma)/(\tfrac14+\gamma^2))_a$ and
$K=\sum_\gamma\norm{c(\gamma)}$, the uniform error of the window mean is at most
$\tau+H^{-1}\sum_{\gamma\in E}\norm{c(\gamma)}/\abs\gamma$, and that of the second moment is at most
$3K\tau+H^{-1}\sum_{\gamma,\eta\in E,\ \gamma+\eta\ne0}\norm{c(\gamma)}\norm{c(\eta)}/\abs{\gamma+\eta}$; any
certified bound for $K$ and positive lower bound for the least eigenvalue may be inserted in
the proof of \cref{zer:lem:ap}.

Averaging rules that are independent of the centre and of the target, and cancel all ordinates up
to a height, are in \cref{fil:sec:primes}: the simultaneous rule of \cref{rec:thm:polylogavg} and the
explicit sinc-product kernel of \cref{zer:thm:kernel}.

\begin{remark}[The remainder of the source proof]\label{zer:rem:sourceremainder}
Let $d\equiv7\pmod8$ be prime, $(M,d)=1$, and let $R_{d,M}$ be the centred remainder of
$\sum_{n\le x,\ n\equiv1\bmod M}\Lambda(n)(1-\chi_d(n))(1/n-1/x)$, whose main term is
$\log x/\varphi(M)$; its constant is the regular part at $1$ of the weighted logarithmic
derivative minus the pole coefficient, as in the source's \code{logDerivMean}. Take $q=Md$ and
$w(a)=\mathbf 1_{a\equiv1\bmod M}(1-\chi_d(a))$ on $(\Z/q\Z)^\times$; then $\sum_aw(a)=d-1$, so
$\sum_aw(a)/n_q=1/\varphi(M)$. The powers of $d$ add
$\log d\sum_{d^k\le x,\ d^k\equiv1\bmod M}(d^{-k}-x^{-1})$, a constant plus $O_{d,M}((1+\log x)/x)$, so
\[
 \sqrt xR_{d,M}(x)=\sum_aw(a)F_a(\log x)+O_{d,M}\bigl((1+\log x)/\sqrt x\bigr).
\]
Since $w\ne0$, \cref{zer:cor:multiscale} applies to the shifted copies
$(\sqrt{e^{h_j}x}\,R_{d,M}(e^{h_j}x)-\sum_aw(a)\beta_a)_j$: their covariance is positive definite,
every nonzero real combination has both signs on sets of positive lower logarithmic density,
and the centred remainder of the source has infinitely many spectral frequencies and is not a
linear combination of finitely many other scales. If $\varphi(M)>1$, the weights of $R_{d,M}$ and
of $R_{d,1}$ (computed on the same modulus, with the finitely many powers of primes dividing $M$
added back) are linearly independent, so the pair $(\sqrt xR_{d,M},\sqrt xR_{d,1})$ has a
two-dimensional ball about its central biases in every sufficiently late window of the fixed
logarithmic length of \cref{zer:cor:smoothbalance}, and both biases are attained by one convex
average of at most three scales.
\end{remark}

\subsection{Positive spectral weights: peaks, sharp windows and resonance}\label{zer:sec:peaks}
{\emergencystretch=6em\sources{prime_peak_rigidity; sparse_grh_balanced_windows; sharp_multiscale_prime_windows; grh_prime_resonance_2026-11-26; prime_endpoints_and_symplectic_compressions}}

On the critical line $1/(\rho(1-\rho))=1/(\tfrac14+\gamma^2)$ is positive. In the identity class,
and for suitably balanced or autocorrelated windows, all spectral coefficients of a smoothed
prime error are then nonnegative. Such a function is bounded by its value at $0$, which is an
arithmetic constant, and returns arbitrarily close to it. This gives exact extremal values,
rigidity near the extremum, sharp inequalities between windows, and a covariance kernel whose
local structure is arithmetic. We keep the notation of \cref{zer:sec:cov}; $K_\chi$ depends only
on $\chi^*$.

\begin{lemma}[Positive spectra; \stU]\label{zer:lem:posspec}
\begin{enumerate}[label=\textup{(\roman*)}]
\item Let $F_1,\dots,F_k:\R\to\C^N$ be absolutely convergent trigonometric series and $\eps>0$. The
set of $\tau$ with $\sup_{u\in\R}\norm{F_i(u+\tau)-F_i(u)}<\eps$ for all $i$ is relatively dense
and unbounded above.
\item Let $f(t)=\sum_\gamma c_\gamma e^{i\gamma t}$ with $c_\gamma\ge0$ and $\sum c_\gamma<\infty$.
Then $\abs{f(t)}\le f(0)$ for all $t$. If $t_j\to\infty$ and $t_{j+1}-t_j\to0$, then
$\limsup_j\Re f(t_j)=f(0)$.
\item Let $\widehat G$ be the characters of a finite abelian group $G$, let
$w_{\chi,\gamma}\ge0$ be summable, and put
\[
 \mathcal K(a,t)=\sum_{\chi\in\widehat G}\sum_\gamma w_{\chi,\gamma}\overline{\chi(a)}e^{i\gamma t}.
\]
For $(a_i,t_i)\in G\times\R$ and $c_i\in\C$ put
\[
 Q=\sum_{i,j}c_i\overline{c_j}\mathcal K(a_ia_j^{-1},t_i-t_j),\qquad
 \mathcal S(T)=\sum_{i,j}c_i\overline{c_j}\mathcal K(a_ia_j^{-1},T+t_i-t_j).
\]
Then $Q\ge0$, $\abs{\mathcal S(T)}\le Q$, and $\limsup_{T\to\infty}\abs{\mathcal S(T)}=Q$.
\end{enumerate}
\end{lemma}
\begin{proof}
(i) Discard tails of total coefficient norm below $\eps/4$. For the finitely many remaining
frequencies, the closure of $\set{(e^{i\gamma t})_\gamma:t\in\R}$ is a compact subgroup $\mathcal T$
of a torus. For a symmetric neighbourhood $U$ of the identity, the translates by the dense
orbit cover $\mathcal T$, so finitely many, by $s_1,\dots,s_m\in[A,B]$, already do; for every $t$
some $t-s_j$ has its phases in $U$, so returns occur in every interval of length $B-A+1$, and
arbitrarily far to the right. No independence of the frequencies is used. (ii) The bound is
the triangle inequality. By (i) there are arbitrarily large $\tau$ with $\abs{f(\tau)-f(0)}$
small, and uniform continuity transfers this to a mesh with gaps tending to zero.
(iii) $\mathcal S(T)=\sum_{\chi,\gamma}w_{\chi,\gamma}e^{i\gamma T}\abs{\sum_ic_i\overline{\chi(a_i)}e^{i\gamma t_i}}^2$
has nonnegative coefficients and $\mathcal S(0)=Q$; apply (ii).
\end{proof}

\begin{lemma}[Mass of the zero spectrum; \stG]\label{zer:lem:mass}
For a primitive character $\chi$ of conductor $f$ and parity $b$,
\[
 M_\chi:=K_\chi(0)=\sum_\gamma\frac{m_\chi(\gamma)}{1/4+\gamma^2}=
 \begin{cases}2+\zergE-\log(4\pi),&f=1,\\[2pt]
 \log(f/\pi)+\psi\bigl((1+b)/2\bigr)+2\Re\dfrac{L'}{L}(1,\chi),&f>1,\end{cases}
\]
and $M_\chi>0$. For $f=1$ and $t\ge0$ one has exactly
\begin{equation}\label{zer:eq:zetaexact}
 K_\zeta(t)=e^{t/2}\Bigl[S(e^t)-t+1+\zergE-e^{-t}\log(2\pi)+\mathcal T(e^{-t})\Bigr],\qquad
 \mathcal T(z)=\sum_{k\ge1}\frac{z^{2k+1}}{2k(2k+1)},
\end{equation}
where $S(x)=\sum_{n\le x}\Lambda(n)(1/n-1/x)$ and
$\mathcal T(z)=z+\tfrac12(1-z)\log(1-z)-\tfrac12(1+z)\log(1+z)$.
\end{lemma}
\begin{proof}
By \cref{ana:lem:hadamardmass}, $\sum_\rho m_\chi(\rho)/(\rho(1-\rho))=2\Re\Xi_\chi'/\Xi_\chi(1)$ for the
completed entire function $\Xi_\chi$, and its gamma-factor evaluation gives the displayed values;
for $f=1$ one uses $\xi'/\xi(1)=1+\zergE-\tfrac12\log\pi+\tfrac12\psi(\tfrac12)$ and
$\psi(\tfrac12)=-\zergE-2\log2$. By \cref{src:thm:dgrh}, $\rho(1-\rho)=\tfrac14+\gamma^2$, so the sum is
$M_\chi$, a convergent sum of positive terms. For
\eqref{zer:eq:zetaexact}, the contour argument of \cref{zer:lem:ramp} for $-\zeta'/\zeta$, carried to
the left past all trivial zeros, gives residues $\log x-1-\zergE$ at $1$,
$x^{\rho-1}/(\rho(1-\rho))$ at each zero, $x^{-1}\log(2\pi)$ at $0$ and $-x^{-2k-1}/(2k(2k+1))$ at $-2k$;
multiply by $x^{1/2}$. At $t=0$, $S(1)=0$ and $\mathcal T(1)=1-\log2$.
\end{proof}

\paragraph{Peak rigidity.}
Put $E_a(x)=\sqrt x\bigl(S_a(x)-\log x/n_q-c_a\bigr)$, so $E_a(e^t)=Z_a(t)+\beta_a$, and
\[
 B_a(t)=\frac1{n_q}\sum_\chi\overline{\chi(a)}K_\chi(t),\qquad v_a=B_a(0),\qquad D(t)=v_1-B_1(t),
\]
where $B_1$ is the identity class. By \cref{zer:lem:ramp},
$E_a(x)=B_a(\log x)+O_q((1+\log x)/\sqrt x)$. Conjugation shows that $B_a$ and $v_a$ are real and
$v_a=v_{a^{-1}}$; $v_1=n_q^{-1}\sum_\chi M_{\chi^*}>0$, and every $v_a$ is a finite Fourier transform
of the masses of \cref{zer:lem:mass}.

\begin{theorem}[Peak rigidity; \stG]\label{zer:thm:peak}
The following hold for every $q\ge1$.
\begin{enumerate}[label=\textup{(\roman*)}]
\item $\limsup_{x\to\infty}E_1(x)=v_1$, and $\limsup_{x\to\infty}\abs{E_a(x)}\le v_1$ for every
$a$, with $\abs{v_a}<v_1$ for $a\ne1$.
\item $D(t)\ge0$, and for all $a$ and all real $t,u$,
\[
 \abs{B_a(t+u)-B_a(u)}^2\le2v_1D(t).
\]
\item For all $a$ and $t$,
\[
 \Bigl|\frac{B_a(t)+B_{a^{-1}}(t)}2-v_a\Bigr|\le D(t),\qquad
 \frac1{n_q}\sum_\chi\frac{\abs{K_\chi(t)-M_{\chi^*}}^2}{M_{\chi^*}}\le2D(t).
\]
The matrices $V_{a,b}=v_{ab^{-1}}$ and $H(t)_{a,b}=\tfrac12(B_{ab^{-1}}(t)+B_{ba^{-1}}(t))$ satisfy
$V\succ0$ and $V-H(t)\succeq0$, and $\tr(V-H(t))=n_qD(t)$.
\item For every $\eps>0$ the set of $t$ with $\sup_{a,u}\abs{B_a(t+u)-B_a(u)}<\eps$ is relatively
dense. Hence the near-peak scales of $E_1$ are relatively dense on the logarithmic axis.
\end{enumerate}
\end{theorem}
\begin{proof}
With $w_\gamma=(\tfrac14+\gamma^2)^{-1}$ and multiplicities retained, exactly
\[
 D(t)=\frac1{n_q}\sum_\chi\sum_\gamma w_\gamma(1-\cos\gamma t)=\frac1{2n_q}\sum_\chi\sum_\gamma w_\gamma\abs{e^{i\gamma t}-1}^2\ge0 .
\]
The translate difference is $n_q^{-1}\sum_{\chi,\gamma}w_\gamma\overline{\chi(a)}e^{i\gamma u}(e^{i\gamma t}-1)$, and
Cauchy--Schwarz for the measure of mass $w_\gamma/n_q$, of total mass $v_1$, gives (ii). Put
$D_\chi(t)=M_{\chi^*}-\Re K_\chi(t)\ge0$, whose average is $D(t)$. Then
$\abs{M_{\chi^*}-K_\chi(t)}^2\le M_{\chi^*}\sum_\gamma w_\gamma\abs{1-e^{i\gamma t}}^2=2M_{\chi^*}D_\chi(t)$, and
$v_a-\tfrac12(B_a(t)+B_{a^{-1}}(t))=n_q^{-1}\sum_\chi\Re\chi(a)D_\chi(t)$, which give (iii); the
vectors $(\overline{\chi(a)})_a$ diagonalize $V$ and $H(t)$ with eigenvalues $M_{\chi^*}$ and
$\Re K_\chi(t)$. \Cref{zer:lem:posspec}(i) makes $D(t)$ small on a relatively dense set, which with
(ii) gives (iv), and then $\limsup E_1=v_1$. Since all $M_{\chi^*}>0$ and some character separates
$a\ne1$ from $1$, the triangle inequality is strict in $\abs{v_a}<v_1$.
\end{proof}

\begin{corollary}[A unique exposed peak; \stG]\label{zer:cor:peak}
For $x_j\to\infty$ the following are equivalent: $E_1(x_j)\to v_1$;
$(E_a(x_j))_a\to(v_a)_a$; $E_a(x_je^u)\to B_a(u)$ for every $a$, uniformly for
$u\ge-\tfrac12\log x_j$. The hyperplane where the identity coordinate equals $v_1$ meets the
cluster set of $(E_a(x))_a$ in one point. The same holds along integers $x$.
\end{corollary}
\begin{proof}
If $E_1(x_j)\to v_1$ then $D(\log x_j)\to0$, and (ii) gives uniformity in $u$; the error in
$E_a=B_a+O(\cdot)$ tends to zero uniformly for $x_je^u\ge\sqrt{x_j}$. Rounding $e^t$ to an integer
changes $t$ by $o(1)$, and the $B_a$ are uniformly continuous.
\end{proof}

\begin{theorem}[Universal simultaneous peaks; \stG]\label{zer:thm:univpeak}
There are integers $n_j\to\infty$ such that for every $q\ge1$ and $a\in(\Z/q\Z)^\times$,
$E_{q,a}(n_je^u)-B_{q,a}(u)\to0$ uniformly for $u\ge-\tfrac12\log n_j$. In particular
$E_{q,1}(n_j)\to v_{q,1}$ for every $q$. Let $F:\Z\to\C$ have period $Q$, put
$\mu_F=\varphi(Q)^{-1}\sum_{a\in G_Q}F(a)$,
$c_F=\sum_{a\in G_Q}F(a)c_{Q,a}+\sum_{p\mid Q}\log p\sum_{k\ge1}F(p^k)p^{-k}$, and
$E_F(x)=\sqrt x\bigl(\sum_{n\le x}\Lambda(n)F(n)(1/n-1/x)-\mu_F\log x-c_F\bigr)$. Then
$E_F(x)=B_F(\log x)+O_{F}((1+\log x)/\sqrt x)$ with $B_F=\sum_{a\in G_Q}F(a)B_{Q,a}$, and
$E_F(n_je^u)\to B_F(u)$ uniformly for $u\ge-\tfrac12\log n_j$.
\end{theorem}
\begin{proof}
For moduli $1,\dots,j$ take, by \cref{zer:lem:posspec}(i) applied to the finitely many series
$B_{q,a}$, a return $t_j$ with uniform translate errors below $1/(2j)$, large enough that the
arithmetic errors are below $1/j$ for $x\ge\tfrac12e^{t_j/2}$ and that rounding $e^{t_j}$ to an
integer $n_j>n_{j-1}$ moves all these series by less than $1/(2j)$. For a periodic weight,
nonunits $n$ with $\Lambda(n)\ne0$ are powers of primes dividing $Q$, which give the additional
constant and an error $O((1+\log x)/x)$; the unit classes follow by linearity.
\end{proof}

This includes the weight $\mathbf 1_{n\equiv1\bmod q}(1-\chi_d(n))$ of the source proof. For a signed
combination the recurring value $B_F(0)$ need not be the maximum, but the recurrence of the
whole profile holds. Conversely, an eventual bound $E_1(x)\le C$ already excludes zeros of the
$L(s,\chi)$, $\chi\bmod q$, from $\Re s>\tfrac12$, by \cref{ana:thm:zerotest} with all weights
$1/n_q$. Thus GRH for the characters modulo $q$ holds if and only if $\limsup E_1(x)<\infty$, and
then the limsup is $v_1$. For $q=1$ this is the equivalence of \cref{grp:prop:split} with $K=\Q$,
since the prime powers $p^k$, $k\ge2$, change $\sqrt x$ times the remainder by $O(1)$ only. For
example $v_1=M_\zeta=2+\zergE-\log(4\pi)\approx0.0461914179$ for $q=1$, by \cref{zer:lem:mass}; for $q=3$,
$v_1\approx0.0797106939$ and $v_2\approx-0.0335192760$; and for $q=5$ the classes $1$, $2$,
$3$, $4$ give approximately $0.1522978093$, $-0.0275913840$, $-0.0275913840$ and
$-0.0509236233$. These decimals are
evaluated in floating point from \cref{zer:lem:mass} in \artifact{prime_peak_rigidity/check.py}.

\paragraph{Negative excursions of the source remainder.}
The remainder of the source is bounded from above by the source itself. Its excursions below
zero at modulus one can be located as follows. Fix a prime $d\equiv7\pmod8$, let
$\chi=\chi_d$, and write $R_d=R_{d,1}$, so that
\[
 R_d(x)=\sum_{n\le x}\Lambda(n)\bigl(1-\chi(n)\bigr)\Bigl(\frac1n-\frac1x\Bigr)
 -\log x+1+\zergE-\frac{L'(1,\chi)}{L(1,\chi)} .
\]
Put
\[
 B_0=2+\zergE-\log(4\pi)=0.0461914179\ldots,\qquad
 B_d=\log\frac d\pi-\zergE+2\,\frac{L'(1,\chi)}{L(1,\chi)} ;
\]
since $\chi$ is odd, these are the masses $M_\zeta$ and $M_\chi$ of \cref{zer:lem:mass}. The weight
$F=1-\chi$ of \cref{zer:thm:univpeak} has period $d$ and $\mu_F=1$, and $E_F(x)=\sqrt x\,R_d(x)$, both
centring constants being the unique ones for which the remainder tends to $0$; orthogonality gives
$B_F=K_\zeta-K_\chi$. Hence one sequence of integers $n_j$, the same for all such $d$, satisfies
$\sqrt{n_j}\,R_d(n_j)\to B_F(0)=B_0-B_d$ for every $d$ simultaneously; in particular
$\liminf_{x\to\infty}\sqrt x\,R_d(x)\le B_0-B_d$. The following theorem adds the lower bound,
geometric sequences of scales and the growth in $d$.

\begin{theorem}[Excursion bounds; \stG]\label{ana:thm:excursion}
For every prime $d\equiv7\pmod8$,
\[
 -B_d-B_0\le\liminf_{x\to\infty}\sqrt x\,R_d(x)\le-B_d+B_0 .
\]
There is a single sequence of integer exponents $n_j\to\infty$ with
$\sqrt{a^{n_j}}\,R_d(a^{n_j})\to B_0-B_d$ for every integer $a\ge2$ and every $d$. In
particular the lower limit divided by $\log d$ tends to $-1$ as $d\to\infty$.
\end{theorem}

The limits are taken in this order, first in $x$ for fixed $d$. Nothing is asserted about the
first $x$ at which a given excursion occurs.

\begin{proof}
Here $K_\zeta$ and $K_\chi$ are the zero series of \cref{zer:sec:cov}, over all nontrivial zeros
with multiplicity, central zeros included; they are real because $\chi$ is real. For fixed $d$,
\begin{equation}\label{ana:eq:O1}
 \sqrt x\,R_d(x)=K_\zeta(\log x)-K_\chi(\log x)+O_d\bigl(x^{-1/2}\bigr).
\end{equation}
This is the expansion of \cref{zer:thm:univpeak} for $F=1-\chi$ with a sharper remainder: in the
contour argument of \cref{zer:lem:ramp}, applied to $-\zeta'/\zeta+L'/L(\cdot,\chi)$, the kernel has
only a simple pole at $0$, since $\zeta(0)L(0,\chi)\ne0$ for this odd primitive character, and its
residue depends only on $d$. As $\abs{K(u)}\le K(0)$ for both series (\cref{zer:lem:posspec}(ii)),
\eqref{ana:eq:O1} proves the lower bound; the upper bound was noted before the theorem.

For the integer exponents, given $j$ take the finitely many frequencies $\gamma\log a$ with
$2\le a\le j$ and $\gamma$ among the first $j$ ordinates of $\zeta$ and of each $L(\cdot,\chi_d)$ with
$d\le j$. The integer iterates of their phase vector in the finite torus have a convergent
subsequence, so differences of increasingly separated indices return all phases to within $1/j$
of $1$; choose such an exponent $n_j>n_{j-1}$. Absolute convergence of the series gives
$K_\zeta(n_j\log a)-K_\chi(n_j\log a)\to B_0-B_d$ for every $a$ and $d$, and \eqref{ana:eq:O1} at
$x=a^{n_j}$ gives the claim. Finally $L'(1,\chi)/L(1,\chi)=O(\log\log d)$ under \cref{src:thm:dgrh}
\cite[Corollary~3.3.2]{IharaMurtyShimura2009Logarithmic}, which gives $B_d\sim\log d$.
\end{proof}

\begin{proposition}[When the lower endpoint is attained; \stG]\label{ana:prop:parity}
The value $-B_d-B_0$ is a limiting value of $\sqrt x\,R_d(x)$ if and only if every finite
integer relation among the positive ordinates of $\zeta$ and of $L(\cdot,\chi)$ has an even
total coefficient on the ordinates of $\zeta$. Here each occurrence of an ordinate is
labelled by its function, so a positive ordinate shared by the two functions already violates
the condition.
\end{proposition}

\begin{proof}
Equality forces each cosine with positive weight in $K_\zeta$ to approach $-1$ and each
cosine in $K_\chi$ to approach $+1$, because their individually nonnegative
deficits sum to zero in the limit. The function $\zeta$ has no central zero, and central
zeros of $L$ contribute with the required sign. On each finite set of frequencies,
Kronecker's theorem identifies the closure of the real flow with the subtorus cut out by the
integer relations, and the required sign vector lies in that closure exactly under the
stated parity condition. A diagonal sequence and the summable tails give the sufficiency for
all ordinates.
\end{proof}

For $d=7$ the two endpoints are approximately $-0.30142740$ and $-0.20904456$, and for
$d=23$ they are $-0.61549819$ and $-0.52311535$. Further values are in the file
\begin{center}
\artifact{grh_method_transfers_sharp_torsion/checks/oscillation_bounds.txt}.
\end{center}
These numbers evaluate the proven bounds through
\[
 \frac{L'(1,\chi)}{L(1,\chi)}=\log(2\pi)+\zergE-\frac1{L(0,\chi)}\sum_{a=1}^{d-1}\chi(a)\log\Gamma(a/d),
 \qquad L(0,\chi)=-\frac1d\sum_{a=1}^{d-1}a\chi(a).
\]
They are not estimates of a lower limit by sampling.

\paragraph{How often the peak is approached.}
\Cref{zer:thm:peak}(iv) makes the near-peak scales of $E_1$ relatively dense. Their logarithmic
frequency can be bounded below with \cref{src:thm:dgrh} alone, and it can be computed through its
first conductor-dependent term under an additional hypothesis on the rational rank of the
ordinates. Let $P_q^-(\eps)$ and $P_q^+(\eps)$ be the lower and upper limits, as $Y\to\infty$, of
\[
 \frac1Y\int_1^{e^Y}\mathbf 1_{\set{E_1(x)>v_1-\eps}}\,\frac{\dd x}x ,
\]
where $q$ is fixed and limits in $\eps$ are taken after these limits in $Y$. Put
\[
 g(u)=u-\log I_0(u),\qquad A=\int_0^\infty g(v^{-2})\,\dd v,\qquad
 B=\int_0^\infty g(v^{-2})\log v\,\dd v,\qquad C=\frac A{2\pi},
\]
where $I_0$ is the modified Bessel function. Numerically
\[
 C=0.4427822838\ldots,\qquad B/A=-0.6933500542\ldots .
\]
Pool, over the characters modulo $q$, the positive ordinates of the
primitive inducing $L$-functions, counting multiplicities and keeping repetitions; let $N_q(T)$
be the length of this list through height $T$ and $r_q(T)$ the dimension of the $\Q$-span of these
ordinates. Let $\mathrm{rd}_q=q\prod_{p\mid q}p^{-1/(p-1)}$.

\begin{theorem}[Frequency of endpoint visits]\label{zer:thm:visits}
Write $W=\log(1/\eps)$.
\begin{enumerate}[label=\textup{(\roman*)}]
\item \textup{(\stG)} For every $q$,
\[
 \limsup_{\eps\downarrow0}\frac\eps{W^2}\log\frac1{P_q^-(\eps)}\le n_qC^2 .
\]
The same bound holds for the lower fraction of indices $j$ with $E_1(e^{jh})>v_1-\eps$, along every
fixed mesh $h>0$.
\item \textup{(\stG, with the rank estimate as a hypothesis)} Suppose that
\begin{equation}\label{zer:eq:rankloss}
 N_q(T)-r_q(T)=o_q(T/\log T).
\end{equation}
Then, for $P=P_q^-$ and for $P=P_q^+$,
\begin{equation}\label{zer:eq:visitlaw}
 \log\frac1{P(\eps)}=\frac{n_qC^2}{2\eps}\Bigl[W^2+2W\log W
 +2W\Bigl(\log C+\log\frac{\mathrm{rd}_q}{2\pi}+\frac BA\Bigr)+o_q(W)\Bigr].
\end{equation}
If the pooled positive ordinates are rationally independent, $P_q^-(\eps)=P_q^+(\eps)$ for every
$\eps>0$.
\end{enumerate}
\end{theorem}

Condition \eqref{zer:eq:rankloss} allows infinitely many repetitions and rational relations,
including relations between low and high ordinates, provided the cumulative loss of dimension obeys
the bound; for example it suffices that rational independence holds after discarding
$O_q(T/\log^2T)$ ordinates through height $T$. \Cref{src:thm:dgrh} gives no rank estimate of this
strength. The constant in (i) is twice the leading constant of \eqref{zer:eq:visitlaw}.

\begin{corollary}[The modulus from the visit frequency]\label{zer:cor:visitmodulus}
Among the moduli satisfying \eqref{zer:eq:rankloss}, either density germ $P_q^\pm$ determines the
canonical cyclotomic modulus, which is $q/2$ if $q\equiv2\pmod4$ and $q$ otherwise. For instance, if
$q=5$ and $q=8$ both satisfy \eqref{zer:eq:rankloss}, then
\[
 \lim_{\eps\downarrow0}\frac\eps W\log\frac{P_5(\eps)}{P_8(\eps)}=4C^2\log\frac4{5^{3/4}}>0 .
\]
\end{corollary}
\begin{proof}
The limit of $2\eps\log(1/P)/(C^2W^2)$ is $n_q$. Subtracting the first two terms of the bracket in
\eqref{zer:eq:visitlaw} and dividing by $2W$ gives $\log C+\log(\mathrm{rd}_q/(2\pi))+B/A$, hence
$\mathrm{rd}_q$. In the product for $\mathrm{rd}_q$ the exponent of each odd $p\mid q$ is
$v_p(q)-1/(p-1)>0$, and that of $2$ is $v_2(q)-1$ when $q$ is even; unique factorization after
clearing the rational exponents gives every odd prime power and every $2$-part at least four. For
odd $q$ the characters modulo $q$ and $2q$ have the same primitive inducing characters, and their
normalized ramps differ by $o(1)$. For the example, $n_5=n_8=4$ and the root discriminants are
$5^{3/4}$ and $4$.
\end{proof}

By \cref{zer:lem:ramp}, $E_1(e^t)=B_1(t)+O_q((1+t)e^{-t/2})$ and
\[
 D(t)=v_1-B_1(t)=\sum_{\gamma>0}a_\gamma(1-\cos\gamma t),\qquad a_\gamma=\frac2{n_q(1/4+\gamma^2)},
\]
the sum running over the pooled list; conjugation pairs the two signs over the full character
family, and central zeros cancel.

\begin{lemma}[Comparison of tilted means; \stU]\label{zer:lem:tilt}
Let $a_j>0$ with $\sum_ja_j<\infty$, let $\omega_j$ be arbitrary real numbers, and put
$D(t)=\sum_ja_j(1-\cos\omega_jt)$ and $J(s)=\sum_jg(sa_j)$ for $s\ge0$. Let $\nu$ be a probability
measure on $\R$ whose Fourier transform is real and nonnegative, and
$L_\nu(s)=\int e^{-sD(t)}\,\dd\nu(t)$. Then $e^{J(s)}L_\nu(s)$ is nondecreasing in $s$ and equals
one at $s=0$. Consequently, whenever $J'(s)<\eps$,
\begin{equation}\label{zer:eq:tilt}
 \nu\set{D<\eps}\ge e^{-J(s)}\bigl(1-J'(s)/\eps\bigr).
\end{equation}
\end{lemma}
\begin{proof}
For finitely many indices, expanding $e^{u\cos\theta}$ in its Fourier series and dividing by the
constant coefficient gives
\[
 e^{J(s)}L_\nu(s)=\sum_{k\in\Z^m}\widehat\nu\Bigl(\sum_jk_j\omega_j\Bigr)\prod_j\frac{I_{\abs{k_j}}(sa_j)}{I_0(sa_j)} .
\]
Every factor is nonnegative and nondecreasing: $I_b(u)/u^b$ is a positive constant times
$\int_{-1}^1e^{uy}(1-y^2)^{b-1/2}\,\dd y$, its logarithmic derivative is $I_{b+1}(u)/I_b(u)$, the
derivative of this ratio is the variance of $y$ under the tilted integral, and the ratios $I_k/I_0$
telescope into products of these \cite{NISTDLMFBessel}. The infinite case follows by uniform
convergence of $D$ and $0\le g(u)\le u$. Differentiating gives
$\int De^{-sD}\,\dd\nu/L_\nu(s)\le J'(s)$; Markov's inequality in the tilted probability gives mass at
least $1-J'(s)/\eps$ to $\set{D<\eps}$, and the tilt is undone using $e^{sD}\ge1$ there and
$L_\nu\ge e^{-J}$.
\end{proof}

The triangular density $T^{-1}(1-\abs t/T)_+$, the law of the difference of two independent
uniform points of $[0,T]$, has nonnegative Fourier transform, and its mass on an even event is at
most twice the ordinary time fraction of the event in $[0,T]$. Half the right side of
\eqref{zer:eq:tilt} therefore bounds that fraction at every horizon $T$. The difference of two
independent uniform indices of $\set{0,\dots,N-1}$, multiplied by $h>0$, gives the same comparison
on every mesh, with no exclusion for aliased frequencies.

\begin{lemma}[Two-term transform; \stG]\label{zer:lem:visittransform}
Let $f_\chi$ be the conductor of $\chi^*$, $b_q=n_q^{-1}\sum_\chi\log(f_\chi/(2\pi))$, $z=\sqrt{2s/n_q}$ and
$d=b_q+B/A$, and put $J(s)=\sum_{\gamma>0}g(z^2/(\gamma^2+1/4))$ over the pooled list and
$F(s)=n_qCz(\log z+d)$. As $z\to\infty$,
\[
 J(s)=F(s)+O_q(\log^2z),\qquad J'(s)=\frac Cz(\log z+d+1)+O_q(\log^2z/z^2).
\]
Moreover $b_q=\log(\mathrm{rd}_q/(2\pi))$.
\end{lemma}
\begin{proof}
The zero count for both signs, summed over the family and halved by conjugation, gives
\[
 N_q(T)=\frac{n_qT}{2\pi}(\log T+b_q-1)+O_q(\log(2+T))\qquad(T\ge1)
\]
\cite{BennettMartinOBryantRechnitzer2021Counting}; central zeros affect only the remainder. For
$p^a\parallel q$ the number of local characters of conductor dividing $p^j$ is $\varphi(p^j)$, so the
average exponent of the local conductor is $a-p^{a-1}/\varphi(p^a)=a-1/(p-1)$; local components vary
independently, and the geometric mean of the $f_\chi$ is $\mathrm{rd}_q$.

Let $Y=1-\cos\Theta$ with $\Theta$ uniform on $[-\pi,\pi]$, so $g(u)=-\log\mathbb Ee^{-uY}$,
$g'(u)=\mathbb E_uY$ and $-g''(u)=\operatorname{Var}_uY$ for the tilt by $e^{-uY}$. Since
$Y\asymp\Theta^2$, comparison with Gaussian integrals gives $\mathbb E_uY^k\ll_ku^{-k}$ for $k=1,2$
and $u\ge1$, and both moments are bounded for small $u$. For $H(u)=ug'(u)$ and $h=g,H$ it follows that
$\abs{uh'(u)}\ll\min(1,u)$, $g(u)=O(\log(2+u))$ for $u\ge1$, and $g(u),H(u)=O(u)$ for $0<u\le1$; these
also give convergence of $A$ and $B$. Replacing $\gamma^2+\tfrac14$ by $\gamma^2$ in a sum of $h$
costs $O_q(1)$. Stieltjes integration against the zero count gives
\[
 \sum_{\gamma>0}h(z^2/\gamma^2)=\frac{n_qz}{2\pi}\int_0^\infty h(v^{-2})(\log z+b_q+\log v)\,\dd v+O_q(\log^2z),
\]
because $\abs{\partial_th(z^2/t^2)}\ll t^{-1}\min(1,z^2/t^2)$: above one, integration against the
error $O_q(\log(2+t))$ costs $O_q(\log^2z)$, and below the first ordinate the relative counting
error is $O_q(t(1+\abs{\log t}))$. For $h=g$ this is the first estimate. For $h=H$ integrate by parts
using $H(v^{-2})=-(v/2)\partial_vg(v^{-2})$, whose integrals against $1$ and $\log v$ are $A/2$ and
$(A+B)/2$; since $sJ'(s)=\sum H(sa_\gamma)$, this gives the second.
\end{proof}

\begin{proof}[Proof of \cref{zer:thm:visits}\textup{(i)}]
Fix $\lambda>1$ and take $z=\lambda CW/\eps$. Then $J'(s)/\eps\to1/\lambda$ and
$J(s)\sim n_q\lambda C^2W^2/\eps$. Apply \eqref{zer:eq:tilt}, halve it for ordinary time or mesh
averages, and let $\lambda\downarrow1$. This proves the bound for the event $D<\eps$; for the
arithmetic event use $D<\eps(1-1/W^2)$, since $E_1(e^t)$ differs from $B_1(t)$ by $o_q(1)$ at late
times, and this preserves the constant.
\end{proof}

For the sharp law, let $\mathcal T$ be the closure of $t\mapsto(e^{i\gamma t})_\gamma$ in the product
of circles, with Haar probability $\mu$. The deficit extends continuously to $\mathcal T$, long-time
means of continuous functions are Haar integrals (\cref{zer:lem:posspec}(i) and Stone--Weierstrass),
and hence an open set has lower time density at least its Haar measure and a closed set has upper
time density at most its Haar measure; no atom-free limiting law is assumed.

\begin{lemma}[Deleting redundant ordinates; \stG]\label{zer:lem:visitsparse}
Under \eqref{zer:eq:rankloss} the Haar transform $L(s)=\int_{\mathcal T}e^{-sD}\,\dd\mu$ satisfies
$-\log L(s)=F(s)+o_q(z/\log z)$.
\end{lemma}
\begin{proof}
Traverse the list in nondecreasing height, keep a frequency each time it increases the rational
span, and put all other occurrences in a list $\mathcal E$, whose counting function is $N_q-r_q$. The
kept angles have independent uniform finite-dimensional Haar marginals, by the character criterion.
Dropping the nonnegative deficit carried by $\mathcal E$ gives
\[
 e^{-J(s)}\le L(s)\le\exp\bigl[-J(s)+J_{\mathcal E}(s)\bigr],\qquad
 J_{\mathcal E}(s)=\sum_{\gamma\in\mathcal E}g(sa_\gamma),
\]
the lower bound by the Haar mean of the expansion in the proof of \cref{zer:lem:tilt}, the upper
one without any independence between $\mathcal E$ and its complement. With
$\mathcal E(t)=\#\set{\gamma\in\mathcal E:\gamma\le t}$, integration by parts bounds $J_{\mathcal E}$
by $O_q(1+\int_0^\infty t^{-1}\mathcal E(t)\min(1,z^2/t^2)\,\dd t)$. Up to $\sqrt z$ the integral is
$O_q(\sqrt z\log(2z))$; for $t\ge\sqrt z$, $\mathcal E(t)/t=o_q(1/\log z)$ uniformly by
\eqref{zer:eq:rankloss}, and $\int_0^\infty\min(1,z^2/t^2)\,\dd t=2z$. Hence $J_{\mathcal E}=o_q(z/\log z)$,
and \cref{zer:lem:visittransform} concludes.
\end{proof}

\begin{proof}[Proof of \cref{zer:thm:visits}\textup{(ii)}]
Let $z_0$ be the large solution of $\eps z_0=C(\log z_0+d+1)$ and $s_0=n_qz_0^2/2$; it is unique on
the decreasing branch of $(\log z+d+1)/z$, $z_0\asymp W/\eps$, and $F'(s_0)=\eps$. With the slack
$\eta=\eps/W^2$, an arithmetic near-peak visit at a late time implies $D\le\eps+\eta$, so the upper
time density is at most the Haar mass of this closed set. Markov's inequality for $e^{-s_0D}$ and
\cref{zer:lem:visitsparse} give $\log P_q^+(\eps)\le s_0\eps-F(s_0)+o_q(z_0)$, since
$s_0\eta=O_q(z_0/W)$.

For the lower bound take $h=h(\eps)\to0$ so slowly that $h\ge W^{-1/2}$ and
$\sup\abs{\log L(s)+F(s)}\le h^4z_0/W$ over $z_0/2\le\sqrt{2s/n_q}\le2z_0$, which is possible by
\cref{zer:lem:visitsparse} and $\log z_0\sim W$. Reset $\eta=\eps h/W$ and choose $s$ on the same
branch with $F'(s)=u=\eps-2\eta$; then $s/s_0$ and $z/z_0$ are $1+O_q(h/W)$. In the tilted measure
$\dd\mu_s=e^{-sD}\dd\mu/L(s)$ one has $\mu_s\set{\abs{D-u}\ge\eta}=o_q(1)$. To see this without
differentiating the remainder, put $v=sh^{3/2}/W$; the logarithmic moment generating function of
$\pm v(D-u)$ is $\mp vu+\log L(s\mp v)-\log L(s)=O_q(v^2W/z^3+h^4z_0/W)$, by Taylor expansion of $F$
and $\abs{F''}\ll_qW/z^3$ on this range, the linear terms cancelling. Subtracting
$v\eta\asymp_qz_0h^{5/2}/W$ in the exponential Markov bound, the quadratic term is
$O_q(z_0h^3/W)$, so the total error is $o_q(v\eta)$ while $v\eta\to\infty$.

The open event $\abs{D-u}<\eta$ forces an arithmetic visit at all late times in it. Its lower time
density is at least its Haar mass, which after undoing the tilt is at least
$L(s)e^{s(u-\eta)}(1-o_q(1))$, and $s(u-\eta)-F(s)=s_0\eps-F(s_0)+o_q(z_0)$: the tolerance costs
$O_q(s_0\eta)=O_q(z_0h)$ and moving the stationary point costs $O_q((s-s_0)^2W/z_0^3)=o_q(z_0)$. Thus
$\log P_q^\pm(\eps)=s_0\eps-F(s_0)+o_q(z_0)$. With $l=\eps z_0/C$ the defining equation reads
$l=W+\log l+\log C+d+1$, so $l=W+\log W+\log C+d+1+o_q(1)$, and
$F(s_0)-s_0\eps=n_qC^2(l^2-2l)/(2\eps)$. Expanding and inserting $b_q=\log(\mathrm{rd}_q/(2\pi))$ gives
\eqref{zer:eq:visitlaw}. Under independence of the whole list the Haar law of $D$ is that of a sum of
independent variables, one of them continuous, so it has no atoms, and continuous approximation
gives the existence of the density at each threshold.
\end{proof}

The finite-horizon comparison of \cref{zer:lem:tilt} controls centred spectral averages from below
at every horizon. Floating-point checks of the constants $A,B$, of the tilt inequalities on
finite lists with repetitions and exact resonances, and of the zeta ramp against $32$ ordinates are
in the following script, with its output beside it:
\begin{center}\footnotesize
\artifact{prime_endpoints_and_symplectic_compressions/endpoint/check_endpoint.py}
\end{center}
For an
independent series, the leading small-deviation behaviour of regularly varying weights belongs to
the theory of Borovkov and Ruzankin \cite{BorovkovRuzankin2008Small}.

\paragraph{A balanced window with a sharp normalized law.}
Let $b_\chi\ge0$ $(\chi\bmod q)$ be real, not all zero, with $b_{\overline\chi}=b_\chi$, and put
$w(n)=\sum_\chi b_\chi\chi(n)$, $\mu=b_{\chi_0}$, $S_w(x)=\sum_{n\le x}\Lambda(n)w(n)(1/n-1/x)$. For
$c>1$ let $\alpha=c/(c+1)$, $\alpha'=1/(c+1)$ and
\[
 \mathcal W_w(x;c)=S_w(x)-\alpha S_w(cx)-\alpha'S_w(x/c)+\mu(\alpha-\alpha')\log c\qquad(x\ge c),
\]
which equals $E_w(x)-\alpha E_w(cx)-\alpha'E_w(x/c)$ with $E_w=S_w-\mu\log x$; no constant of the
prime sum is needed. Its Mellin transform carries the factor
$P_c(s)=1-\alpha c^{s-1}-\alpha'c^{1-s}=(1-u)(cu-1)/((c+1)u)$, $u=c^{s-1}$, which vanishes only on
$\Re s\in\set{0,1}$ and on the critical line equals $1-\kappa_c\cos(\gamma\log c)>0$, where
$\kappa_c=2\sqrt c/(c+1)<1$.

\begin{theorem}[Sharp balanced-window law; \stG]\label{zer:thm:balanced}
Put
\[
 L_w(c)=\sum_\chi b_\chi\sum_\gamma m_\chi(\gamma)\frac{1-\kappa_c\cos(\gamma\log c)}{1/4+\gamma^2}>0 .
\]
For every strictly increasing $x_j\to\infty$ with $x_{j+1}/x_j\to1$,
\[
 \limsup_{j\to\infty}\sqrt{x_j}\,\mathcal W_w(x_j;c)=\limsup_{j\to\infty}\sqrt{x_j}\,\abs{\mathcal W_w(x_j;c)}=L_w(c),
\]
and the same value is the limsup over all real $x$.
\end{theorem}
\begin{proof}
By \cref{zer:lem:ramp}, $E_w(x)=\sum_\chi b_\chi c_\chi+x^{-1/2}f_w(\log x)+O((1+\log x)/x)$ with
$f_w=\sum_\chi b_\chi K_\chi$. Since $\alpha+\alpha'=1$ and $\alpha c^{-1/2}=\alpha'c^{1/2}=\kappa_c/2$,
\[
 \sqrt x\,\mathcal W_w(x;c)=f_w(t)-\tfrac{\kappa_c}2\bigl(f_w(t+\log c)+f_w(t-\log c)\bigr)
 +O\bigl((1+t)e^{-t/2}\bigr),\qquad t=\log x,
\]
a series with coefficients $b_\chi m_\chi(\gamma)(1-\kappa_c\cos(\gamma\log c))/(\tfrac14+\gamma^2)\ge0$ and value
$L_w(c)$ at $t=0$. Apply \cref{zer:lem:posspec}(ii); $\log x_j$ has gaps tending to zero.
\end{proof}

The positive weights are essential, and the statement uses neither simplicity nor
independence of ordinates. On a mesh with power-saving gaps, the sparse one-sided criterion of
\cref{sec:ana} shows conversely that a finite limsup forces all active zeros onto the critical
line, and that without GRH the limsup is $+\infty$. For $\zeta$ ($q=1$, $w=1$), \eqref{zer:eq:zetaexact}
evaluates the constant from one prime sum:
\[
 L_1(c)=M_\zeta-\kappa_cK_\zeta(\log c)
 =2+\zergE-\log(4\pi)-\frac{2c}{c+1}\Bigl(S(c)-\log c+1+\zergE-\frac{\log(2\pi)}c+\mathcal T(1/c)\Bigr),
\]
and $S(2)=0$ gives $L_1(2)=0.0625323180445\ldots$ (a numerical evaluation of this closed form).
Consequently, for every fixed integer $k\ge1$,
\[
 \limsup_{n\to\infty}n^{k/2}\Bigl(S(n^k)-\tfrac23S(2n^k)-\tfrac13S(n^k/2)+\tfrac13\log2\Bigr)=L_1(2).
\]
Since $P_c(0)=0$, the term $\log(2\pi)/x$ of \eqref{zer:eq:zetaexact} cancels exactly in
$\mathcal W_1$, so for all $x\ge c$ the corrected quantity
$\sqrt x\bigl(\mathcal W_1(x;c)+\mathcal T(1/x)-\alpha\mathcal T(1/(cx))-\alpha'\mathcal T(c/x)\bigr)$ is exactly the
positive-coefficient series at $\log x$ and has absolute value at most $L_1(c)$; the constant is
optimal on every mesh of the theorem.

\paragraph{Sharp inequalities between compact prime windows.}
Fix an odd prime $q$, $d=q-1$, $G=(\Z/q\Z)^\times$ and $0<h<\log2$. Define
\[
 g_h(u)=\begin{cases}e^{h-\abs u/2}-e^{\abs u/2},&\abs u\le h,\\0,&\abs u>h,\end{cases}\qquad
 \mathsf H=e^h-1,\qquad\mathsf B=h\mathsf H,
\]
and $b_j(u)=e^{-(j+1/2)u}/(1-e^{-2u})$ for $u>0$ and $j=0,1$.
\[
\]
For $a\in G$ and $t>h$ put
\begin{align*}
 P(a,t)&=\sum_{n\equiv a\bmod q}\frac{\Lambda(n)}{\sqrt n}g_h(\log n-t)+\frac1d\sum_{k\ge1}\frac{\log q}{q^{k/2}}g_h(k\log q-t),\\
 A(a,t)&=\frac12\int_{t-h}^{t+h}g_h(u-t)\bigl[\mathbf 1_{a=1}(b_0+b_1)(u)+\mathbf 1_{a=-1}(b_0-b_1)(u)\bigr]\dd u,\\
 C(t)&=\frac{\mathsf B}d\bigl(e^{t/2}+e^{-t/2}\bigr),\qquad \mathsf E(a,t)=C(t)-P(a,t)-A(a,t).
\end{align*}
The prime sum samples $e^{t-h}\le n\le e^{t+h}$; the second term comes from $\zeta$, the
primitive character inducing the principal one. With
$J_j=2\int_0^\infty(\mathsf H-g_h(u))b_j(u)\,\dd u$,
$\mathcal D(f,j,\delta)=\mathsf H(\log(f/\pi)+\psi(\tfrac14+\tfrac j2))+J_j+2\delta\mathsf B$ and
\[
 \mathsf M=\frac1d\Bigl(\mathcal D(1,0,1)+\frac{q-3}2\mathcal D(q,0,0)+\frac{q-1}2\mathcal D(q,1,0)\Bigr),
\]
no prime or zero enters $\mathsf M$, because $h<\log2$.

\begin{theorem}[Sharp multiscale contrast; \stG]\label{zer:thm:sharpwin}
$\mathsf E$ extends to a continuous real positive-definite function $\mathcal K$ on $G\times\R$,
\[
 \mathcal K(a,t)=\frac1d\sum_{\chi\bmod q}\overline{\chi(a)}\sum_\gamma m_\chi(\gamma)w_h(\gamma)e^{i\gamma t},\qquad
 w_h(y)=\frac{e^h+1-2e^{h/2}\cos(hy)}{1/4+y^2}>0,
\]
with $\mathcal K(1,0)=\mathsf M>0$ and $\mathcal K(a^{-1},-t)=\mathcal K(a,t)$. Hence
\[
 \abs{\mathsf E(a,t)}\le\mathsf M\quad(t>h),\qquad\limsup_{t\to\infty}\mathsf E(1,t)=\mathsf M .
\]
For a finite family $(a_i,t_i)$ and $c_i\in\C$ let
$Q=\sum_{i,j}c_i\overline{c_j}\mathcal K(a_ia_j^{-1},t_i-t_j)\ge0$. For
$T>h+\max\abs{t_i-t_j}$,
\[
 \Bigl|\sum_{i,j}c_i\overline{c_j}\,\mathsf E(a_ia_j^{-1},T+t_i-t_j)\Bigr|\le Q,
\]
and the limsup of the left side as $T\to\infty$ is $Q$. If all differences $\abs{t_i-t_j}$, $i\ne j$,
exceed $h$, then $Q$ is computed from the displayed prime sums, integrals and $\mathsf M$.
\end{theorem}
\begin{proof}
With $k(v)=e^{v/2}\mathbf 1_{[0,h]}(v)$ one has $g_h(u)=\int k(v)k(v+u)\,\dd v$, so
$\widehat g_h(y)=\int g_h(u)e^{iyu}\dd u=\abs{(e^{(1/2+iy)h}-1)/(1/2+iy)}^2=w_h(y)$, and
$\int g_h(u)e^{\pm u/2}\dd u=\mathsf B$. For a primitive character of conductor $f$, parity $j$ and
$\delta=1$ for conductor one, and for a real compactly supported Lipschitz $F$ that is piecewise
smooth, the Weil explicit formula reads
\begin{align*}
 \sum_\gamma m(\gamma)\widehat F(\gamma)={}&\delta\int F(u)(e^{u/2}+e^{-u/2})\dd u
 +F(0)\bigl(\log(f/\pi)+\psi(\tfrac14+\tfrac j2)\bigr)\\
 &+\int_0^\infty(2F(0)-F(u)-F(-u))b_j(u)\dd u\\
 &-\sum_{n\ge2}\frac{\Lambda(n)}{\sqrt n}\bigl(\chi(n)F(\log n)+\overline{\chi(n)}F(-\log n)\bigr),
\end{align*}
obtained from the residue theorem for the completed logarithmic derivative, integrated
against the transform $\int F(u)e^{(s-1/2)u}\dd u$, with the digamma integral representation for the Gamma
terms \cite{Koukoulopoulos2019Distribution}; smooth compact mollification and the bound
$\widehat F(y)=O((1+y^2)^{-1})$ justify it for these $F$. Take $F=g_h(\cdot-t)$ with $t>h$, so
$F(0)=F(-\log n)=0$, and average with $\overline{\chi(a)}/d$: the pole terms give $C(t)$;
orthogonality gives the class $a$ for $(n,q)=1$, and $n=q^k$ only through $\zeta$, with weight
$1/d$, which is $P(a,t)$; and writing
\[
 b_j=\tfrac12(b_0+b_1)+\tfrac12\chi(-1)(b_0-b_1)
\]
gives $A(a,t)$.
So $\mathcal K=\mathsf E$ for $t>h$. At $(1,0)$ the prime terms vanish since $h<\log2$, each character
gives $\mathcal D(f,j,\delta)$, and the $d/2$ even characters include the principal one. The
remaining assertions are \cref{zer:lem:posspec}(iii), the weights $m_\chi(\gamma)w_h(\gamma)/d$ being
positive.
\end{proof}

\begin{corollary}[Two and three windows; \stG]\label{zer:cor:sharpwin}
For $u>h$, $a\in G$ and $T>u+h$,
\[
 \abs{2\mathsf E(1,T)-\mathsf E(a,T+u)-\mathsf E(a^{-1},T-u)}\le2\bigl(\mathsf M-\mathsf E(a,u)\bigr),
\]
sharp in the limsup, which improves the triangle bound $4\mathsf M$ whenever $\mathsf E(a,u)>-\mathsf M$.
For $g_1,g_2\in G\times\R$ with $X=\mathcal K(g_1)$, $Y=\mathcal K(g_2)$, $Z=\mathcal K(g_1g_2^{-1})$,
\[
 \mathsf M^3+2XYZ-\mathsf M(X^2+Y^2+Z^2)\ge0,
\]
equivalently, if $\abs Z<\mathsf M$, $\dfrac{(X+Y)^2}{2\mathsf M(\mathsf M+Z)}+\dfrac{(X-Y)^2}{2\mathsf M(\mathsf M-Z)}\le1$.
\end{corollary}
\begin{proof}
Take $(1,0),(a,u)$ with coefficients $(1,-1)$, so $Q=2(\mathsf M-\mathcal K(a,u))$; and use the
nonnegative determinant of the Gram matrix of $(1,0),g_1,g_2$.
\end{proof}

For $q=5$, $h=0.4$, $a=2$, $u=\log10$ and $T=\log1000$, the prime-sum evaluation gives
$\mathsf M\approx0.44998$, a translated contrast $\approx-0.21147$ and the sharp bound $\approx0.63564$,
about $35\%$ of the triangle bound $4\mathsf M$; these are floating-point values from
\begin{center}\artifact{sharp_multiscale_prime_windows/check_windows.py}.\end{center}

\paragraph{Arithmetic tangents of the zero spectrum.}
For a primitive real character $\chi$ of conductor $q$ (conductor one gives $\zeta$) the kernel
$K_\chi$ is real, even, and positive semidefinite:
\[
 \sum_{i,j}c_i\overline{c_j}K_\chi(t_i-t_j)=\sum_\gamma m_\chi(\gamma)\frac{\abs{\sum_ic_ie^{i\gamma t_i}}^2}{1/4+\gamma^2}.
\]
It is
the covariance of a centred stationary Gaussian process $X_\chi$: take, for each distinct
positive ordinate, the coordinates $\sqrt{2m_\chi(\gamma)/(1/4+\gamma^2)}(\cos\gamma t,\sin\gamma t)$, and
$2\sqrt{m_\chi(0)}$ for a central zero, and multiply by independent standard normal variables.
These coordinates also give a deterministic Hilbert-space realization with
$\langle v_\chi(t),v_\chi(u)\rangle=K_\chi(t-u)$. By \cref{zer:lem:ramp}, $K_\chi(\log x)$ is the
normalized prime error $\sqrt x(S_\chi(x)-\delta_\chi\log x-c_\chi)$ up to $O((1+\log x)/\sqrt x)$. Put
\[
 d_q=1-\zergE-\log(2\pi)+\log q,\qquad J_\chi(a)=\begin{cases}\chi(n)\Lambda(n)/\sqrt n,&e^a=n\in\N,\ n\ge2,\\0,&\text{otherwise},\end{cases}
\]
and $B(u,v)=-\tfrac12(u\log u+v\log v-\abs{u-v}\log\abs{u-v})$ with $0\log0=0$.

\begin{lemma}[Cusps of the kernel; \stG]\label{zer:lem:cusp}
As $t\downarrow0$, $K_\chi(t)=M_\chi+\tfrac t2(\log t-d_q)+O_\chi(t^2)$. Near each $a>0$,
$K_\chi(t)=G_a(t)+J_\chi(a)(t-a)_+$ with $G_a$ of class $C^2$. For $q>1$ and parity $b$, for $t\ge0$,
\begin{align*}
 K_\chi(t)={}&K_\zeta(t)+(M_\chi-M_\zeta)\cosh(t/2)-\log q\,\sinh(t/2)
 +(2-b)t\sinh(t/2)\\
 &+2b\cosh(t/2)\log\cosh(t/2)-\sum_{n\le e^t}\frac{2\Lambda(n)(1-\chi(n))}{\sqrt n}\sinh\Bigl(\frac{t-\log n}2\Bigr).
\end{align*}
\end{lemma}
\begin{proof}
For $0<t<\log2$ the prime sum in \eqref{zer:eq:zetaexact} is empty, and expanding
$1-e^{-t}=t(1-t/2+O(t^2))$ in $\mathcal T$ gives the zeta case. The prime power $n$ enters
\eqref{zer:eq:zetaexact} through $2\Lambda(n)n^{-1/2}\sinh((t-\log n)/2)\mathbf 1_{t\ge\log n}$, whose linear
part is $J_\zeta(\log n)(t-\log n)_+$ and whose remainder is $C^2$. For $\chi$, let
$F_\chi(s)=(q/\pi)^{(s+b)/2}\Gamma((s+b)/2)L(s,\chi)$ and $H_\chi(s)=F_\chi'/F_\chi(\tfrac12+s)=\sum_\gamma s/(s^2+\gamma^2)$
for $s>\tfrac12$, pairing $\pm\gamma$ (a central zero contributes $m/s$); the same holds for $\xi$. Pairing
gives $\int_0^\infty e^{-st}K_\chi(t)\dd t=(sM_\chi-H_\chi(s))/(s^2-\tfrac14)$. The Dirichlet series and
\[
 \tfrac12\bigl[\psi(\tfrac{z+1}2)-\psi(\tfrac z2)\bigr]=\int_0^\infty\frac{e^{-zu}}{1+e^{-u}}\dd u,\qquad z=\tfrac12+s,
\]
give $H_\chi-H_\zeta=\tfrac12\log q+\int_0^\infty e^{-st}\bigl(b/(2\cosh(t/2))-2\cosh(t/2)\bigr)\dd t
+\sum_n\Lambda(n)(1-\chi(n))n^{-1/2-s}$, where $2\cosh(t/2)$ inverts the terms $1/z+1/(z-1)$ of
$\xi'/\xi$. Inverting with $(s^2-\tfrac14)^{-1}\leftrightarrow2\sinh(t/2)$ gives the formula, and the
convolution integral evaluates to the two elementary terms. For $t<\log2$ the prime sum is empty
and the local expansion follows from the zeta case; at $\log n$ the two prime contributions
combine to $2\chi(n)\Lambda(n)n^{-1/2}\sinh((t-\log n)/2)\mathbf 1_{t\ge\log n}$.
\end{proof}

Consequently $V_\chi(t)=2(K_\chi(0)-K_\chi(t))=t(\log(1/t)+d_q)+O_\chi(t^2)$, and $X_\chi$ has a
continuous version.

\begin{theorem}[Brownian tangents with an arithmetic correction; \stG]\label{zer:thm:tangent}
Fix $U>0$ and distinct reals $a_1,\dots,a_m$, and let $A_\chi(\boldsymbol a)$ be the symmetric matrix
with zero diagonal and entries $J_\chi(\abs{a_i-a_j})$. As $h\downarrow0$, uniformly for $u,v\in[0,U]$,
\begin{multline*}
 \frac1h\operatorname{Cov}\bigl(X_\chi(a_i+hu)-X_\chi(a_i),X_\chi(a_j+hv)-X_\chi(a_j)\bigr)\\
 =\delta_{ij}\bigl[(\log(1/h)+d_q)\min(u,v)+B(u,v)\bigr]-A_\chi(\boldsymbol a)_{ij}\min(u,v)+O(h).
\end{multline*}
The processes $Y_{i,h}(u)=(X_\chi(a_i+hu)-X_\chi(a_i))/\sqrt{h\log(1/h)}$ converge jointly in
$C([0,U])^m$ to independent standard Brownian motions. Moreover, for $a=\log n$ and $\abs v\le M$,
\[
 \frac{2K_\chi(a+vh)-K_\chi(a+(v+1)h)-K_\chi(a+(v-1)h)}h=-\frac{\chi(n)\Lambda(n)}{\sqrt n}(1-\abs v)_++O(h),
\]
and the numerator is $O(h^2)$ at a fixed positive lag outside the logarithms of prime powers
with $\chi(n)\ne0$.
\end{theorem}
\begin{proof}
By stationarity the covariance is
$K(\Delta+h(u-v))-K(\Delta+hu)-K(\Delta-hv)+K(\Delta)$ with $\Delta=a_i-a_j$. For $i=j$ it is
$\tfrac12(V(hu)+V(hv)-V(h\abs{u-v}))$, and the local expansion with $u+v-\abs{u-v}=2\min(u,v)$
gives the diagonal, with error $O(h^2)$. For $i\ne j$ use \cref{zer:lem:cusp} near $\Delta$ (the
even reflection has the same jump at negative lags): the $C^2$ part contributes $O(h^2uv)$, and the
ramp contributes $A_{ij}h((u-v)_+-u)=-A_{ij}h\min(u,v)$. After division by $h\log(1/h)$ the
covariance tends to $\delta_{ij}\min(u,v)$, which gives finite-dimensional convergence of these
centred Gaussian processes, and $\mathbb E\abs{Y_{i,h}(u)-Y_{i,h}(v)}^4=3(V(h\abs{u-v})/(h\log(1/h)))^2
\ll\abs{u-v}^{3/2}$ gives tightness. The triangular profile follows from
$2v_+-(v+1)_+-(v-1)_+=-(1-\abs v)_+$.
\end{proof}

The Gaussian process is built from the zero spectrum; nothing is asserted about the
distribution of prime errors themselves. Take now $a_i=\log i$, $1\le i\le N$, and write
$A_{\chi,N}$ for the matrix: for $i<j$ its entry is $\chi(p)^k\log p/p^{k/2}$ if $j=ip^k$, and $0$
otherwise. Let $V_\chi(h)=2(K_\chi(0)-K_\chi(h))$,
$R_{\chi,N}(h)_{ij}=\bigl(2K_\chi(\log(i/j))-K_\chi(\log(i/j)+h)-K_\chi(\log(i/j)-h)\bigr)/V_\chi(h)$, the
correlation matrix of $N$ increments of length $h$, and $\ell_q(h)=\log(1/h)+d_q$. Let $A_N^{(q)}$ be
the zeta matrix $A_{\zeta,N}$ with every edge at a power of a prime dividing $q$ removed.

\begin{theorem}[Determinants to every logarithmic order; \stG]\label{zer:thm:det}
For fixed $\chi$ and $N$, as $h\downarrow0$, $\ell_q(h)(I-R_{\chi,N}(h))=A_{\chi,N}+O(h)$ in operator
norm, $R_{\chi,N}(h)$ is positive definite for small $h$, and for $\ell_q(h)>\norm{A_{\chi,N}}$
\[
 \log\det R_{\chi,N}(h)=-\sum_{r\ge2}\frac{\tr(A_{\chi,N}^r)}{r\,\ell_q(h)^r}+O\Bigl(\frac h{\ell_q(h)}\Bigr),
\]
with a convergent series. There is a diagonal matrix $D_\chi$ with entries $\pm1$ such that
$A_{\chi,N}=D_\chi A^{(q)}_ND_\chi$. Consequently every coefficient of the expansion depends only on
the set of primes dividing $q$, and if $\chi,\psi$ have conductors $q,r$ with the same prime
divisors,
\[
 \log\det R_{\chi,N}(q\eps)-\log\det R_{\psi,N}(r\eps)=O\Bigl(\frac\eps{\log(1/\eps)}\Bigr)\qquad(\eps\downarrow0).
\]
Explicitly, with
\[
 S_2(N;q)=\sum_{\substack{p^k\le N\\p\nmid q}}\Bigl\lfloor\frac N{p^k}\Bigr\rfloor\frac{(\log p)^2}{p^k},\qquad
 S_3(N;q)=\sum_{\substack{p^k\le N\\p\nmid q}}(k-1)\Bigl\lfloor\frac N{p^k}\Bigr\rfloor\frac{(\log p)^3}{p^k},
\]
\[
 \log\det R_{\chi,N}(h)=-\frac{S_2(N;q)}{\ell_q(h)^2}-\frac{2S_3(N;q)}{\ell_q(h)^3}
 +O\bigl(\ell_q(h)^{-4}+h/\ell_q(h)\bigr).
\]
\end{theorem}
\begin{proof}
With $u=v=1$, $B(1,1)=0$ and \cref{zer:thm:tangent}, the covariance matrix of the increments is
$h(\ell I-A+E_h)$ with $\norm{E_h}=O(h)$, and $V_\chi(h)=h(\ell+O(h))$; hence
$R=I-A/\ell+F_h$ with $\norm{F_h}=O(h/\ell)$. For large $\ell$ the smallest eigenvalues of $I-A/\ell$
and $R$ exceed $\tfrac12$ and $\tfrac14$; the derivative of $\log\det$ along the segment gives
$\abs{\log\det R-\log\det(I-A/\ell)}\le4N\norm{F_h}$, and
$\log\det(I-A/\ell)=-\sum_{r\ge1}\tr(A^r)/(r\ell^r)$ with vanishing $r=1$ term. For the gauge write
$n=n_{\mathrm{ram}}n_{\mathrm{unr}}$ with the primes of $n_{\mathrm{ram}}$ dividing $q$ and
$(n_{\mathrm{unr}},q)=1$, and let $D_\chi=\operatorname{diag}(\chi(n_{\mathrm{unr}}))_{n\le N}$; if $j=ip^k$
with $p\nmid q$ then $\chi(i_{\mathrm{unr}})\chi(j_{\mathrm{unr}})=\chi(p)^k$, and edges at $p\mid q$ vanish
in both matrices. Similar matrices have equal traces of powers, and
$\ell_q(q\eps)=\ell_r(r\eps)=\log(1/\eps)+1-\zergE-\log(2\pi)$. Finally each edge $\set{i,ip^k}$
contributes twice to $\tr(A^2)$, so $\tr(A^2)=2S_2$; in a triangle $i<j<l$ the ratios $j/i=p^a$,
$l/j=r^b$ and $l/i$ are prime powers only if $p=r$, so the triangles are $i,ip^a,ip^{a+b}$, each
traversed six times, and $\tr(A^3)=6S_3$.
\end{proof}

For $\zeta$ and $N=8$ one finds $\tr(A^2)/2=3.1248179\ldots$ and
$\tr(A^3)/3=0.4995370\ldots$. The convergence in
$1/\log(1/h)$ is slow, as the numerical record
\begin{center}\artifact{grh_prime_resonance_2026-11-26/results/verification.json}\end{center}
shows. At higher orders
$\tr(A^r)$ is the weighted number of closed walks of length $r$ in the graph whose edges join $i$
and $ip^k$.

\begin{theorem}[Coalescing Gram determinants; \stG]\label{zer:thm:gramcoal}
Fix $x_0<x_1<\dots<x_n$, let $\Delta_j=x_j-x_{j-1}$ and $L=\log(1/h)$. As $h\downarrow0$,
\[
 \det\bigl(K_\chi(h(x_i-x_j))\bigr)_{0\le i,j\le n}
 =M_\chi(hL)^n\prod_{j=1}^n\Delta_j\Bigl[1+\frac{nd_q-\sum_j\log\Delta_j}L+O(L^{-2})\Bigr].
\]
\end{theorem}
\begin{proof}
Replace $v_\chi(hx_j)$, $j\ge1$, by the consecutive differences, a change of determinant one. The
top-left entry is $M_\chi$; by \cref{zer:lem:cusp} the lower block is
$h(L\operatorname{diag}(\Delta_j)+Q+O(h))$ with $Q_{jj}=\Delta_j(d_q-\log\Delta_j)$, the leading
off-diagonal terms cancelling because disjoint consecutive intervals are orthogonal for
$\min(u,v)$; the top-right entries are $O(hL)$. The Schur complement changes the block by
$O(h^2L^2)$, and $\det(D+Q/L+O(hL))=\det D\,[1+\tr(D^{-1}Q)/L+O(L^{-2})]$.
\end{proof}

This kernel is the prime error itself; the autocorrelation of the prime error at merging scales
behaves differently and is treated next.

\subsection{Merging scales: a cubic logarithmic law}\label{zer:sec:coalesce}
\sources{prime_scale_coalescence_2026-11-26; dirichlet-zero-multiplicity-and-reconstruction; grh-two-open-questions; dirichlet-coupled-zeros-and-prime-variances}

\Cref{zer:thm:coercive} keeps the scales fixed. For the zeta function the rate at which the
covariance degenerates as the scales merge is determined exactly up to a bounded factor. Let
\[
 E(t)=e^{t/2}\bigl(S(e^t)-t+\zergE+1\bigr),\qquad S(x)=\sum_{n\le x}\Lambda(n)\Bigl(\frac1n-\frac1x\Bigr),
\]
let $\gamma>0$ run over the distinct positive ordinates of zeta zeros, $m_\gamma$ the
multiplicities, $N(T)=\sum_{0<\gamma\le T}m_\gamma$ and $M_2(T)=\sum_{0<\gamma\le T}m_\gamma^2$. Since
$\zeta(\tfrac12)\ne0$ (the alternating eta series is positive at $\tfrac12$), \cref{zer:lem:ramp}
gives $E(t)=2\sum_{\gamma>0}m_\gamma\cos(\gamma t)/(\tfrac14+\gamma^2)+O((1+t)e^{-t/2})$, and by
\cref{zer:lem:ap}(i) the long-time covariance of $E(t+h_i)$ and $E(t+h_j)$ is $C(h_i-h_j)$, where
\[
 C(h)=2\sum_{\gamma>0}\frac{m_\gamma^2\cos(\gamma h)}{(1/4+\gamma^2)^2},\qquad C_0=C(0),\qquad
 D_2=2\sum_{\gamma>0}\frac{m_\gamma^2\gamma^2}{(1/4+\gamma^2)^2}.
\]
The multiplicity is squared because coincident frequencies are combined first. Fix $r\ge3$
distinct reals $u_1,\dots,u_r$, put $G_\eps=(C(\eps(u_i-u_j)))_{i,j}$, $\tau_\eps=\eps^3\log(1/\eps)$,
$\bar u$ the mean of the $u_i$, $v=u-\bar u\mathbf 1$, and
$W=\set{c\in\R^r:c\cdot\mathbf 1=c\cdot u=0}$. For $c\in W$ let
\[
 H_c(t)=\sum_ic_i(t-u_i)_+,\qquad \mathcal B(c)=\int_\R H_c(t)^2\dd t=\frac1{12}\sum_{i,j}c_ic_j\abs{u_i-u_j}^3,
\]
and let $b_1\ge\dots\ge b_{r-2}>0$ be the eigenvalues of $\mathcal B$ on $W$.

\begin{theorem}[Arithmetic coalescence law; \stG]\label{zer:thm:coalesce}
The long-time limit is taken first, then $\eps\downarrow0$.
\begin{enumerate}[label=\textup{(\roman*)}]
\item Uniformly for $c$ in the unit sphere of $W$,
$\mathcal B(c)-o(1)\le c^{\mathsf T}G_\eps c/\tau_\eps\le1.3208\,\mathcal B(c)+o(1)$.
\item $\lambda_1(G_\eps)\sim rC_0$, $\lambda_2(G_\eps)\sim D_2\norm v^2\eps^2$, and for $1\le j\le r-2$
\[
 b_j\le\liminf_{\eps\downarrow0}\frac{\lambda_{j+2}(G_\eps)}{\tau_\eps}\le\limsup_{\eps\downarrow0}\frac{\lambda_{j+2}(G_\eps)}{\tau_\eps}\le1.3208\,b_j .
\]
\item $\det G_\eps\asymp_u\eps^{3r-4}\log(1/\eps)^{r-2}$ and the condition number is
$\asymp_u\eps^{-3}/\log(1/\eps)$.
\end{enumerate}
All conclusions persist after deleting any fixed finite set of positive ordinates with their
negatives; $C_0$ and $D_2$ change, while $\mathcal B$ and the factor $1.3208$ do not.
\end{theorem}

With two offsets only the first two statements of (ii) hold. No simplicity, independence or
pair-correlation conjecture is used; the factor $1.3208$ is an explicit comparison, not a
claim of optimality. The arithmetic input is Montgomery's theorem, proved under RH
\cite{Montgomery1973Pair}, which gives the classical bound $M_2(T)\le(\tfrac43+o(1))N(T)$
\cite{ConreyGhoshGonek1998Simple}. A cumulative bound does not compare nonmonotone test
functions; the following local form does.

\begin{lemma}[Local squared multiplicities; \stG]\label{zer:lem:localM2}
For fixed $0<a<b$,
\[
 b-a\le\liminf_{T\to\infty}\frac{M_2(bT)-M_2(aT)}{T\log T/(2\pi)}\le
 \limsup_{T\to\infty}\frac{M_2(bT)-M_2(aT)}{T\log T/(2\pi)}\le Q(f)\,(b-a)
\]
for every real even continuous $f\ge0$ with $f(0)=1$, $f,\widehat f\in L^1$ and $\widehat f(v)\le0$ for
$\abs v\ge1$, where $\widehat f(v)=\int f(x)e^{-2\pi ivx}\dd x$ and
$Q(f)=\widehat f(0)+\int_{-1}^1\abs v\widehat f(v)\dd v$. The Fej\'er kernel gives $Q=\tfrac43$, a
closed-form test gives $Q=\tfrac12+\tfrac1{\sqrt2}\cot\tfrac1{\sqrt2}=1.3274992963\ldots$, and a test of
Chirre, Gon\c calves and de Laat gives $Q\le1.3208$ \cite{ChirreGoncalvesdeLaat2020Pair}; so the
upper constant is at most $1.3208$. Moreover $M_2(T)=O(T\log(2+T))$.
\end{lemma}
\begin{proof}
Let $w(y)=4/(4+y^2)=\int_\R e^{-2\abs s}e^{isy}\dd s$, $D_T=T\log T/(2\pi)$, and for $I=(aT,bT]$
\[
 F_{I,T}(v)=D_T^{-1}\sum_{\gamma,\gamma'\in I}e^{iv(\gamma-\gamma')\log T}w(\gamma-\gamma'),
\]
each zero repeated with its multiplicity; $F_{I,T}\ge0$, being an integral of absolute squares. On each
compact subinterval of $(-1,1)$ the measures $F_{I,T}(v)\dd v$ converge weakly to
$(b-a)(\delta_0+\abs v\dd v)$. Indeed, for the prefix $(0,cT]$ Montgomery's theorem
\cite{Montgomery1973Pair} at height $X=cT$ gives the form factor
$(1+o(1))X^{-2\abs\beta}\log X+\abs\beta+o(1)$ uniformly on $\abs\beta\le1-\delta$; change variables to
$\beta=v\log T/\log X$ and multiply by $D_X/D_T\to c$, the peak having mass tending to one. Subtracting
the prefixes $c=b$ and $c=a$ leaves the pairs on opposite sides of $aT$, whose total absolute
contribution is $O_{a,b}(\log^3T)$ uniformly in $v$: each unit distance bin on either side holds
$O_b(\log T)$ zeros, and $\sum_{0\le j,k\le bT+1}(1+(j+k)^2)^{-1}=O_b(\log T)$.

The diagonal pairs contribute $M_2(bT)-M_2(aT)$. For $0<d<1$ and $f_d(x)=f(dx)$, positivity of $f$ in the
ordinate variable, Fourier inversion, $F_{I,T}\ge0$ and the sign of $\widehat f_d$ outside $[-d,d]$ give
\[
 \frac{M_2(bT)-M_2(aT)}{D_T}\le\int_\R\widehat f_d(v)F_{I,T}(v)\dd v\le\int_{-d}^d\widehat f_d(v)F_{I,T}(v)\dd v ,
\]
and the weak limit applies to the sharp cutoff, whose endpoints carry no limiting mass. Letting
$T\to\infty$ and then $d\uparrow1$ gives the bound $(b-a)Q(f)$. For the Fej\'er kernel $\widehat f$ is the
triangle $(1-\abs v)_+$ and $Q=\tfrac43$. For $g(t)=\cos(\sqrt2t)/(\sqrt2\sin(1/\sqrt2))$ on
$J=[-\tfrac12,\tfrac12]$, zero elsewhere, and $f=\abs{\widehat g}^2$, one has $\int g=1$, $f(0)=1$, $\widehat f$
supported in $[-1,1]$ and $Q(f)=\int_Jg(t)[g(t)+\int_J\abs{t-u}g(u)\dd u]\dd t$; the bracket has second
derivative $g''+2g=0$ inside $J$ and is even, hence constant, equal to $g(\tfrac12)+\tfrac12$. The
rational choice $g(t)=1+z(\tfrac1{12}-t^2)$ gives $Q=\tfrac43-z/90+z^2/189$, which is $\tfrac{531}{400}$ at
$z=\tfrac{21}{20}$. The test of \cite[proof of Theorem~1]{ChirreGoncalvesdeLaat2020Pair} is
$f(x)=\widehat h(x/R)/\widehat h(0)$ for their even polynomial times a Gaussian $h$ with $h(x)\le0$ for
$\abs x\ge R$ and $\widehat h\ge0$; then $Q(f)$ is their rigorously bounded quantity, at most $1.3208$.
The lower bound follows from $m_\gamma^2\ge m_\gamma$ and the Riemann--von Mangoldt formula, and the global
bound from dyadic summation.
\end{proof}

A bound for $M_2(T)$ alone would not bound these differences, since two upper bounds cannot be
subtracted; the proof uses the local pair correlation of one height interval.

\begin{proof}[Proof of \cref{zer:thm:coalesce}]
\emph{Spline identity.} For $c\in W$, $H_c$ vanishes outside $[\min u_i,\max u_i]$ by the two
moment conditions. With $J_c(t)=\sum_ic_i\abs{t-u_i}^3$ one has $J_c''=12H_c$ and
$H_c''=\sum_ic_i\delta_{u_i}$, so $12\int H_c^2=\sum_ic_iJ_c(u_i)$, which is the displayed cubic form;
it is positive for $c\ne0$, since $H_c=0$ would make all slope jumps vanish. By Parseval,
$\mathcal B(c)=\pi^{-1}\int_0^\infty\abs{p_c(x)}^2x^{-4}\dd x$ with $p_c(x)=\sum_ic_ie^{iu_ix}$.

\emph{Scaled measures.} Let $\nu_T=\frac{2\pi}{T\log T}\sum_{\gamma>0}m_\gamma^2\delta_{\gamma/T}$. By
\cref{zer:lem:localM2} and approximation by step functions,
$\int f\le\liminf\int f\dd\nu_T\le\limsup\int f\dd\nu_T\le1.3208\int f$ for continuous nonnegative
compactly supported $f$ on $(0,\infty)$. Since $\abs{p_c(x)}\le A_u\norm c\min(x^2,1)$ by the two
cancellations, and $\limsup\nu_T((0,\delta])=O(\delta)$,
$\limsup\int_R^\infty x^{-4}\dd\nu_T=O((1+\log R)/R^3)$, the comparison extends to
$f_c=\abs{p_c}^2/x^4$, also with $x^4$ replaced by $(x^2+1/(4T^2))^2$, uniformly over unit $c\in W$
by a finite net. Since
\[
 \frac{c^{\mathsf T}G_\eps c}{\tau_\eps}=\frac1\pi\int_0^\infty\frac{\abs{p_c(x)}^2}{(x^2+1/(4T^2))^2}\dd\nu_T(x),\qquad T=1/\eps,
\]
this proves (i).

\emph{Separating the scales.} With $K(x)=\cos x-1+x^2/2\in[0,\min(x^4/24,x^2/2)]$ and
$M_2(T)=O(T\log T)$, $\sum_{\gamma>T}m_\gamma^2/\gamma^2=O(\log T/T)$, so
$C(h)=C_0-\tfrac12D_2h^2+O(\abs h^3\log(1/\abs h))$. Centre the nodes, so $v=u$. Then
$G_\eps=C_0\mathbf 1\mathbf 1^{\mathsf T}+\eps^2A+O(\tau_\eps)$ with
$A=-\tfrac12D_2(u^{\circ2}\mathbf 1^{\mathsf T}+\mathbf 1(u^{\circ2})^{\mathsf T}-2uu^{\mathsf T})$, whose
compression to $\mathbf 1^\perp$ is $D_2uu^{\mathsf T}$ and to $W$ is $0$. In an orthonormal basis
$\mathbf 1/\sqrt r$, $u/\norm u$, $W$, the matrix has blocks $rC_0+O(\eps^2)$,
$D_2\norm u^2\eps^2+O(\tau_\eps)$ and $\tau_\eps Q_\eps$ with $\mathcal B-o(1)\preceq Q_\eps\preceq1.3208\,\mathcal B+o(1)$,
off-diagonal blocks $O(\eps^2)$ against the first and $O(\tau_\eps)$ between the last two (this
uses both cancellations). Eliminating the first block changes the $W$-block by $O(\eps^4)$, and the
second by $O(\tau_\eps^2/\eps^2)=o(\tau_\eps)$, through congruences tending to the identity. By
min--max, congruence by matrices with singular values tending to $1$ changes positive
eigenvalues by factors $1+o(1)$, which gives (ii), and (iii) follows by multiplying and dividing
eigenvalues. Deleting finitely many frequencies changes $N$ and $M_2$ by bounded amounts, and
each deleted frequency contributes $O(\eps^4)$ to $c^{\mathsf T}G_\eps c$ for $c\in W$.
\end{proof}

For fixed $c\ne0$ the mean square of $\sum_ic_iE(t+\eps u_i)$ thus has exactly three possible
orders: $C_0(\sum c_i)^2$ if $\sum c_i\ne0$; $\sim D_2(\sum u_ic_i)^2\eps^2$ if $\sum c_i=0\ne\sum u_ic_i$;
and between $\mathcal B(c)\tau_\eps$ and $1.3208\,\mathcal B(c)\tau_\eps$ asymptotically if $c\in W$. Further
moment cancellations do not lower the order, since $\mathcal B>0$ on $W$.

\begin{corollary}[Raw finite differences; \stG]\label{zer:cor:rawdiff}
For a fixed integer $d\ge2$, the long-time mean square $V_d(\eps)$ of
$\sum_{j=0}^d(-1)^{d-j}\binom djE(t+j\eps)$ satisfies
\[
 \mathcal B_d\le\liminf\frac{V_d(\eps)}{\tau_\eps}\le\limsup\frac{V_d(\eps)}{\tau_\eps}\le1.3208\,\mathcal B_d,\qquad
 \mathcal B_d=\frac1{12}\sum_{i,j=0}^d(-1)^{i+j}\binom di\binom dj\abs{i-j}^3>0,
\]
with $\mathcal B_2,\mathcal B_3,\mathcal B_4,\mathcal B_5=\tfrac23,1,\tfrac83,\tfrac{25}3$.
\end{corollary}
\begin{proof}
The difference coefficients annihilate constants and linear functions, so they lie in $W$.
\end{proof}

These are raw differences; a finite truncation of the zero series would give $O(\eps^{2d})$, so
no fixed table of zeros can confirm the cubic law.

\begin{theorem}[Saturation of linear prediction; \stG]\label{zer:thm:predict}
For a target node $k$ let $V_k^*(\eps)$ be the infimum over real $a_i$ $(i\ne k)$, which may depend
on $\eps$, of the long-time mean square of $E(t+\eps u_k)-\sum_{i\ne k}a_iE(t+\eps u_i)$, and
$B_k^*=\min\set{\mathcal B(c):c\in W,\ c_k=1}>0$. Then
\[
 B_k^*\le\liminf\frac{V_k^*(\eps)}{\tau_\eps}\le\limsup\frac{V_k^*(\eps)}{\tau_\eps}\le1.3208\,B_k^*,
\]
and an intercept does not help, since all observations have mean zero.
\end{theorem}
\begin{proof}
The projection $w_k$ of $e_k$ to $W$ is nonzero, since otherwise a nonzero affine function would
vanish at the other $r-1\ge2$ nodes. Then $V_k^*=1/(e_k^{\mathsf T}G_\eps^{-1}e_k)$ and
$B_k^*=1/(w_k^{\mathsf T}\mathcal B^{-1}w_k)$. The block elimination gives
$\tau_\eps G_\eps^{-1}=P_WQ_\eps^{-1}P_W+o(1)$, since $\tau_\eps/\eps^2\to0$; invert the bounds on $Q_\eps$.
\end{proof}

For $u=(-1,0,1)$, $W$ is spanned by $(1,-2,1)$, $\mathcal B(1,-2,1)=\tfrac23$, $b_1=\tfrac19$ and
$B^*_{\mathrm{middle}}=\tfrac16$. So $\lambda_3(G_\eps)/\tau_\eps$ lies asymptotically in
$[\tfrac19,1.3208/9]$, and the optimal prediction error of the middle value in
$[\tfrac16,1.3208/6]$ after division by $\tau_\eps$; the midpoint predictor $\tfrac12(E(t-\eps)+E(t+\eps))$
is optimal up to a factor $1+o(1)$, its excess error being
$[C_0+C(2\eps)-2C(\eps)]^2/(2[C_0+C(2\eps)])=O(\eps^4)$. These finite constants are checked in exact
rational arithmetic in
\begin{center}\artifact{prime_scale_coalescence_2026-11-26/verification/verify_geometry.py}.\end{center}

\begin{proposition}[Exact constants; \stG, with the asymptotic for $M_2$ as a hypothesis]\label{zer:prop:kappa}
If moreover $M_2(T)\sim\kappa T\log T/(2\pi)$, then $1\le\kappa\le1.3208$,
$G_\eps|_W/\tau_\eps\to\kappa\mathcal B$,
$\lambda_{j+2}(G_\eps)/\tau_\eps\to\kappa b_j$, $V_k^*(\eps)/\tau_\eps\to\kappa B_k^*$, the optimal predictor
converges to the minimizer defining $B_k^*$, and
$\det G_\eps/(\eps^{3r-4}\log(1/\eps)^{r-2})\to rC_0D_2\norm v^2\kappa^{r-2}\det\mathcal B$. Also
$C(h)=C_0-\tfrac12D_2h^2+\tfrac\kappa{12}\abs h^3\log(1/\abs h)+o(\abs h^3\log(1/\abs h))$.
\end{proposition}
\begin{proof}
Now $\nu_T$ converges vaguely to $\kappa\,\dd x$, and the tail bounds upgrade this to the quadratic
forms; the block elimination and inversion give the limits, the optimal contrast being
\[
 \frac{G_\eps^{-1}e_k}{e_k^{\mathsf T}G_\eps^{-1}e_k}\to\frac{\mathcal B^{-1}w_k}{w_k^{\mathsf T}\mathcal B^{-1}w_k}.
\]
For $C(h)$ scale the remainder $K$, using
$\int_0^\infty K(x)x^{-4}\dd x=\tfrac16\int_0^\infty(1-\cos x)x^{-2}\dd x=\pi/12$.
\end{proof}

By \cref{zer:thm:coalesceinverse} below, $\kappa=1$ holds exactly when simple zeros have density one;
squared multiplicities are controlled by the density of multiple zeros through Korolev's uniform
multiplicity tail. The mean-square derivative
of the limiting signal exists, its squared coefficients summing to $D_2$, but no second
derivative does, since $\sum m_\gamma^2\gamma^4/(\tfrac14+\gamma^2)^2=\infty$; this roughness at the scale
$\gamma\asymp1/\eps$, together with the density of zeros, produces the factor $\eps^3\log(1/\eps)$.

\paragraph{The converse: which stencils read multiplicities from one variance.}
The forward comparison has a converse, and the stencils for which it holds can be read off a finite
symbol. A \emph{stencil} is a nonzero $c\in W$ for distinct nodes $u_1<\dots<u_r$; translating the
nodes changes nothing below, so take $u_1\ge0$. Put $p_c(x)=\sum_ic_ie^{iu_ix}$,
\[
 K_c(x)=\frac{\abs{p_c(x)}^2}{x^4},\qquad
 \zerSym_c(\xi)=\sum_{i<j}c_ic_j(u_j-u_i)^{3+i\xi}\quad(\xi\in\R),\qquad
 \Phi(y)=\cos y-1+\frac{y^2}2 ,
\]
with $K_c$ extended continuously to $x=0$, and recall
$\mathcal B(c)=\pi^{-1}\int_0^\infty K_c(x)\dd x=\zerSym_c(0)/6$ from the spline identity in the proof of
\cref{zer:thm:coalesce}. For a positive locally finite Borel measure $A$ on $[0,\infty)$ with
$\int_1^\infty x^{-4}\dd A(x)<\infty$ put $A(T)=A([0,T])$, $U_c(h)=h^4\int K_c(hx)\dd A(x)$ and
$\tau(h)=h^3\log(1/h)$; an atom at $0$ contributes $O(h^4)$.

\begin{theorem}[Finite families of stencils; \stU]\label{zer:thm:stencilbank}
For a nonempty finite family of stencils the following property holds exactly when the symbols
$\zerSym_c$ have no common real zero: for all such $A$ and every finite $d\ge0$,
\[
 \frac{\pi A(T)}{T\log T}\longrightarrow d
 \quad\Longleftrightarrow\quad
 \frac{U_c(h)}{\tau(h)}\longrightarrow d\,\mathcal B(c)\quad\text{for every stencil in the family.}
\]
For the reverse implication the variances need only be sampled along one decreasing sequence
$h_n\to0$ with $h_{n+1}/h_n\to1$. Each $\mathcal B(c)=\zerSym_c(0)/6$ is positive.
\end{theorem}
\begin{proof}
The cancellations give $p_c(x)=O(\min(x^2,1))$, and expanding the squared modulus gives
\begin{equation}\label{zer:eq:stencilcomp}
 \abs{p_c(y)}^2=2\sum_{i<j}c_ic_j\Phi((u_j-u_i)y),
\end{equation}
the constant and quadratic corrections summing to zero because $\sum c_i=\sum u_ic_i=0$. Two
integrations by parts give $\int_0^\infty\Phi(x)x^{-4}\dd x=\frac16\int_0^\infty(1-\cos x)x^{-2}\dd x=\pi/12$,
so \eqref{zer:eq:stencilcomp} gives $\mathcal B(c)=\zerSym_c(0)/6$; it is positive since $K_c\ge0$ and
exponentials at distinct nodes are linearly independent.

Let $\mu_T(S)=\pi A(TS)/(T\log T)$ for Borel $S\subseteq[0,\infty)$ and $T>2$. A counting limit $d$ gives $\mu_T\to d\,\dd x$ vaguely, by
evaluation on intervals, and $A(T)\ll(1+T)\log(2+T)$. For $R\ge1$ the $\mu_T$-mass of $[R,2R]$ is
$O(R(1+\log 2R))$ uniformly for $T>2$, so dyadic summation gives
$\int_R^\infty x^{-4}\dd\mu_T(x)=O(R^{-3}(1+\log 2R))$; with $h=1/T$ and $K_c(x)=O(x^{-4})$ at infinity this
gives the forward implication.

Conversely, fix one stencil. For small $\delta>0$, $\abs{p_c}$ is bounded below on $[\delta/3,2\delta]$ by
analyticity and nontriviality. Taking mesh indices with $h_n\sim\delta/T$, $h_nx$ stays in that interval for
$T/2\le x\le T$, and $U_c(h_n)=O(\tau(h_n))$ gives $A([T/2,T])\ll T\log T$; dyadic summation gives the
cumulative growth bound above, also when the stencil has higher-order cancellation at the origin. By
local uniform boundedness the $\mu_T$ have vaguely convergent subsequences, and every limit $\nu$ has
$\nu([0,R])\ll R$ for $R>0$. For $s>0$ choose mesh indices with $Th_n\to s$; the kernels converge on
compacts, the tail bounds hold uniformly for $Th_n$ near $s$, $\log(1/h_n)/\log T\to1$, and
\[
 \frac{U_c(s/T)}{T^{-3}\log T}=\frac1\pi\int_0^\infty\abs{p_c(sx)}^2x^{-4}\dd\mu_T(x).
\]
Hence the signed measure $\sigma=\nu-d\,\dd x$ satisfies $\int\abs{p_c(sx)}^2x^{-4}\dd\sigma(x)=0$ for every
$s>0$ and each stencil of the family.

Put $R(s)=\int\Phi(sx)x^{-4}\dd\sigma(x)$. Since $\abs\sigma([0,X])\ll X$ and $\abs{\Phi(y)}\ll\min(y^4,y^2)$,
splitting at $1/\abs s$ and summing the tail gives $R(s)=O(\abs s^3)$; $R$ is continuous and even, and
\eqref{zer:eq:stencilcomp} gives $\sum_{i<j}c_ic_jR((u_j-u_i)s)=0$. In the variable $w=\log s$ this is a
finite translation equation for the bounded function $H(w)=e^{-3w}R(e^w)$; with
$\widehat f(\xi)=\int e^{-iw\xi}f(w)\dd w$, its distributional transform satisfies
$\zerSym_c(\xi)\widehat H(\xi)=0$. If the symbols have no common real zero, test against compactly supported
smooth functions divided by $\sum_c\abs{\zerSym_c}^2$ and sum the equations multiplied by
$\overline{\zerSym_c}$: then $\widehat H=0$ and $R=0$. Finally $R^{(4)}$ is the Fourier transform of the
symmetrization of $\sigma$ (each of the two reflections with weight $\tfrac12$), since
$\partial_s^4\Phi(sx)=x^4\cos(sx)$; truncation and testing against Schwartz functions justify this, the
growth bounds making the measures tempered. Fourier injectivity gives $\sigma=0$. Every subsequential
limit is therefore $d\,\dd x$, and evaluation on $[0,1]$ gives the counting limit.

For necessity let $\omega$ be a common real zero; $\omega\ne0$ since $\mathcal B(c)>0$. Substitution in
\eqref{zer:eq:stencilcomp} gives, for real $\xi$, with absolutely convergent integrals,
\begin{equation}\label{zer:eq:stencilmellin}
 \int_0^\infty K_c(y)y^{i\xi}\dd y=2\zerSym_c(-\xi)\int_0^\infty\Phi(y)y^{-4+i\xi}\dd y .
\end{equation}
On $x\ge1$ take $\dd A(x)=\pi^{-1}\log x\,[1+\eps\cos(\omega\log x)]\dd x$ with $0<\eps<1$. Its normalized
count is $1+\eps\Re\{T^{i\omega}/(1+i\omega)\}+o(1)$, which has no limit. Yet the substitution $y=hx$ gives
$U_c(h)/\tau(h)\to\mathcal B(c)$ for every stencil of the family: the oscillatory coefficient in
\eqref{zer:eq:stencilmellin} vanishes, $K_c(y)\abs{\log y}$ is integrable, and the integral over $0<y<h$ lost
by the cutoff tends to zero.
\end{proof}

\begin{remark}\label{zer:rem:stencilexist}
For a family with no common zero, mere existence of all the finite limits of $U_c(h)/\tau(h)$ forces
them to be proportional to $\mathcal B(c)$. Repeat the compactness argument with the individual limits:
each tangent $\nu$ gives the same value as its dilation $\nu(s\,\cdot)/s$ in every scaled kernel
integral, both being tangents of the same family, and the injectivity argument applied to their
difference gives $\nu(s\,\cdot)=s\nu$ for every $s>0$. So $\nu$ has constant density on the half-line,
and any one measured limit fixes that constant for all tangents.
\end{remark}

\begin{corollary}[Three observations; \stU]\label{zer:thm:inversemeasure}
For $a,b>0$, every stencil on three nodes satisfies \cref{zer:thm:stencilbank}. For the nodes
$(0,a,a+b)$ the weights are proportional to $(b,-(a+b),a)$, and the second difference
$D_hf(t)=bf(t)-(a+b)f(t+ah)+af(t+(a+b)h)$ has
$\mathcal B=B_{a,b}=a^2b^2(a+b)/3$.
\end{corollary}
\begin{proof}
In $\zerSym_c$ the term of the longest lag strictly dominates the other two:
$ab(a+b)^3>b(a+b)a^3+(a+b)ab^3$ because $(a+b)^2>a^2+b^2$. Hence $\zerSym_c$ has no real zero, and
$\zerSym_c(0)=ab(a+b)\bigl((a+b)^2-a^2-b^2\bigr)=2a^2b^2(a+b)$. (Equivalently, $H_c$ is a tent of base $a+b$
and height $ab$.)
\end{proof}

With fewer than three nodes $W=0$, so three is the least number of observations for this recovery,
and four is the least at which it can fail. On the nodes $0,1,2,3$ the cancellations say exactly
that $\sum_{j=0}^3c_jz^j=(1-z)^2(a+bz)$, that is $c=(a,b-2a,a-2b,b)$. Reflection interchanges $a$ and
$b$; scaling and reflection preserve $t=ab/(a^2+b^2)\in[-\tfrac12,\tfrac12]$, and conversely $t$
determines the normalized response
\begin{equation}\label{zer:eq:fourresponse}
 \frac{\abs{p_c(x)}^2}{a^2+b^2}=\abs{1-e^{ix}}^4(1+2t\cos x).
\end{equation}
Call the shape $t$ \emph{determining} if its stencil satisfies \cref{zer:thm:stencilbank}.

\begin{theorem}[Four consecutive observations; \stU]\label{zer:thm:fourpoint}
Put $I=(\tfrac1{13},\tfrac2{11})$ and $\alpha=\log3/\log2$. Every $t\notin I$ is determining. For
$t\in I$ put $A_t=4-7t$, $B_t=8-32t$, $C_t=27t$; there are unique $u\in(0,\pi)$ and $v\in(-\pi,0)$
with $B_te^{iu}+C_te^{iv}=A_t$, and the shape fails to be determining precisely when
\begin{equation}\label{zer:eq:fourexception}
 \frac{v-\alpha u}{2\pi}\in\Z+\alpha\Z .
\end{equation}
These exceptional shapes form a countable dense subset of $I$, and each misses exactly one pair of
frequencies $\set{\pm\omega}$: its symbol vanishes at $\pm\omega$ and nowhere else on $\R$. For every
shape, $\mathcal B(c)=(a^2+b^2)(4+2t)/6$.
\end{theorem}
\begin{proof}
The lag autocorrelations $\sum_ic_ic_{i+k}$ for $k=1,2,3$ are $7ab-4(a^2+b^2)$, $a^2+b^2-4ab$ and $ab$,
so after division by $a^2+b^2$ the symbol is
\begin{equation}\label{zer:eq:foursymbol}
 \zerSym_t(\xi)=7t-4+8(1-4t)e^{i\xi\log2}+27te^{i\xi\log3},
\end{equation}
with $\zerSym_t(0)=4+2t$. For $t\le0$ the modulus of the second term exceeds the sum of the other two
by $4+2t$; for $t\ge\tfrac14$ the third exceeds the other two by $4+2t$. For $0<t<\tfrac14$ the moduli
$A_t,B_t,C_t$ are positive and
\[
 B_t+C_t-A_t=4+2t,\qquad B_t-C_t-A_t=4-52t,\qquad C_t-B_t-A_t=66t-12 ,
\]
so they form a strict triangle precisely on $I$. At its endpoints a zero would need the two phases
to be $(1,-1)$ or $(-1,1)$, impossible since $\alpha$ is irrational by unique factorization.

On $I$ the phases $(e^{i\xi\log2},e^{i\xi\log3})$ at a zero must form one of the two conjugate
triangle configurations. For the configuration $(e^{iu},e^{iv})$, lifting gives
$\xi\log2=u+2\pi m$ and $\xi\log3=v+2\pi n$ with integers $m,n$, which is \eqref{zer:eq:fourexception} in
both directions; irrationality of $\alpha$ makes $(m,n)$ unique, and conjugation gives the opposite
frequency. The cosine rule makes $u,v$ real analytic on $I$, with $(u,v)\to(0,-\pi)$ at the left
endpoint and $(u,v)\to(\pi,0)$ at the right, so $v-\alpha u$ is nonconstant on every open subinterval
and its level sets are discrete in $I$. The group $\Z+\alpha\Z$ is countable and dense, so its
preimage is countable and dense. Finally $\mathcal B(c)=\zerSym_c(0)/6=(a^2+b^2)(4+2t)/6$.
\end{proof}

Since $(2^{i\xi},3^{i\xi})$ is dense on the two-torus, $\inf_\xi\abs{\zerSym_t(\xi)}>0$ outside $\overline I$ by
strict domination, while on $\overline I$, its two determining endpoints included, the infimum is $0$
because the torus contains a cancelling triangle. So almost every shape in $I$ recovers the
limiting count although its symbol takes arbitrarily small values.

\begin{example}[An exceptional shape]\label{zer:ex:fourexceptional}
Let $F(t)=(v-\alpha u)/(2\pi)$ on $I$. Its endpoint values are $-\tfrac12$ and $-\tfrac\alpha2$, and the group
element $1-\alpha$ lies strictly between them because $2^3<3^2$ and $3<2^2$; the intermediate value
theorem gives an exceptional shape. For one solution $t_*$ the positive missed frequency is
$\omega_*=(2\pi-u)/\log2$, and interval arithmetic, enclosing opposite signs of $F(t)-(1-\alpha)$ at two
rational endpoints, gives
\begin{gather*}
 0.16788340058442558<t_*<0.16788340058442560,\\
 6.23399540333078<\omega_*<6.23399540333080 .
\end{gather*}
One stencil is $a=1$, $b=2t_*/(1+\sqrt{1-4t_*^2})$, with all four entries nonzero. The enclosures are in
\begin{center}\artifact{grh-two-open-questions/zero-stencils/checks.txt}.\end{center}
\end{example}

Two shapes suffice together. With $\widetilde V_t(h)=U_c(h)/(a^2+b^2)$, \eqref{zer:eq:fourresponse} holds
pointwise under the spectral integral, so distinct shapes $t_1\ne t_2$ recover the ordinary second
difference at each $h$ by $\widetilde V_0(h)=(t_2\widetilde V_{t_1}(h)-t_1\widetilde V_{t_2}(h))/(t_2-t_1)$; finite
normalized limits for both shapes give the counting limit, by the determining shape $t=0$, and fix both
limits. The coefficients show how additive errors are amplified when the shapes are close.

The following five-point contrast has a symbol with a real zero, and the necessity
construction of \cref{zer:thm:stencilbank} is then explicit. With $s=(1+\sqrt3)/2$ and $Q(z)$ the real
quadratic with coefficients $(s+\sqrt s)/2$, $1/(2s)$, $(s-\sqrt s)/2$, one has
$\abs{Q(z)}^2=1+\abs{1+z}^4/8$ on $\abs z=1$, and the coefficients of $(1-z)^2Q(z)$ annihilate affine
functions, with squared modulus $\abs{1-z}^4+\abs{1-z^2}^4/8$. So $U_c(h)=U_{(1,-2,1)}(h)+U_{(1,-2,1)}(2h)/8$, its
symbol is $(1+2^{i\xi})\zerSym_{(1,-2,1)}(\xi)$, and it vanishes at $\omega=\pi/\log2$: for
$\dd A(x)=\log(2+x)(1+\eps\cos(\omega\log x))\dd x$ on $x\ge1$, replacing $h$ by $2h$ reverses the
oscillation of $U_{(1,-2,1)}(h)/\tau(h)$, so the five-point normalized variance converges while
$A(T)/(T\log T)=1+\eps\Re\{T^{i\omega}/(1+i\omega)\}+o(1)$ does not.

When $x^{-4}$ is replaced by $(\tfrac14+x^2)^{-2}$ and $\int(1+x^2)^{-1}\dd A<\infty$, the change costs
$O(h^4)=o(\tau(h))$, since $\abs{p_c(hx)}^2\ll h^4x^4$ and $0\le1-x^4/(\tfrac14+x^2)^2\ll(1+x^2)^{-1}$. Write
$V_c(h)$ for the long-time mean square of $\sum_ic_iE(t+hu_i)$ and $\kappa(T)=2\pi M_2(T)/(T\log T)$.

\begin{theorem}[Reading multiplicities from the variance; \stG]\label{zer:thm:coalesceinverse}
\begin{enumerate}[label=\textup{(\roman*)}]
\item Let a finite family of stencils have symbols without a common real zero; for instance one
three-point stencil, or one four-point stencil on four equally spaced nodes whose shape is not
exceptional in
\cref{zer:thm:fourpoint}. For every finite $\kappa\ge0$, $V_c(h)/\tau(h)\to\mathcal B(c)\kappa$ as $h\downarrow0$
for every stencil of the family if and only if $\kappa(T)\to\kappa$. On the left it suffices to take
$h=h_n$ along a decreasing sequence with $h_{n+1}/h_n\to1$, for instance $h=1/k$, and by
\cref{zer:rem:stencilexist} existence of the finite limits suffices. In particular, for a three-point
stencil with $a,b>0$, $V_{a,b}(h)/\tau(h)\to B_{a,b}\kappa$ if and only if $\kappa(T)\to\kappa$.
\item For every stencil $c$, with any nodes, the following are equivalent: simple zeros of
$\zeta$ have density one among all zeros counted with multiplicity; $M_2(T)/N(T)\to1$;
$c^{\mathsf T}G_\eps c/\tau_\eps\to\mathcal B(c)$. In the last condition it suffices to take $\eps=h_n$ along
any decreasing sequence with $h_n/h_{n+1}$ bounded, for instance $h_n=2^{-n}$.
\end{enumerate}
\end{theorem}

For $a=b=1$ the measured coefficient is $\tfrac23\kappa$: knowing that one such limit exists gives the
limiting coefficient for every contrast. The criterion does not decide whether it holds.

\begin{proof}
(i) By the spectral formula, $V_c(h)=2\sum_{\gamma>0}m_\gamma^2\abs{p_c(h\gamma)}^2/(\tfrac14+\gamma^2)^2$. Apply
\cref{zer:thm:stencilbank} to $A=2\sum_{\gamma>0}m_\gamma^2\delta_\gamma$, for which $\pi A(T)/(T\log T)=\kappa(T)$
and $\int(1+x^2)^{-1}\dd A<\infty$ by $M_2(T)=O(T\log(2+T))$ (\cref{zer:lem:localM2}), with the change of
denominator.

(ii) Let $D_j(T)$ count distinct zeros of multiplicity $j$ up to height $T$. Korolev's theorem gives
absolute $A_0,c_0>0$ with $D_j(T)\le A_0N(T)e^{-c_0j}$ for all $j\ge1$ and large $T$
\cite[Theorem~A]{Korolev2006Multiple}. With $d_T=1-D_1(T)/N(T)$ and $e_T=M_2(T)/N(T)-1$, weighting by
$j(j-1)$ against $j$ for $j\ge2$ and separating at $K$ gives
$d_T\le e_T\le Kd_T+A_0\sum_{j>K}j^2e^{-c_0j}$, which proves the equivalence of the first two
conditions (and $e_T=O(d_T\log(e/d_T))$). Replacing $m_\gamma^2$ by $m_\gamma$ in the spectral formula gives a
baseline with limit $\mathcal B(c)$ by the Riemann--von Mangoldt formula and the scaling of the proof
of \cref{zer:thm:coalesce}; the excess has weights $m_\gamma(m_\gamma-1)\ge0$, so the second condition
implies the third. Conversely, $p_c$ is a nonzero exponential polynomial whose
zeros do not accumulate at $0$; choose $z>0$ with $p_c\ne0$ on $(0,z]$. For large $T$ take the first
$n$ with $h_n\le z/T$, so $z/R\le h_nT\le z$ with $R=\sup h_n/h_{n+1}$; for $T/2<\gamma\le T$, $\abs{p_c(h_n\gamma)}$
is bounded below and $(\tfrac14+\gamma^2)^{-2}\gg T^{-4}$, and $\tau(h_n)\asymp T^{-3}\log T$, so the vanishing
normalized excess gives $\sum_{T/2<\gamma\le T}m_\gamma(m_\gamma-1)=o(T\log T)$; dyadic summation gives
$M_2(T)-N(T)=o(T\log T)$.
\end{proof}

\paragraph{What an exceptional variance still determines.}
At an exceptional four-point shape one variance no longer fixes the count, but it fixes everything
except one sinusoidal term in $\log$ height.

\begin{theorem}[One missing frequency; \stU]\label{zer:thm:missingfreq}
Let $t$ be an exceptional shape of \cref{zer:thm:fourpoint} with positive missed frequency $\omega$, and
suppose $U_c(h)/\tau(h)\to d\,\mathcal B(c)$ along a decreasing sequence $h_n\to0$ with $h_{n+1}/h_n\to1$.
Every subsequential vague limit of $\mu_T(S)=\pi A(TS)/(T\log T)$, as $T\to\infty$, has the form
\begin{equation}\label{zer:eq:tangentform}
 [d+\Re(Cx^{i\omega})]\dd x,\qquad C\in\C,\qquad\abs C\le d .
\end{equation}
There is $r\in[0,d]$ such that the set of subsequential limits of $a(T)=\pi A(T)/(T\log T)$ is the
interval $[d-r/\sqrt{1+\omega^2},\,d+r/\sqrt{1+\omega^2}]$. With $L=\exp(\pi/\omega)$ moreover
$a(T)+a(LT)\to2d$.
\end{theorem}
\begin{proof}
The compactness and translation argument of \cref{zer:thm:stencilbank} applies to every tangent $\nu$.
For $\sigma=\nu-d\,\dd x$ the bounded function $H$ has Fourier transform supported on $\set{\pm\omega}$. A
distribution supported on finitely many points is a sum of derivatives of delta masses, so $H$ is a sum
of exponentials at these frequencies times polynomials; boundedness forces the polynomials to be
constant, since the top coefficients form a nonzero periodic trigonometric polynomial bounded away from
zero along some sequence of period translates; reality gives conjugate coefficients.

For the measures $\Re(Cx^{i\omega})\dd x$ the map $\sigma\mapsto H$ gives $\Re(CJe^{-i\omega w})$ with
$J=\int_0^\infty\Phi(x)x^{-4+i\omega}\dd x$, and $J\ne0$: otherwise $R$ for $\cos(\omega\log x)\dd x$ would vanish
identically, contradicting the injectivity of $\sigma\mapsto R$ shown with its fourth derivative. Hence
this two-dimensional real family realizes every possible $H$; injectivity gives \eqref{zer:eq:tangentform},
and positivity gives $\abs C\le d$.

The set of tangents is compact for vague convergence, and $C$ depends continuously on the tangent,
since two compactly supported test integrals distinguish the independent functions
$\cos(\omega\log x)$ and $\sin(\omega\log x)$. Let $r$ be the largest modulus attained. Replacing $T$ by $sT$
replaces $\nu(S)$ by $\nu(sS)/s$ and rotates $C$ by $s^{i\omega}$, so the full circle of radius $r$ occurs.
Evaluation on $[0,1]$ gives $d+\Re(C/(1+i\omega))$, which proves the interval statement. For $s=L$ the
rotation negates $C$, so $a(T)+a(LT)\to2d$ along every tangent subsequence, hence along all heights by
compactness.
\end{proof}

\begin{proposition}[The first missing frequency; \stU]\label{zer:prop:firstres}
The least positive frequency missed by any four-point shape is the number $\omega_*$ of
\cref{zer:ex:fourexceptional}.
\end{proposition}
\begin{proof}
On $I$ one has $A_t=B_t/2+C_t/3$ with $B_t,C_t>0$, so a zero of \eqref{zer:eq:foursymbol} requires
$2^{i\omega}-\tfrac12$ and $3^{i\omega}-\tfrac13$ to point in opposite directions. Write $L=2\pi/\log3$ and
$U=3\pi/(2\log2)$. For $0<\omega\le L$ the arguments $\omega\log2$ and $\omega\log3$ bound an arc of length at
most $2\pi\log(3/2)/\log3<\pi$; the positive combination $B_t2^{i\omega}+C_t3^{i\omega}$ points strictly inside
that arc, whose interior excludes the positive real axis, so this range contains no missed frequency.

Let $\theta_p(\omega)$ be the continuous argument of $p^{i\omega}-1/p$ starting at $0$. With $\phi=\omega\log p$,
\[
 \theta_p'=(\log p)\frac{1-p^{-1}\cos\phi}{1-2p^{-1}\cos\phi+p^{-2}} .
\]
On $[L,U]$ the angle $\omega\log2$ lies between $\pi$ and $3\pi/2$, so $\theta_2'\le\log2$, while the same
formula gives $\theta_3'\ge\tfrac34\log3>\log2$; thus $\theta_3-\theta_2$ strictly increases. At $L$ its value
lies in $(0,\pi)$, since $\theta_3(L)=2\pi$ and $\pi<\theta_2(L)<2\pi$. At $U$, $\theta_2(U)=\pi+\arctan2$, whereas
$\theta_3(U)=2\pi+\arg(e^{i\psi}-\tfrac13)$ with $\psi=(3\alpha/2-2)\pi\in(\pi/3,\pi/2)$, using $3^9>2^{14}$ and
$3^3<2^5$; here $\sin\psi>2(\cos\psi-\tfrac13)$, so $\pi<\theta_3(U)-\theta_2(U)<2\pi$. The unique crossing of
$\pi$ is the first possible opposition, and the pair $(t_*,\omega_*)$ of \cref{zer:ex:fourexceptional}
realizes a zero in $(L,U)$, hence gives this crossing.
\end{proof}

\begin{proposition}[Integer multiplicities with the counting main terms; \stU]\label{zer:prop:discrete}
For every exceptional shape, with missed frequency $\omega$, every $0<q_0<1$ and every
$0<\eps\le\min(q_0,1-q_0)$, there is a locally finite set of positive ordinates with multiplicities
$m_x\in\set{1,2}$ such that
\[
 \sum_{x\le T}m_x=\frac T{2\pi}\log\frac T{2\pi}-\frac T{2\pi}+O(1),\qquad
 \frac{2\pi}{T\log T}\sum_{x\le T}m_x^2=1+q_0+\eps\Re\frac{T^{i\omega}}{1+i\omega}+o(1),
\]
and the spectral variance $2\sum_xm_x^2\abs{p_c(hx)}^2/(\tfrac14+x^2)^2$ divided by $\tau(h)$ tends to
$(1+q_0)\mathcal B(c)$. So the factor $1/\sqrt{1+\omega^2}$ of \cref{zer:thm:missingfreq} is attained by an
integer spectrum whose ordinary count has both Riemann--von Mangoldt main terms.
\end{proposition}
\begin{proof}
For $x\ge10$ put $\lambda(x)=\log(x/2\pi)/(2\pi)$ and $q(x)=q_0+\eps\cos(\omega\log x)$, and define the strictly
increasing unbounded functions
\[
 F_1(T)=\int_{10}^T(1-q(x))\lambda(x)\dd x,\qquad F_2(T)=\frac12\int_{10}^Tq(x)\lambda(x)\dd x .
\]
Place simple points at the crossings of the positive integers by $F_1$ and double points at the
crossings of the positive elements of $\Z+\theta$ by $F_2$, with $0<\theta<1$ outside the countable set of
fractional parts of $F_2(F_1^{-1}(n))$, $n\ge1$. The two sets are disjoint and their counts differ from
$F_1,F_2$ by $O(1)$. So the ordinary count is $\int_{10}^T\lambda+O(1)$ and the squared count is
$\int_{10}^T(1+q)\lambda+O(1)$, and direct integration gives the two formulas. The continuous
squared-counting density $2(1+q)\lambda$ has variance limit $(1+q_0)\mathcal B(c)$ by
\eqref{zer:eq:stencilmellin}, since the symbol vanishes at $\pm\omega$ and the constant in $\log(x/2\pi)$
contributes $O(h^3)$. Replacing it by $2\sum_xm_x^2\delta_x$ changes the cumulative function by $O(1)$;
integration by parts against $h^4K_c(hx)$ costs $O(h^4\int_0^\infty\abs{K_c'(y)}\dd y)=O(h^4)$, as $K_c$ is
smooth at $0$ and $K_c'=O(y^{-4})$ at infinity, and the change of denominator also costs $O(h^4)$.
\end{proof}

With $q_0=\eps=0.1604$ the density $1+q$ lies in $[1,1.3208]$, so the model keeps the local limiting
bounds of \cref{zer:lem:localM2} on every fixed interval while attaining the full half-width those
bounds allow. The ordinates are a model spectrum; an Euler product realizing them would be an
additional arithmetic constraint.

\begin{corollary}[Exceptional four-point variances for $\zeta$; \stG]\label{zer:cor:missingzeta}
Let $c$ be a four-point stencil, exceptional with missed frequency $\omega$, and suppose
$V_c(h)/\tau(h)\to d\,\mathcal B(c)$ along a decreasing sequence with $h_{n+1}/h_n\to1$. Then the lower and
upper limits of $\kappa(T)$ are $d\mp r/\sqrt{1+\omega^2}$ with
\[
 r\le\min(d-1,\,1.3208-d),
\]
so the variance fixes the midpoint $d$ of the limiting normalized counts, and $\kappa(T)$ converges if
$d=1$ or $d=1.3208$; by \cref{zer:thm:coalesceinverse}(ii), $d=1$ means that simple zeros have density
one. Together with \cref{zer:thm:coalesceinverse}(i): if any one four-point variance on four equally
spaced nodes has a normalized
limit along such a sequence, then
\[
 \limsup_{T\to\infty}\kappa(T)-\liminf_{T\to\infty}\kappa(T)\le\frac{2\cdot0.1604}{\sqrt{1+\omega_*^2}},\qquad
 \frac{0.1604}{\sqrt{1+\omega_*^2}}=0.025405106245\ldots,
\]
and this half-width is attained, by \cref{zer:prop:discrete}, among spectra satisfying the local bounds of
\cref{zer:lem:localM2}.
\end{corollary}
\begin{proof}
By the proof of \cref{zer:thm:coalesceinverse}(i), \cref{zer:thm:missingfreq} applies to
$A=2\sum_{\gamma>0}m_\gamma^2\delta_\gamma$. By \cref{zer:lem:localM2} every tangent lies between $\dd x$ and
$1.3208\,\dd x$ on each fixed positive interval; these interval bounds bound the continuous density
\eqref{zer:eq:tangentform} everywhere on $(0,\infty)$, and its extreme values are $d\pm\abs C$. This gives the
bound on $r$, and $r\le(1.3208-1)/2=0.1604$. At a determining shape $\kappa(T)$ converges, and at an
exceptional one $\omega\ge\omega_*$ by \cref{zer:prop:firstres}. The numerical value is enclosed in
\artifact{grh-two-open-questions/zero-stencils/checks.txt}.
\end{proof}

\paragraph{Five equally spaced observations.}
On the nodes $0,1,2,3,4$ write $P(z)=\sum_{j=0}^4c_jz^j$ for a stencil, so that $P(1)=P'(1)=0$, and put
$b_k=k^3\sum_{j=0}^{4-k}c_jc_{j+k}$, so that
$\zerSym_c(\xi)=b_1+b_22^{i\xi}+b_33^{i\xi}+b_44^{i\xi}$. A \emph{missing positive frequency} is an $\omega>0$
with $\zerSym_c(\omega)=0$; by \cref{zer:thm:stencilbank} the stencil is determining exactly when there is
none, and at a missing frequency the perturbation $\cos(\omega\log x)$ of the leading density is invisible
to the variance. Put $\alpha=\log3/\log2$.

\begin{theorem}[Five consecutive observations; \stU]\label{zer:thm:fivepoint}
If $b_3=0$, there are infinitely many missing positive frequencies when $b_1+b_2z+b_4z^2$ has a root on the
unit circle, and none otherwise. If $b_3\ne0$, solve
\begin{equation}\label{zer:eq:fivequad}
 4b_1b_4X^2+2b_2(b_1+b_4)X+(b_1-b_4)^2+b_2^2-b_3^2=0 .
\end{equation}
For each distinct real root $-1<X<1$ put $z=X+i\sqrt{1-X^2}=e^{i\theta}$ with $0<\theta<\pi$, and
$w=-(b_1+b_2z+b_4z^2)/b_3=e^{i\phi}$; count the root once exactly when
\begin{equation}\label{zer:eq:fivegroup}
 \frac{\phi-\alpha\theta}{2\pi}\in\Z+\alpha\Z .
\end{equation}
The resulting count is the exact number of missing positive frequencies, and it is at most two. If
$\phi-\alpha\theta=2\pi(\alpha m-n)$ with integers $m,n$, the corresponding frequency is $\abs{\theta+2\pi m}/\log2$.
\end{theorem}

Padding a four-node stencil with $c_4=0$ gives $b_4=0$; then \eqref{zer:eq:fivequad} becomes the
cosine rule of the proof of \cref{zer:thm:fourpoint}, and \eqref{zer:eq:fivegroup} becomes
\eqref{zer:eq:fourexception}. \Cref{zer:thm:fourpoint} keeps what is special to four nodes: the explicit
interval $I$, the density of the exceptional shapes, $\mathcal B(c)$ and the certified example.
\begin{proof}
The cancellations give $\sum_{i,j}c_ic_j(i-j)^2=0$, that is $b_1+b_2/2+b_3/3+b_4/4=0$, and
$\zerSym_c(0)=6\mathcal B(c)>0$. When $b_3=0$ the symbol is periodic in $\xi$ with period $2\pi/\log2$, which gives
the first assertion. When $b_3\ne0$ let $z=2^{i\xi}$, $w=3^{i\xi}$. The equation $\zerSym_c(\xi)=0$ determines
$w$ from $z$, and on $\abs z=1$ the condition $\abs w=1$ expands to \eqref{zer:eq:fivequad} with $X=\Re z$. This
quadratic does not vanish identically: otherwise $\abs{b_1+b_2z+b_4z^2}$ would be constant $\abs{b_3}$ on the
circle, so the polynomial would be a monomial, since the autocorrelation at the difference of its extreme
degrees is the product of the extreme coefficients; and for a monomial the linear relation above forces
the coefficient and $b_3$ to have distinct magnitudes, the associated distance $1,2$ or $4$ differing from
$3$. If $z$ were real, $w$ would be real, making $\xi\log2$ and $\xi\log3$ multiples of $\pi$, so $\xi=0$ by
irrationality of $\alpha$, which is excluded. For each interior root the two choices of $z$ give conjugate
configurations. The lifting step of the proof of \cref{zer:thm:fourpoint}, with $(\theta,\phi)$ in place of
$(u,v)$, shows that the upper one lifts to a frequency exactly when \eqref{zer:eq:fivegroup} holds, for at most
one pair $(m,n)$; conjugation accounts for the negative frequency, and distinct roots give distinct
positive frequencies.
\end{proof}

The dyadic five-point contrast above has $b_3=0$: its symbol is $(1+2^{i\xi})\zerSym_{(1,-2,1)}(\xi)$, and it misses
exactly the odd multiples of $\pi/\log2$.

\begin{theorem}[Two missing frequencies; \stU]\label{zer:thm:fivetwo}
Let $P_{s,t}(z)=(1-z)^2(1+sz+tz^2)$. In an open neighbourhood of $(s,t)=(8,14)$ every stencil has five nonzero
coefficients and $b_3\ne0$, the shapes missing exactly two positive frequencies form a countable dense subset,
the shapes missing exactly one are dense, and almost every shape misses none.
\end{theorem}
\begin{proof}
At the centre the coefficients are $(1,6,-1,-20,14)$, the vector $(b_1,b_2,b_3,b_4)$ is
$(-260,-1080,1728,896)$, and the two
roots of \eqref{zer:eq:fivequad} lie in $(-0.895,-0.894)$ and $(-0.580,-0.579)$. Near the centre both roots and
the phases $\theta_k,\phi_k$ of their upper configurations are real analytic, and the map
$(s,t)\mapsto(\phi_1-\alpha\theta_1,\phi_2-\alpha\theta_2)$ has Jacobian determinant in $(-0.008,-0.007)$ at the
centre. To verify this, write $[u,v]=\Re u\Im v-\Im u\Re v$ and put $U=i(b_2z+2b_4z^2)$, $V=ib_3w$ and
$Q=b_1'+b_2'z+b_3'w+b_4'z^2$ at each solution, a prime denoting a parameter derivative; then
$U\theta'+V\phi'=-Q$ and $\phi'-\alpha\theta'=(-[U,Q]+\alpha[Q,V])/[U,V]$. The coefficient derivatives in the order
$s,t$ are $(41,-352,405,0)$ and $(-60,16,108,64)$, and rational interval substitution gives the entries, roots
in increasing order,
\[
 \begin{pmatrix}(0.0405,0.0408)&(-0.129,-0.128)\\(-0.0161,-0.0158)&(-0.144,-0.142)\end{pmatrix},
\]
with $[U,V]\ne0$ in both rows; square roots are enclosed between consecutive multiples of $10^{-12}$ and
$1.5849<\alpha<1.5850$ by integer powers in
\begin{center}{\small\artifact{dirichlet-coupled-zeros-and-prime-variances/five_node_filters/check_five_nodes.py}}.\end{center}
The inverse
function theorem identifies a neighbourhood of the centre with an open set of phase mismatches. The group
$2\pi(\Z+\alpha\Z)$ is countable and dense, so its square cuts out a countable dense set of parameters
passing both tests, the points with exactly one coordinate in it are dense, and the union of the two
families of coordinate lines has measure zero, also after the local change of variables.
\end{proof}

So a five-point variance can converge on a spectrum with two independently adjustable invisible
oscillations: with missing frequencies $\omega_1,\omega_2$, a leading density proportional to
$1+a\cos(\omega_1\log x)+b\cos(\omega_2\log x)$, $\abs a+\abs b<1$, leaves the variance unchanged while the counting
statistic retains both terms, with factors $(1+i\omega_k)^{-1}$ in complex form. Whether the ratio of the two
frequencies has a special arithmetic nature is not known.

\paragraph{How finely must the scales be sampled?}
Call a stencil \emph{elliptic} if $\inf_{\xi\in\R}\abs{\zerSym_c(\xi)}>0$; every three-point stencil is elliptic
by the domination in \cref{zer:thm:inversemeasure}, and so is every four-point shape outside $\overline I$ in
\cref{zer:thm:fourpoint}. Write $a(T)=\pi A(T)/(T\log T)$.

\begin{theorem}[Exact mesh condition; \stU]\label{zer:thm:mesh}
Let $h_n\downarrow0$ and let $c$ be elliptic. The equivalence
\[
 U_c(h_n)/\tau(h_n)\to d\,\mathcal B(c)\iff a(T)\to d
\]
holds for every admissible positive measure $A$ and every finite $d\ge0$ if and only if
$\log(h_n/h_{n+1})\to0$. If the mesh condition fails, there is even a discrete counterexample
$A=2\sum_xm_x^2\delta_x$ with $m_x\in\set{1,2}$ whose ordinary count has the two Riemann--von Mangoldt main
terms, for which the sampled variances converge and $a(T)$ does not. The same holds under the weaker
assumption $\abs{\zerSym_c(\xi)}\ge c_0(1+\abs\xi)^{-N}$.
\end{theorem}

Put $k(s)=e^{-3s}\Phi(e^s)$ and $g_c(s)=e^sK_c(e^s)$; by \eqref{zer:eq:stencilcomp}, with $r_{ij}=u_j-u_i$,
\begin{equation}\label{zer:eq:logkernel}
 g_c(s)=2\sum_{i<j}c_ic_jr_{ij}^3\,k(s+\log r_{ij}),
\end{equation}
both kernels having exponentially decreasing envelopes at both ends. With $\check g(s)=g(-s)$ and
$\widehat f(\xi)=\int f(s)e^{-is\xi}\,\dd s$,
\begin{equation}\label{zer:eq:gammamult}
 \widehat{\check k}(\xi)=\int_0^\infty\Phi(x)x^{-4+i\xi}\,\dd x=\frac\pi{2\Gamma(4-i\xi)\cosh(\pi\xi/2)},\qquad
 \widehat{\check g_c}(\xi)=2\zerSym_c(-\xi)\widehat{\check k}(\xi):
\end{equation}
two integrations by parts reduce the first integral to $(\varsigma(\varsigma-1))^{-1}\int_0^\infty(1-\cos x)x^{1-\varsigma}\,\dd x$ at
$\varsigma=3-i\xi$, and for $0<\Re v<2$ the Laplace representation gives
$\int_0^\infty(1-\cos x)x^{-1-v}\,\dd x=\pi/(2\Gamma(1+v)\sin(\pi v/2))$; take $v=\varsigma-2$.

\begin{lemma}[Prescribing a localized response; \stU]\label{zer:lem:invert}
For every real $\phi\in C_c^\infty(\R)$ there is a real Schwartz function $f$ with $\check k*f=\phi$, and, if $c$
satisfies the polynomial lower bound, a real Schwartz $f$ with $\check g_c*f=\phi$.
\end{lemma}
\begin{proof}
Divide $\widehat\phi$ by the multiplier in \eqref{zer:eq:gammamult}. Stirling's formula gives
$\abs{\Gamma(4-i\xi)\cosh(\pi\xi/2)}\ll(1+\abs\xi)^{7/2}$ on the real axis and a polynomial bound in a narrow strip,
so Cauchy's estimate bounds every derivative polynomially; the exponential sum has bounded derivatives, so
the lower bound controls the derivatives of its reciprocal. The quotient is Schwartz, and Hermitian symmetry
makes $f$ real.
\end{proof}

\begin{lemma}[Local perturbations; \stU]\label{zer:lem:pulse}
Let $x_j\to\infty$ be separated by at least one with $x_{j+1}-x_j\to\infty$, let $f$ be real Schwartz,
$b(t)=\sum_jf(t-x_j)$, $0<q_0<1$, and $\eps>0$ so small that $q=q_0+\eps b\in(0,1)$. With
$\lambda(x)=(2\pi)^{-1}\log(x/(2\pi))$ and $\dd A(x)=2(1+q(\log x))\lambda(x)\,\dd x$ on $x\ge10$, as $t\to\infty$,
\begin{gather*}
 \frac{U_c(e^{-t})}{\tau(e^{-t})}=(1+q_0)\mathcal B(c)+\frac\eps\pi(\check g_c*b)(t)+O\bigl(t^{-1}(1+\norm b_\infty)\bigr),\\
 a(e^t)=1+q_0+\eps(\mathcal H*b)(t)+O\bigl(t^{-1}(1+\norm b_\infty)\bigr),
\end{gather*}
with $\mathcal H(s)=e^{-s}\mathbf 1_{s\ge0}$. For $f\ne0$ the second expression has at least two distinct
subsequential limits.
\end{lemma}
\begin{proof}
The separation bounds sums of translates of $(1+\abs t)^{-2}$. For the first formula substitute $x=e^{t+s}$ and
use the finite first absolute moments of $g_c$; the lower tail is exponentially small. For the second
substitute $x=e^{t-s}$ and use $\int\mathcal H=\int s\mathcal H(s)\,\dd s=1$. Convolution by $\mathcal H$, with
multiplier $(1+i\xi)^{-1}$, preserves Schwartz functions and is injective; choose $v$ with
$(\mathcal H*f)(v)\ne0$. At $x_j+v$ the other translates contribute $o(1)$, and at $(x_j+x_{j+1})/2$ the total is
$o(1)$, so the two limits differ.
\end{proof}

\begin{proof}[Proof of \cref{zer:thm:mesh}]
If the logarithmic gaps tend to zero, \cref{zer:thm:stencilbank} gives the equivalence. Otherwise there are
$\delta>0$ and centres $x_j$ of intervals of length $2\delta$ missed by all $-\log h_n$, with spacings tending
to infinity. Choose a nonzero real smooth $\phi$ supported in $(-\delta,\delta)$, apply \cref{zer:lem:invert}
with $g_c$, and sum translates; Fubini is justified by $g_c\in L^1$ and the uniform bound on the translates of
$\abs f$. Then $\check g_c*b=\sum_j\phi(\cdot-x_j)$ vanishes at every sample, and \cref{zer:lem:pulse} gives the
two behaviours. The discrete version is the construction below.
\end{proof}

\begin{theorem}[Several stencils and grids; \stU]\label{zer:thm:bankmesh}
For each of finitely many stencils $c$ let $S_c\subset[0,\infty)$ be its observed log-times $t=\log(1/h)$, and
put $\mathcal E=\bigcup_c\bigcup_{i<j,\,c_ic_j\ne0}(S_c-\log r_{ij})$. If $\mathcal E$ misses intervals of some fixed
positive length at arbitrarily large heights, there is a single positive measure, also realizable with integer
multiplicities as below, for which every stencil along its own observations has normalized variance tending to
$(1+q_0)\mathcal B(c)$, while $a(T)$ has no limit. This holds whenever each $S_c$ is locally finite with
$\limsup_T\#(S_c\cap[0,T])/T<\infty$; in particular the ratios and offsets of finitely many geometric grids may
be chosen independently.
\end{theorem}
\begin{proof}
Choose the centres and $\phi$ inside gaps of $\mathcal E$ and apply \cref{zer:lem:invert} with the common kernel
$k$. By \eqref{zer:eq:logkernel}, $\check g_c*b(t)=2\sum_{i<j}c_ic_jr_{ij}^3\sum_m\phi(t-\log r_{ij}-x_m)$, which
vanishes on $S_c$, and \cref{zer:lem:pulse} applies. Under the density hypothesis the finite union of
translates has bounded average count per unit length; if its late gaps were all shorter than every fixed
$\delta$, its lower average count would exceed $1/(2\delta)$ for every $\delta$.
\end{proof}

For integer multiplicities in both theorems, take $\eps$ smaller so that $q$ stays away from $0$ and $1$ and run
the construction of \cref{zer:prop:discrete} with $q(\log x)$ in place of $q(x)$: the ordinary count is
$\frac T{2\pi}\log\frac T{2\pi}-\frac T{2\pi}+O(1)$, the cumulative function of $2\sum_xm_x^2\delta_x$ differs from
$A(T)$ by $O(1)$, and integration by parts and the zeta denominator cost $O(h^4)$ in $U_c(h)$. These spectra
keep the two classical counting terms; realizing them as zeros of an Euler product would need further
arithmetic input. The counterexamples have $d=1+q_0>1$; for the value $d=1$, positivity of the excess
$m_\gamma(m_\gamma-1)$ makes a coarse mesh suffice, as in \cref{zer:thm:coalesceinverse}(ii).

\paragraph{A rate for the inversion.}
A rate for the variance at all small scales gives a rate for the sharp count, even for a measure with atoms.
Use the three-point stencil on the nodes $(0,a,L)$, $L=a+b$, with weights $(b,-L,a)$, so that
$\mathcal B(c)=a^2b^2L/3$ by \cref{zer:thm:inversemeasure}.

\begin{theorem}[Quantitative inversion for every three-point stencil; \stU]\label{zer:thm:rate}
Let $0<\delta\le1$ and $d\ge0$, and suppose that for all sufficiently small $h>0$
\[
 U_c(h)=d\,\mathcal B(c)h^3\log(1/h)+O\bigl(h^3(\log(1/h))^{1-\delta}\bigr).
\]
Then $\pi A(T)/(T\log T)=d+O((\log T)^{-\delta/4})$ as $T\to\infty$, the implied constant depending on the
spacings, $d$, $\delta$, the starting scale, the bound in the hypothesis and the mass on a fixed compact
interval.
\end{theorem}

\begin{proof}
Write $K=K_c$. Since $K(0)=a^2b^2L^2/4>0$, integration over $[0,\eps_0/h]$ gives $A(X)\ll(1+X)\log(2+X)$.
Compare $A$ with $\dd A_0(x)=(d/\pi)\mathbf 1_{[1,\infty)}(x)\log x\,\dd x$: from $K(x)\ll\min(1,x^{-4})$,
$\int_0^\infty K(y)\abs{\log y}\,\dd y<\infty$ and $\int_0^h\abs{\log(y/h)}\,\dd y=h$, one gets
$U_{c,A_0}(h)=d\,\mathcal B(c)h^3\log(1/h)+O(h^3)$. For $\sigma=A-A_0$ put
\[
 R(t)=\int_0^\infty\frac{\cos(tx)-1+(tx)^2/2}{x^4}\,\dd\sigma(x),\qquad r(t)=R(t)/t^3\quad(t>0),
\]
the integrand having continuous value $t^4/24$ at zero and being $O(\min(\abs{tx}^4,\abs{tx}^2))$. Splitting at
$1/t$ and integrating the tail by parts gives $\abs{r(t)}\ll1+\log(1/t)$ for small $t$, and $\abs{R(t)}\ll_{t_0}t^4$
for $t\ge t_0$. Expanding the variance cancels the constant and quadratic terms: with
$u(h)=(U_c(h)-U_{c,A_0}(h))/h^3$,
\[
 r(t)=\frac{u(t/L)}{2abL^3}+\frac{a^2}{L^2}r(at/L)+\frac{b^2}{L^2}r(bt/L),
\]
with coefficient sum $\rho=(a^2+b^2)/L^2<1$. With $p=1-\delta$ and $C_0=\log(L/\min(a,b))$, after $j$ substitutions
the forcing terms total $O(\rho^j(1+\log(1/t)+jC_0)^p)$ for $t<t_0$, and the unreplaced remainder
$O(\rho^j(1+\log(1/t)+jC_0))$ tends to zero; so $\abs{R(t)}\ll t^3(1+\log(1/t))^{1-\delta}$ for $0<t<t_0$.

Put $\dd\sigma_T(x)=\pi\,\dd\sigma(Tx)/(T\log T)$, with transform $R_T(t)=\pi T^3R(t/T)/\log T$. Then
$\abs{R_T(t)}\ll\abs t^3(\log T)^{-\delta}$ for $1\le\abs t\le t_0T$, the same times $1+\abs{\log\abs t}$ for
$\abs t<1$, and $\abs{R_T(t)}\ll\abs t^4/(T\log T)$ for $\abs t\ge t_0T$. For even smooth compactly supported real
$f$, with $\widehat g(t)=\int g(x)e^{-itx}\,\dd x$,
\[
 \int_0^\infty f\,\dd\sigma_T=\frac1{2\pi}\int_\R R_T(t)\,\widehat{x^4f}(t)\,\dd t,
\]
since $g=x^4f$ has $g(0)=g''(0)=0$, so the subtracted Taylor terms leave the inverse Fourier integral unchanged;
Fubini follows from the numerator bound and the rapid decay of $\widehat g$. Mollify the indicators of
$[-1+\eps,1-\eps]$ and $[-1-\eps,1+\eps]$ at width $\eps<\tfrac14$ to get $f_-\le\mathbf 1_{[-1,1]}\le f_+$ with
$\abs{\widehat{x^4f_\pm}(t)}\ll_N(1+\abs t)^{-1}(1+\eps\abs t)^{-N}$. The part $\abs t\le t_0T$ of the integral is
$O((\log T)^{-\delta}\eps^{-3})$, and for $\eps=(\log T)^{-\delta/4}$ and $N=8$ the tail is
$O(\eps^{-8}T^{-5}/\log T)$. The smooth tests against the rescaled $A_0$ differ from its interval count by
$O(\eps)$, positivity of $A$ sandwiches its count, and $\pi A_0(T)/(T\log T)=d+O(1/\log T)$.
\end{proof}

\begin{corollary}[A count remainder from a variance remainder; \stG, with the variance remainder as a hypothesis]\label{zer:cor:ratezeta}
For a three-point stencil, if $V_c(h)=d\,\mathcal B(c)h^3\log(1/h)+O(h^3(\log(1/h))^{1-\delta})$ with
$0<\delta\le1$, then $2\pi M_2(T)=dT\log T+O(T(\log T)^{1-\delta/4})$; a measured remainder $O(h^3)$ gives
$O(T(\log T)^{3/4})$.
\end{corollary}
\begin{proof}
Apply \cref{zer:thm:rate} to $A=2\sum_{\gamma>0}m_\gamma^2\delta_\gamma$; the change from $\gamma^4$ to
$(\tfrac14+\gamma^2)^2$ costs $O(h^4)$, as before \cref{zer:thm:coalesceinverse}.
\end{proof}

Establishing such a remainder for the actual zeta zeros is a further arithmetic problem, and the exponent
$\delta/4$, which comes from smoothing an interval of relative width $\eps$, is not claimed optimal; for Fourier
Tauberian estimates for monotone counting functions see \cite{Safarov2001Fourier}.

\paragraph{Residue classes.}
Fix $q$, let $\chi$ run over the $n_q$ characters modulo $q$, with the multiplicities $m_\chi(\gamma)$ of the
primitive inducers at distinct signed ordinates and $M(\gamma)=\sum_\chi m_\chi(\gamma)$, and let $Z_a$ be as in
\eqref{zer:eq:va}. By \cref{zer:lem:ramp}, $F_\chi=\sum_a\chi(a)Z_a$ is a constant plus
$\sum_{\gamma\ne0}m_\chi(\gamma)(\tfrac14+\gamma^2)^{-1}e^{i\gamma t}$ plus a decaying term, and $D_h$ kills the constant.
With long-time inner products,
\[
 W_{\chi\psi}(h)=\langle D_hF_\chi,D_hF_\psi\rangle=\sum_{\gamma\ne0}\frac{\abs{P(h\gamma)}^2m_\chi(\gamma)m_\psi(\gamma)}{(1/4+\gamma^2)^2},
 \qquad C_{ab}(h)=\langle D_hZ_a,D_hZ_b\rangle ,
\]
and $C=HWH^*$ with $H_{a\chi}=\overline{\chi(a)}/n_q$, $HH^*=I/n_q$.

\begin{theorem}[Pairwise inversion and the identity class; \stG]\label{zer:thm:identityclass}
\begin{enumerate}[label=\textup{(\roman*)}]
\item For characters $\chi,\psi$ modulo $q$ and finite $d\ge0$, $W_{\chi\psi}(h)/\tau(h)\to B_{a,b}d$ if and only
if
\[
 \frac\pi{T\log T}\sum_{0<\abs\gamma\le T}m_\chi(\gamma)m_\psi(\gamma)\to d ;
\]
$h=1/k$ suffices. Hence $C/\tau$
converges entrywise exactly when all these folded pair counts converge, with limit $B_{a,b}HDH^*$ for
the matrix $D$ of their limits.
\item $\liminf C_{11}(h)/\tau(h)\ge B_{a,b}/n_q$, and the following are equivalent:
$C_{11}(h)/\tau(h)\to B_{a,b}/n_q$ for one pair $a,b>0$; $\sum_{0<\abs\gamma\le T}M(\gamma)(M(\gamma)-1)=o(T\log T)$;
$C(h)/\tau(h)\to(B_{a,b}/n_q)I$ for every $a,b>0$; simple zeros have density one among the zeros of the
Dedekind zeta function of $\Q(\zeta_q)$, counted with multiplicity. In the first condition it suffices
to take $h=2^{-k}$.
\end{enumerate}
\end{theorem}
\begin{proof}
(i) is \cref{zer:thm:inversemeasure} with mass $m_\chi(\gamma)m_\psi(\gamma)$ at $\abs\gamma$; the unit-bin mass is
$O_q(\log^2(2+\abs y))$. For (ii) let $J_\chi$ be $W_{\chi\chi}$ with $m_\chi^2$ replaced by $m_\chi$;
Riemann--von Mangoldt gives $J_\chi/\tau\to B_{a,b}$, and $C^0=H\operatorname{diag}(J_\chi)H^*$ has
$C^0/\tau\to B_{a,b}I/n_q$. All entries of $W-\operatorname{diag}(J_\chi)$ are nonnegative, every entry of $H$
has modulus $1/n_q$ and its identity row is constant, so
\[
 \abs{C_{uv}-C^0_{uv}}\le C_{11}-C^0_{11}=\frac1{n_q^2}\sum_{\gamma\ne0}\frac{\abs{P(h\gamma)}^2M(\gamma)(M(\gamma)-1)}{(1/4+\gamma^2)^2}.
\]
This gives the lower bound, and \cref{zer:thm:inversemeasure} with $d=0$ applied to the excess gives the
equivalence of the first two conditions and then the third. The dyadic sufficiency follows from the
shell argument of \cref{zer:thm:coalesceinverse}(ii) with $\abs P$ bounded below on
$[2^{-s-1},2^{-s}]$. By the cyclotomic factorization $M(\gamma)$ is the multiplicity of the Dedekind zeta
zero, whose count is $\sim n_qT\log T/\pi$. The bound $\sum_{0<\abs\gamma\le T}M(\gamma)^4\ll_qT\log T$ (from the
fourth moment of the positive test counts $Y_j$ in the proof of \cref{zer:thm:familyrate} below, and
$(\sum_jm_j)^4\le n_q^3\sum_jm_j^4$) and Cauchy--Schwarz show that density-one simplicity makes the
total squared multiplicity of the multiple zeros $o(T\log T)$; conversely $M\le M(M-1)$ when $M\ge2$.
\end{proof}

For $q=4$ the inducers are $\zeta$ and the Dirichlet beta function; with multiplicities $m,r$ at an
ordinate, $(m+r)(m+r-1)=m(m-1)+r(r-1)+2mr$, so the coefficient $\tfrac13$ for the second difference in
the class $1\bmod4$ forces all three weighted excesses to be $o(T\log T)$, the same coefficient in the
class $3\bmod4$ and zero limiting cross coefficient.

\begin{theorem}[Local bounds for the multiplicity matrix; \stG]\label{zer:thm:familyrate}
Let $\chi_1,\dots,\chi_d$ be distinct primitive characters, the conductor-one character allowed, and
$m(\gamma)=(m_j(\gamma))_j$. There are $0<c<C$ depending only on the family such that for fixed
$0<A<B$ and all large $T$,
\[
 c(B-A)T\log T\,I_d\preceq\sum_{AT<\gamma\le BT}m(\gamma)m(\gamma)^{\mathsf T}\preceq C(B-A)T\log T\,I_d ,
\]
as Hermitian forms, and likewise on negative heights. The same holds on $[T,T+H]$ with
$T^\delta\le H\le T$, with bounds multiples of $H\log T$ and constants depending also on $\delta$.
\end{theorem}
\begin{proof}
Let $L=\log T$, let $w\ne0$ be smooth with compact support in $(\alpha,\beta)$, $0<\alpha<\beta<\tfrac18$, and
$f(x)=\int w(v)e^{ivx}\dd v$. For a unit $c\in\C^d$ put $z(\gamma)=\sum_jc_jm_j(\gamma)$,
$U(n)=\sum_jc_j\chi_j(n)$ and $X(t)=\sum_\gamma z(\gamma)f(L(\gamma-t))$. The explicit formula
\cite[Lemma~2.1]{LiuYe2005Distribution} for each primitive inducer gives, uniformly for $t/T$ in compact
positive ranges,
$X(t)=-L^{-1}\sum_{n\le T^\beta}\Lambda(n)U(n)n^{-1/2}w(\log n/L)n^{-it}+o(1)$: only positive frequencies
occur, the archimedean constant vanishes since $\int f=2\pi w(0)=0$, and the pole terms of the
conductor-one character are $o(1)$ by integration by parts. On an interval $J$ of length $H\asymp T$,
$H^{-1}\int_J\abs{\sum_{n\le N}d_nn^{-it}}^2\dd t=(1+O(N\log(2N)/H))\sum\abs{d_n}^2$. For the prime part $P$,
the prime number theorem in reduced classes and orthogonality of the distinct induced characters give
$\sum_p\abs{d_p}^2\to\int v\abs{w(v)}^2\dd v>0$ uniformly in $c$, while higher prime powers contribute
$O(L^{-2})$; so $\norm{X+P}_{L^2(J,\dd t/H)}=o(1)$. The same estimate for $P^2$ and unique factorization
give $\norm P_4=O(1)$, and $\norm P_2\le\norm P_1^{1/3}\norm P_4^{2/3}$ gives $H^{-1}\int_J\abs X\gg1$. With $J$
the middle half of $[AT,BT]$, the ordinates outside contribute $o(1)$ by rapid decay, and each ordinate
inside at most $\abs{z(\gamma)}\norm f_1/L$, so $\sum_{AT<\gamma\le BT}\abs{z(\gamma)}\gg(B-A)TL$; the
Riemann--von Mangoldt count and Cauchy--Schwarz give the lower bound. For the upper bound take a
nonnegative Schwartz $g$ with $g(0)>0$ and smooth Fourier transform supported in $(-\beta,\beta)$; for
$Y_j(t)=\sum_\gamma m_j(\gamma)g(L(\gamma-t))$ the same formula gives two short polynomials with bounded
coefficient square sums and a bounded archimedean term, so $\int_{AT-1}^{BT+1}\abs{Y_j}^2\ll(B-A)T$, and the
diagonal terms bound this below by $L^{-1}\int_{-1}^1g^2\cdot\sum_{AT<\gamma\le BT}m_j(\gamma)^2$; sum over $j$ and use
$\abs{z(\gamma)}^2\le\sum_jm_j(\gamma)^2$. The same positivity applied to $Y_j^4$, whose prime polynomials have
bounded fourth moments, gives $\sum_{T<\gamma\le2T}m_j(\gamma)^4\ll T\log T$. For short intervals choose
$2\beta<\delta$.
\end{proof}

For two functions and the difference of their multiplicities, lower bounds of order $T\log T$ are due
to Fujii and to Bombieri and Perelli \cite{Fujii1976Zeros,BombieriPerelli1998Distinct}; the matrix form
is uniform in the direction.

\begin{corollary}[Orders of the residue-class covariance; \stG]\label{zer:cor:classorders}
Let $d=n_q$, fix distinct offsets $u_1,\dots,u_r$, and let $b_j\in\R^d$ satisfy $\sum_jb_j=\sum_ju_jb_j=0$,
with $H_b(v)=\sum_jb_j(v-u_j)_+$. Uniformly in such $b$, as $h\downarrow0$,
\[
 \mathcal M\Bigl\lvert\sum_jb_j^{\mathsf T}Z(t+hu_j)\Bigr\rvert^2\asymp_q h^3\log(1/h)\int_\R\norm{H_b(v)}^2\dd v ,
\]
$\mathcal M$ the long-time mean. For $r\ge3$ the $rd$-dimensional covariance of
$(Z(t+hu_1),\dots,Z(t+hu_r))$ has $d$ eigenvalues of order $1$, $d$ of order $h^2$ and $(r-2)d$ of order
$\tau(h)$, even with coincidences among ordinates; its determinant has order
$h^{(3r-4)d}\log(1/h)^{(r-2)d}$.
\end{corollary}
\begin{proof}
Place $m(\gamma)m(\gamma)^{\mathsf T}/(T\log T)$ at $\gamma/T$, $T=1/h$. By \cref{zer:thm:familyrate} every vague
limit lies between $cI_d\,\dd x$ and $CI_d\,\dd x$, and dyadic summation gives
$\sum_{0<\abs\gamma\le X}\norm{m(\gamma)}^2\ll X\log(2+X)$, which controls the ends since
$\norm{\sum_jb_je^{iu_jx}}\ll\norm b\min(x^2,1)$. The character-to-class matrix and Parseval bound the
normalized variance by constant multiples of $d^{-1}\int\norm{\sum_jb_je^{iu_jx}}^2x^{-4}\dd x
=2\pi d^{-1}\int\norm{H_b}^2$, a positive definite form on the contrast space. For the eigenvalues use an
orthonormal offset basis beginning with the constant and the centred vectors. The first observation
vector converges in $\ell^2$ of the ordinates to a multiple of the ramp vector, the second divided by
$h$ to a multiple of its derivative, and their joint limiting Gram is positive definite: a relation
would give $m(\gamma)^{\mathsf T}(v+i\gamma w)=0$ for all ordinates, and on $[T,2T]$
\cref{zer:thm:familyrate} forces $w=0$, then $v=0$. The remaining vectors divided by $\sqrt{\tau}$ have
Gram matrices bounded above and below by the first part and converge weakly to zero, so the mixed
entries vanish; min--max gives the orders.
\end{proof}

\begin{corollary}[A criterion including the critical line; \stU]\label{zer:cor:critline}
For $a=b=1$ the following are equivalent: the mean squares $C_{11}(2^{-k})$ exist for all large $k$ and
$C_{11}(2^{-k})/\tau(2^{-k})\to2/(3n_q)$; every nontrivial zero of $\zeta_{\Q(\zeta_q)}$ lies on the
critical line and its simple zeros have density one.
\end{corollary}
\begin{proof}
Use the unconditional centering constant given by the constant Laurent coefficients of $-L'/L$ at
$s=1$. With $\lambda=s-\tfrac12$, the Laplace transform of the identity-class ramp is
$\sum_\chi(-L'/L)(1+\lambda,\chi)/(n_q\lambda(1+\lambda))-1/(n_q\lambda^2)-c/\lambda$, and that of its second
difference is $(e^{hs}-1)^2$ times this plus an entire function. A finite long-time mean square at one
$h$ makes the latter holomorphic on $\Re s>0$ by Cauchy--Schwarz, whereas a zero with real part above
$\tfrac12$ would give a pole, since multiplicities add. The functional equation and the cyclotomic
factorization give the critical line, and \cref{zer:thm:identityclass}(ii), whose proof uses the Riemann
hypothesis only for the primitive inducers modulo $q$, gives the rest.
\end{proof}

The infinite-time average removes all unequal frequencies before the lag tends to zero, so these
statements are insensitive to any finite set of zeros and differ from the short-interval variance
problems of full pair correlation. Multiplicative cancellation, Mellin nonvanishing and logarithmically
periodic counterexamples occur in inverse problems for regular variation
\cite{JacobsenMikoschRosinskiSamorodnitsky2009Inverse,DamekMikoschRosinskiSamorodnitsky2014General}, and
Greenhall classified the spectral ambiguity of the Allan variance, another scale-dependent filtered
variance \cite{Greenhall1998Spectral}. Exact checks of
the spline norms $\mathcal B_2,\dots,\mathcal B_5$ and of the minimal covariance $\tfrac16I_4$ at modulus
$8$ are in
\begin{center}\artifact{dirichlet-zero-multiplicity-and-reconstruction/checks/contrast-results.txt}.\end{center}

\subsection{Open questions}\label{flu:sec:open}
{\raggedright\sources{grh_cm_four_terms_and_prime_bias; grh_multiscale_covariance; multiscale_prime_fluctuations; grh_spectral_rigidity; grh_filtered_prime_races; grh_prime_error_balancing; prime_peak_rigidity; sparse_grh_balanced_windows; sharp_multiscale_prime_windows; prime_scale_coalescence_2026-11-26; prime_endpoints_and_symplectic_compressions; grh_extensions_2026-11-27; dirichlet-zero-multiplicity-and-reconstruction; grh-two-open-questions; dirichlet-coupled-zeros-and-prime-variances; grh_method_transfers_sharp_torsion}}

\begin{enumerate}
\item \emph{Effective constants.} Can the radius $r_q$ and ratio $C_q$ of
\cref{zer:thm:balancing}, the window $H$ and radius of \cref{zer:cor:smoothbalance}, the
constants $c,C,K$ of \cref{zer:thm:coercive}, the least eigenvalues in
\cref{zer:thm:jets,zer:thm:filtered} and a certified window length $L_0$, and the return lengths
near the peak in \cref{zer:thm:peak}, be made explicit, ideally uniformly in $q$, in the number
of scales, or in the cut $H$? Near-coincident ordinates must then be separated rigorously, and
efficient probe primes chosen.
\item \emph{The eigenvalue gap.} The two bounds of \eqref{zer:eq:coercive} differ by three powers
of $\log T$. Can information on pair correlation and multiplicities close the gap, and what are
the sharp asymptotics for arbitrary shifts under stronger hypotheses?
\item \emph{Merging scales.} \Cref{zer:thm:coalesce} settles the rate for $\zeta$, and
\cref{zer:cor:classorders} the order of every eigenvalue of the residue-class matrix. What is the best
upper factor in \cref{zer:lem:localM2}, now at most $1.3208$, over all arithmetic methods, and does
$M_2(T)\sim\kappa T\log T/(2\pi)$ hold, as \cref{zer:prop:kappa} requires (equivalently, by
\cref{zer:thm:coalesceinverse}, does one three-point variance have a normalized limit)? What are the
exact leading constants of the residue-class matrix, which by \cref{zer:thm:identityclass} are the
limits of the folded pair counts? \Cref{zer:thm:rate} turns a variance remainder with relative saving
$(\log(1/h))^{-\delta}$ into a count remainder with saving $(\log T)^{-\delta/4}$: is $\delta/4$ optimal, and
can such a remainder be proved for $\zeta$? Families of stencils characterize all finite limits exactly
when their symbols have no common real zero (\cref{zer:thm:stencilbank}), on meshes whose logarithmic gaps
tend to zero, which is necessary for elliptic stencils (\cref{zer:thm:mesh}); four and five equally spaced
observations are classified by \cref{zer:thm:fourpoint,zer:thm:fivepoint}. Which small-divisor conditions on
symbols with arbitrarily small values suffice for the exact mesh characterization, and what condition is
necessary and sufficient for asynchronously sampled families whose individual meshes have gaps while the
set $\mathcal E$ of \cref{zer:thm:bankmesh} has none? For $\zeta$ itself, can arithmetic constraints on
multiple zeros beyond \cref{zer:lem:localM2} exclude a nonconstant tangent density in
\cref{zer:thm:missingfreq}, or the pulses of \cref{zer:thm:mesh}? That would remove the exceptional shapes
from \cref{zer:thm:coalesceinverse}(i).
\item \emph{Beyond covariance.} Can the convex-hull statements be upgraded to the interior of the
actual limiting support, to every ordering of the unsmoothed counts, or to inclusivity of
ordinary prime races, without a hypothesis on the ordinates? Which nonlinear constraints hold on
the cluster sets? Positive covariance excludes affine hyperplanes only.
\item \emph{Less smoothing and other families.} \Cref{zer:thm:unsmoothed} gives strict positivity
of the multiscale covariance for the unsmoothed counts. Does the jet covariance of
\cref{zer:thm:jets} survive for unsmoothed counts after the prime-square biases are removed? Can the constants of \cref{zer:thm:unsmoothed} be made uniform as shifts coalesce or as the
number of scales grows? Does \cref{zer:thm:polyrigid} extend to families with
other Gamma factors and distinct Euler coefficients, and do the probe arguments extend to
Frobenius-class functions of Artin $L$-functions under their own Riemann hypothesis (which
\cref{src:thm:dgrh} does not provide)?
\item \emph{Kernels and filters.} Which compact filters retain every noncentral frequency with the
shortest memory or the smallest error, and which other Mellin kernels have nonnegative spectral
weights on the critical line, with sharper laws or shorter recurrence? Can the compact window
kernel of \cref{zer:thm:sharpwin} be optimized for given contrasts, or constrained so that the pole
terms cancel, and do its finite Gram conditions on separated windows have a converse forcing GRH
without the full positivity criterion of the explicit formula?
\item \emph{Endpoint visits.} Is the constant $n_qC^2$ of \cref{zer:thm:visits}(i), twice the sharp
leading constant under \eqref{zer:eq:rankloss}, improvable under \cref{src:thm:dgrh} alone? Does a
substantially weaker bound on the rational rank defect preserve the conductor term in
\eqref{zer:eq:visitlaw}? Are there useful finite-height rates for the actual densities without any
assumption on small integer combinations of ordinates?
\end{enumerate}

\begin{question}[The lower limit of the source remainder]\label{ana:q:parity}
Does the parity condition of \cref{ana:prop:parity} hold for the ordinates of $\zeta$ and of
$L(\cdot,\chi_d)$? If it does, $\liminf_{x\to\infty}\sqrt x\,R_d(x)=-B_d-B_0$ exactly. In
general, where in the interval of \cref{ana:thm:excursion} does the lower limit lie?
\end{question}
